\documentclass[11pt]{article}
\usepackage[T1]{fontenc}
\usepackage{amsmath,amssymb,amsthm}
\usepackage[margin=1in]{geometry}
\usepackage{hyperref}
\newtheorem{theorem}{Theorem}
\newtheorem{lemma}{Lemma}
\newtheorem{corollary}{Corollary}
\newcommand{\RV}{\mathrm{RV}}
\newcommand{\RK}{\mathrm{RK}}
\title{Equal-capacity greedy matching:\\a proof via conditional association}
\author{MathIsEvenEasier\\\small Prepared with OpenAI Codex (GPT-6 Astra)}
\date{5 October 2026}
\begin{document}
\maketitle

\noindent\textbf{Status.} The equal-capacity comparison is proved for
all finite numbers of bins and jobs, every common positive integer capacity,
and every fixed degree sequence. The complete source-model statement,
including independent uniform neighborhoods and uniform random arrival
and priority orders, is now formalized in Lean as
\texttt{source\_matching\_tail\_comparison}
(Theorem~\ref{thm:leansource}). The written proof uses the published
Cohen--Sackrowitz association lemma; the Lean proof establishes the needed
finite rational association theorem directly by induction on the common
cap, without assuming the literature result as an axiom. Verification and
axiom auditing run in bounded Azure jobs. No independent human review or
novelty claim is made. Earlier partial regimes and experiments remain below
as a research record, not as premises of the final proof.

\section{The model and the target}
There are $m\geq 1$ days, each of integer capacity $C\geq 1$.
Job $j$ has degree $d_j\in\{0,\ldots,m\}$ and independently chooses a
uniform $d_j$-subset of the days as its feasible set. Jobs are processed in
a uniformly random order, independent of these sets.
RANDOM-VERTEX chooses uniformly among the job's feasible days with space.
RANKING draws a single uniform priority order of days and always chooses
the highest-priority feasible day with space. A job is rejected if no
such day exists. Write $M_\RV$ and $M_\RK$ for the numbers accepted.

Arnosti proves equality of the output distributions for $C=1$ and
stochastic dominance of RANDOM-VERTEX for $C=2$. The conjecture following
Theorem 2 in his paper asks for the same dominance for every $C>2$:
\[
 \Pr(M_\RV\geq t)\geq\Pr(M_\RK\geq t)\qquad\text{for every integer }t.
\]
All expectations and probabilities here average over the random graph
and the algorithms' random choices. An adverse fixed graph is not a
counterexample to this assertion. See
\href{https://nickarnosti.com/ResearchPapers/Arnosti_GreedyMatching.pdf}{\emph{Greedy Matching in Bipartite Random Graphs}},
Stochastic Systems 12(2), 133--150 (2022),
\href{https://doi.org/10.1287/stsy.2021.0082}{doi:10.1287/stsy.2021.0082}.

\section{The all-capacity closure from a published association theorem}
\label{sec:literatureclosure}
This section closes the former unbounded obligation. The proof uses a
stationary correlation theorem, so it does not require monotonicity of
the auxiliary redistribution process or its Poisson potential.
For vectors of equal total, $x$ majorizes $y$ if
$\sum_{i=1}^k x_{(i)}\geq\sum_{i=1}^k y_{(i)}$ for every $k$,
where the coordinates in parentheses are sorted in descending order.
Thus a majorization-increasing function increases when the load profile
becomes more concentrated.

\begin{lemma}[Cohen--Sackrowitz association, specialized]\label{lem:csassociation}
Fix integers $b\geq1$, $q\geq0$, and $0\leq r\leq bq$. Give
\[
 \Omega_{b,r,q}=\{x\in\{0,\ldots,q\}^b:\textstyle\sum_i x_i=r\}
\]
the probability law $\mu_{b,r,q}(x)$ proportional to
$\prod_i1/x_i!$. If $F,G$ are symmetric and increasing under
majorization on this space, then
\[
 \operatorname{Cov}_{\mu_{b,r,q}}(F,G)\geq0.
\]
\end{lemma}
\begin{proof}
The source is Arthur Cohen and Harold B. Sackrowitz,
\href{https://doi.org/10.1214/aos/1176350376}
{\emph{Unbiasedness of Tests for Homogeneity}},
\emph{The Annals of Statistics} 15(2), 805--816 (1987).
Lemma 4.1, equation (4.1), on pages 811--812 establishes conditional
association of the upper sums of the order statistics of iid variables
with a common $PF_2$ carrier, given their total. Definition 2.2,
the definition of $\mathcal H_k$ on page 807, and Remark 2.5 on page 808
identify symmetric Schur-convex functions with the required increasing
functions of these upper sums. Page 809 explicitly includes the
integer-valued case in the statements of Sections 3 and 4; Section 5,
pages 812--815, supplies its tie-aware extension. Equation (1.1) on
page 805 permits counting measure, so the application here is discrete.

To check the distributional hypotheses, use the common carrier
\[
 h_q(k)=\begin{cases}1/k!,&0\leq k\leq q,\\0,&\text{otherwise},\end{cases}
 \qquad k\in\mathbb Z.
\]
Its support is an interval of integers and its successive ratios
$h_q(k+1)/h_q(k)=1/(k+1)$ decrease. Thus it is log-concave, equivalently
$PF_2$ in the source's terminology. Let $X_1,\ldots,X_b$ be independent
with mass $h_q(k)/\sum_{j=0}^q h_q(j)$. Their conditional law given
$\sum_iX_i=r$ is exactly $\mu_{b,r,q}$, and this conditioning event
has positive probability. The source lemma applies to $F,G$ written as
functions of the upper sums of the sorted $X_i$. If a function is
required outside the support, extend it at an upper-sum vector $t$
by the maximum of its values at supported vectors coordinatewise
below $t$, using its minimum supported value when that set is empty.
This is an increasing extension and agrees with the original function
on the support. All functions are bounded, so integrability is automatic.
The cases $b=1$ or $q=0$ are deterministic.

This checks the cap by including it in the carrier \emph{before}
applying the theorem; it does not assume that arbitrary conditioning
preserves association. Sorting automatically includes the multinomial
orbit multiplicities. No uniform measure on integer partitions is used.
\end{proof}

\begin{lemma}[A shared two-bin calculation]\label{lem:pairmeans}
For $u_j\geq v_j\geq0$, $1\leq j\leq t$, put
\[
 Q(X,Y)=\prod_{j=1}^t(u_jX+v_jY),\qquad
 g_k=\binom tk^{-1}[X^kY^{t-k}]Q\quad(0\leq k\leq t).
\]
The sequence $g_k$ is nondecreasing and discretely convex:
$\Delta g_k\geq0$ and $\Delta^2g_k\geq0$ wherever defined.
\end{lemma}
\begin{proof}
For a uniform permutation $\pi$ of $1,\ldots,t$,
$g_k=\mathbb E[\prod_{\ell\leq k}u_{\pi(\ell)}
\prod_{\ell>k}v_{\pi(\ell)}]$. Direct expansion gives, for $h=1,2$,
\[
 \Delta^h g_k=\mathbb E\left[
 \prod_{\ell\leq k}u_{\pi(\ell)}
 \prod_{k<\ell\leq k+h}(u_{\pi(\ell)}-v_{\pi(\ell)})
 \prod_{\ell>k+h}v_{\pi(\ell)}\right]\geq0.
\]
For $h=2$ the two middle terms have the same expectation by exchanging
positions $k+1,k+2$. This proves both assertions, including zero factors.
\end{proof}

\begin{theorem}[Uniform conditional concentration]\label{thm:allconcentration}
Let $r$ independent categorical choices on $b$ bins have probability
rows nonincreasing in the same bin order. Write $N$ for their load vector,
and suppose $E=\{\max_iN_i\leq q\}$ has positive probability.
For every symmetric majorization-increasing function $F$ on
$\Omega_{b,r,q}$,
\[
 \mathbb E(F(N)\mid E)\geq\mathbb E_{\mu_{b,r,q}}F.
\]
\end{theorem}
\begin{proof}
Let $P(z)=\prod_{j=1}^r\sum_{i=1}^b p_{j,i}z_i$, and let
$\overline P$ be its average over all permutations of the variables.
Symmetrizing the bin labels changes neither $F$ nor the cap event.
Put
\[
 L(x)=\left(\prod_i x_i!\right)[z^x]\overline P(z).
\]
The capped symmetrized load law has density $L/\mathbb E_\mu L$
relative to $\mu=\mu_{b,r,q}$: its mass is $[z^x]\overline P$,
whereas $\mu(x)$ is proportional to $1/\prod_i x_i!$.

We show that $L$ increases in majorization. Consider balancing two
loads $a\geq c+2$ to $a-1,c+1$. In the permutation average, pair
terms that exchange the corresponding variables $X,Y$, and fix the
choices assigned to the other variables. The remaining $t=a+c$
choices contribute $Q(X,Y)+Q(Y,X)$, where
\[
 Q(X,Y)=\prod_{j=1}^t(u_jX+v_jY),\qquad u_j\geq v_j\geq0.
\]
By the common row ordering and Lemma~\ref{lem:pairmeans}, the
normalized coefficients $g_k$ of $Q$ are convex. The contribution
of this pair to $L$ is a nonnegative constant times
\[
 a!c!\,[X^aY^c](Q(X,Y)+Q(Y,X))=t!\,(g_a+g_c).
\]
Convexity gives $g_a+g_c\geq g_{a-1}+g_{c+1}$. The other
coordinate factors are unchanged. Thus balancing decreases $L$;
such transfers generate integer majorization.

Lemma~\ref{lem:csassociation} now yields
\[
 \mathbb E(F(N)\mid E)
   =\frac{\mathbb E_\mu(FL)}{\mathbb E_\mu L}
   \geq\mathbb E_\mu F.
\]
The denominator is positive because $\Pr(E)>0$.
\end{proof}

\begin{lemma}[The uniform next choice minimizes the hazard]\label{lem:nextuniform}
Let independent old choices on $b$ bins have probability rows
nonincreasing in the same order. Condition on every load being at
most $q$, where the conditioning event has positive probability.
Then $p_i:=\Pr(N_i=q\mid\max_jN_j\leq q)$ is nonincreasing in $i$.
Consequently any independent next choice with a nonincreasing row
has saturation probability at least $b^{-1}\sum_i p_i$.
\end{lemma}
\begin{proof}
Take bins $i<j$. Condition on which old choices landed in the pair
and on all other destinations. If there are $t$ choices in the pair,
their probabilities of choosing $i,j$ satisfy $u_\ell\geq v_\ell$.
Thus the count $X$ of choices of $i$ has
$\Pr(X=k)=\binom tk g_k$, with $g$ from Lemma~\ref{lem:pairmeans}.
For $t<q$ neither bin reaches $q$; positive cap probability rules out
$t>2q$. Otherwise $t-q\leq q\leq t$. Lemma~\ref{lem:pairmeans}
and $\binom tq=\binom t{t-q}$ give
\[
 \Pr(X=q)=\binom tq g_q\geq\binom t{t-q}g_{t-q}=\Pr(X=t-q).
\]
Conditioning on the common interval $t-q\leq X\leq q$ preserves
this comparison, as does averaging over the finer information.
Thus $p_i\geq p_j$. For the next row $v$, the two similarly ordered
sequences satisfy
\[
 \sum_i v(i)p_i\geq\frac1b\left(\sum_i v(i)\right)
                                  \left(\sum_i p_i\right)
 =\frac1b\sum_i p_i,
\]
because the difference between the two sides is
$\frac1b\sum_{i<j}(v(i)-v(j))(p_i-p_j)\geq0$.
\end{proof}

For use in the matching process, define the uniform saturation hazard by
\[
 h_{b,C}(r)=\frac1b\,
 \mathbb E_{\mu_{b,r,C-1}}\#\{i:x_i=C-1\}.
\]
It is the expected fraction of bins one choice away from saturation.
The coupling needs only that this number depends on $b,r,C$;
no explicit coefficient formula is needed.

\begin{corollary}[The conditional hazard for every capacity]\label{cor:allhazard}
For $b\geq1$, $C\geq1$ and $0\leq r\leq b(C-1)$, take independent
old choices $Z_1,\ldots,Z_r$ and a next choice $Y$ on $b$ bins,
all with probability rows nonincreasing in the same order. If
$N_i=\#\{j:Z_j=i\}$ and $E=\{\max_iN_i<C\}$ has positive probability,
then
\[
 \Pr(N_Y=C-1\mid E)\geq h_{b,C}(r).
\]
Equality holds when every row is uniform. For $C=1$ necessarily
$r=0$ and the hazard is one. For $C\geq2$ this is the assertion
$\mathcal H(b,C)$ used in the historical reductions below.
\end{corollary}
\begin{proof}
Put $q=C-1$ and $B(x)=\#\{i:x_i=q\}$. A balancing transfer in
the capped simplex cannot increase $B$, so $B$ is
majorization-increasing. Theorem~\ref{thm:allconcentration} gives
$\mathbb E(B(N)\mid E)\geq\mathbb E_\mu B$.
Lemma~\ref{lem:nextuniform} shows that replacing the independent ordered
next choice by a uniform choice cannot increase its saturation hazard.
That uniform-choice hazard is $\mathbb E(B(N)\mid E)/b$.
By definition $\mathbb E_\mu B/b=h_{b,C}(r)$, proving the claim,
including $r=0$ and $r=bq$.
\end{proof}

\subsection{The matching bridge, without a first-saturation coupling}
The universal hazard bound also covers the case of \emph{zero} full
days. This lets us start directly at zero accepted jobs, without
conditioning on a future first-saturation time. The following argument
is the complete matching bridge used in the main theorem; it does not
depend on the intermediate matching results in the research record.

Delete degree-zero arrivals and number the remaining arrivals
$1,\ldots,n$. For RANKING fix the priority order $1,\ldots,m$;
relabeling invariance of the neighborhood model makes every fixed
order have the same output law. Let $\sigma_K$ be the arrival accepting
the $K$th job, with $\sigma_K=\infty$ if there is no such arrival.
When $\sigma_K<\infty$, let $A_K$ count the full days just after it.
Put $(\sigma_0,A_0)=(0,0)$. We compare the distribution of this pair
at each $K$; a level never reached is an absorbing cemetery state.
Joint distributions of whole trajectories are not needed.

\begin{lemma}[The conditional transition at every accepted count]
\label{lem:allhistory}
Fix $K<mC$ and a positive-probability state
$(\sigma_K,A_K)=(s_0,a)$ with $s_0<\infty$.
Put $b=m-a$ and $r=K-aC$.
For either algorithm the next-success index $S=\sigma_{K+1}$ has law
\begin{equation}\label{eq:waitingkernel}
 \Pr(S=s\mid\sigma_K=s_0,A_K=a)=
 \left(\prod_{j=s_0+1}^{s-1} f_j(a)\right)(1-f_s(a)),
 \qquad f_j(a)=\frac{\binom a{d_j}}{\binom m{d_j}},
\end{equation}
for $s_0<s\leq n$, and mass
$\prod_{j=s_0+1}^n f_j(a)$ at infinity.
Conditional also on $S=s<\infty$, the mark increment
$A_{K+1}-a$ is Bernoulli with probability exactly $h_{b,C}(r)$
for RANDOM-VERTEX and at least $h_{b,C}(r)$ for RANKING.
This includes $a=0$ and $b=1$.
\end{lemma}
\begin{proof}
Refine the current state to data $\rho$ specifying the set $D$ of
days full at $\sigma_K$, the
destination of every job assigned to $D$, and the acceptance or
rejection of every arrival through $\sigma_K$. Thus $\rho$ also
determines each full day's filling time and the full-day set
$D_{j-1}$ before every earlier arrival $j$. Let $I$ be the final
$b$ nonfull days and $J$ the accepted jobs whose destinations in
$I$ are not yet specified. There are $|J|=K-aC=r$ such jobs.

For a fixed availability set, let $u_j(i)$ be the unconditional
probability that arrival $j$ chooses available day $i$. If $v$ days
are full, these probabilities are
\[
 u_j(i)=\frac{1-f_j(v)}{m-v}\quad(\RV),\qquad
 u_j(i)=\frac{\binom{m-s_i}{d_j-1}}{\binom m{d_j}}
       \quad(\RK),
\]
where $s_i$ is the position of $i$ among the available days in
priority order. The second formula counts neighborhoods containing
$i$ and avoiding the $s_i-1$ earlier available days; full days
are allowed among the other neighbors. Thus the RV row is uniform
and the RK row is nonincreasing. These formulas depend on which
days are full, not on the partial loads of available days.

For proposed destinations $z=(z_j:j\in J)\in I^J$, write
$N_i(z)=\#\{j\in J:z_j=i\}$. The joint path probability factors as
\begin{equation}\label{eq:refinedfactor}
 \Pr(\rho,z)=\kappa_\rho
      \prod_{j\in J}u_j^{D_{j-1}}(z_j)
      \mathbf1_{\{\max_{i\in I}N_i(z)<C\}},
\end{equation}
where $\kappa_\rho$ contains the probabilities of prescribed
destinations in $D$ and prescribed rejections. Independence of fresh
neighborhoods and tie-breaking gives this product along the prescribed
availability path. Conversely, every $z$ satisfying the cap realizes
that path: loads only increase, so no member of $I$ can have filled
earlier; the assigned job sets in $D$ already fix all other filling
times. The success indicators also fix $\sigma_K$.
There is no additional restriction on $z$.

Normalize each restricted row by
\[
 p_{j,i}=\frac{u_j^{D_{j-1}}(i)}
                  {\sum_{\ell\in I}u_j^{D_{j-1}}(\ell)},
 \qquad i\in I.
\]
Every denominator is positive on a positive-probability refinement.
Equation~\eqref{eq:refinedfactor} proves that the residual loads
are independent choices with these rows, conditioned on their cap.
For RV all rows are uniform on $I$; for RK all decrease in the
same inherited priority order. This argument includes $D=\varnothing$.

Until the next acceptance, rejections leave the loads unchanged.
Each has probability $f_j(a)$, for every residual load vector and
every full-day set of size $a$. This proves
\eqref{eq:waitingkernel}. Conditional on $S=s$, the fresh destination
at $s$ is independent of the old residual choices: its row on $I$
is uniform for RV and nonincreasing for RK. The waiting event has
introduced no extra weighting of those loads. Corollary~\ref{cor:allhazard}
therefore gives the asserted mark probabilities on each refinement.
The bound depends only on $K,a,m,C$, so averaging over the refinements
preserves it. With $b=1$ the mark increment is deterministic.
\end{proof}

\begin{theorem}[Equal-capacity greedy matching, all parameters]\label{thm:fullconjecture}
For every $m,C\geq1$, every finite degree sequence
$d_1,\ldots,d_n\in\{0,\ldots,m\}$ independent of the random
neighborhoods, and every integer $k$,
\[
 \Pr(M_\RV\geq k)\geq\Pr(M_\RK\geq k).
\]
In particular this proves Arnosti's conjecture for uniform random
job order and every common capacity $C>2$.
\end{theorem}
\begin{proof}
Fix the degree order. At each $K$ we construct a coupling of the two
state distributions, maintaining
\[
 \sigma_{K,\RV}\leq\sigma_{K,\RK},\qquad
 A_{K,\RV}\leq A_{K,\RK}
 \quad\text{whenever }\sigma_{K,\RK}<\infty.
\]
Level zero is deterministic. Suppose the invariant holds at $K<mC$.
If RK has not reached $K$, it cannot reach a higher level; continue
RV with its own conditional law; RK stays in the cemetery state.
Otherwise both current hitting times are finite.

The waiting kernel in Lemma~\ref{lem:allhistory} is increasing in
the starting index $\sigma_K$ and the full-day count $a$: for
$s>\sigma_K$ its cumulative probability is
\[
 \Pr(S\leq s\mid\sigma_K,A_K)
       =1-\prod_{j=\sigma_K+1}^s f_j(a),
\]
and each $f_j(a)$ increases in $a$. Starting earlier adds factors
between zero and one. A shared uniform variable and the two inverse
cumulative distribution functions therefore give
$S_\RV\leq S_\RK$, including the atom at infinity.

If $S_\RK=\infty$, no new mark comparison is required. Otherwise
both next hits are finite. If the old full-day counts differ, each
new count is its old value plus zero or one, so the order persists
under any coupling. If both old counts equal $a$, both current states
have the same $b=m-a$ and $r=K-aC$. Conditional on their respective
next-hit indices, Lemma~\ref{lem:allhistory} gives Bernoulli
increment probabilities $h_{b,C}(r)$ for RV and $p\geq h_{b,C}(r)$
for RK. A fresh shared uniform variable orders the increments.
This covers $a=0$ with no separate initialization, and it covers
the last acceptance at level $mC$ as well.

Each coupled transition has the actual conditional law of the next
state given the current state as its marginal. The law of total
probability therefore gives the correct state distributions at $K+1$.
This argument does not assume that either original state process is
Markov. Induction reaches $mC$, and since $\{M\geq k\}=\{\sigma_k\leq n\}$, the hitting-time order
proves every tail $1\leq k\leq mC$; the other tails are trivial.
Averaging over a neighborhood-independent random order gives the
source model.
\end{proof}

\noindent\textbf{What is, and is not, established.}
This is a finite written proof with universally quantified parameters,
using a published association theorem. The literature lemma replaces
the previously missing mathematical premise, not merely a numerical
check. The auxiliary Poisson-potential monotonicity remains undecided
and is no longer needed. The proof does not assert domination on each
fixed graph, arbitrary unequal capacities, or a coupling using the same
graph realization. The complete source-model comparison is now formalized in Lean, including
priority relabeling and random-order averaging (Theorem~\ref{thm:leansource}). No independent human review is claimed.

\medskip\noindent\textbf{Lean formalization and dependency record.}
The intermediate assessments below record the status when each stage was
completed. Their then-open formal obligations are all closed by
Theorem~\ref{thm:leansource}.

The first formal stage establishes finite rational covariance identities,
the weighted Chebyshev inequality, a conditional-covariance composition
step, and log-concavity of the capped carrier $1/k!$. It also proves the
size-bias identity obtained by distinguishing a maximal coordinate and
the resulting stochastic increase of the tie count. These statements
quantify over arbitrary finite parameters; they are not finite screens.
The discrete Efron theorem is now formalized in arbitrary finite
dimension, including arbitrary unequal caps and all feasible pairs of
conditioned sums. Its proof includes preservation of log-concavity under
convolution, exact product-law normalization and one induction in the
number of variables. The restricted-domain extension and the comparison
between two values of a labelled maximum are also formalized, including
the exact conditioning identity and positive conditioning masses.
For symmetric observables, exchangeability now identifies the tie-biased
maximum law with the law obtained by pinning one coordinate at the maximum
and deleting it. The stochastic comparison of tie counts is proved for
the actual residual law.
The full finite rational capped-association statement is now a Lean
theorem, proved by induction on the common cap. This bypasses the
previously planned simultaneous maximum/tie regressions: the lower-cap
association itself supplies the pinning comparison needed for the next cap. Exact tie-count strata are identified by one
induction deleting all maxima, with a proved strict-cap endpoint and
normalizing mass. The product-list and capped-vector laws are now
identified exactly, and the labelled-maximum comparison is transferred
to the actual vector law, including a fixed prefix. The normalized
pair-product coefficients, their monotonicity and convexity, and the
factorial balancing and boundary comparisons are also proved over the
rationals. The equivalence between majorization monotonicity and
permutation invariance plus unit balancing is proved in every finite
dimension. Its descent decreases the integer excess by one at each step.
The actual categorical assignment law and its pair-fiber polynomial
decomposition are identified exactly. This proves the ordered capped
boundary probabilities and the uniform-next-choice comparison. The
symmetrized density is now proved majorization-increasing. Applying
association and identifying the exact assignment conditioning proves
conditional concentration and the full categorical saturation bound.
The exact uniform reference, actual neighborhood kernels and normalized
history laws are now formalized too. The refined next-success laws are
transported through a canonical disjoint partition to coarse accepted-count
prefix events, including the waiting, mark and never-success outcomes.
The finite-horizon stopping-state identification and total-probability
identities are now proved as well. The ordered rational coupling of actual
conditional next-acceptance times is also proved, including its never-hit
atom. The joint next-state coupling, its accepted-count induction and all
matching-size tail inequalities in the fixed-priority history model are
now proved. Priority relabeling, explicit graph sampling and source-model
averaging complete the formal proof in Theorem~\ref{thm:leansource}.
The module list and Azure
verification record are in \texttt{formal/README.rst} and
\texttt{formal/verification-azure.json}. The exact source-model transfer is included in the final certificate.

\begin{lemma}[A formalized two-variable Efron step]\label{lem:leanefronpair}
Fix $q,r\geq0$. For $0\leq s\leq q+r$, let $\nu_s$ give weight
$1/(x!y!)$ to pairs $0\leq x\leq q$, $0\leq y\leq r$, $x+y=s$,
and normalize these weights to sum to one. For every coordinatewise
nondecreasing $\phi$ with rational values and every $s\leq t\leq q+r$,
\[
 \mathbb E_{\nu_s}\phi\leq\mathbb E_{\nu_t}\phi.
\]
\end{lemma}
\begin{proof}
It suffices to compare $s$ and $s+1$. Write $p_i=\nu_s(X=i)$ and
$q_i=\nu_{s+1}(X=i)$, and use exclusive cumulative sums
$P(k)=\sum_{i<k}p_i$, $Q(k)=\sum_{i<k}q_i$.
Decreasing adjacent carrier ratios give the cross-multiplied likelihood
inequalities for $p,q$ and for $q$ versus the right shift of $p$.
The finite likelihood-ratio inequality therefore gives
\[
 Q(k)\leq P(k)\leq Q(k+1).
\]
This argument includes zero point masses; it never divides by them.
Define
\[
 \alpha_i=Q(i+1)-P(i),\qquad
 \beta_i=P(i+1)-Q(i+1)\qquad(0\leq i\leq s).
\]
Both are nonnegative. Send mass $\alpha_i$ from $(i,s-i)$ to
$(i,s+1-i)$ and mass $\beta_i$ to $(i+1,s-i)$.
Since $\alpha_i+\beta_i=p_i$, these are the required source masses.
The target masses telescope to $q_i$, including both endpoints.
Each move increases exactly one coordinate, proving the adjacent
expectation inequality. Induction gives the claim for every $s\leq t$.

The Lean theorem is \texttt{capped\_efron\_monotone} in
\texttt{formal/CappedEfron.lean}. The carrier inequalities and positivity
of every feasible normalizing mass are proved there, not assumed.
The more general adjacent-sum theorem in \texttt{EfronTwo.lean} takes
explicit cross-multiplied log-concavity hypotheses. This establishes the
two-variable ingredient of Efron's theorem, not a new result or the
arbitrary-dimensional association theorem.
\end{proof}

\begin{theorem}[The finite-dimensional Efron ingredient in Lean]\label{thm:leanefronproduct}
Let $q_1,\ldots,q_n$ be arbitrary nonnegative integer caps, including the
empty list when $n=0$. For each $0\leq s\leq\sum_iq_i$, give a load vector
$x$ with $0\leq x_i\leq q_i$ and $\sum_i x_i=s$ probability proportional
to $\prod_i1/x_i!$. For every coordinatewise nondecreasing rational-valued
function $\phi$ and all $s\leq t\leq\sum_iq_i$,
\[
 \mathbb E[\phi(X)\mid\textstyle\sum_iX_i=s]
 \leq \mathbb E[\phi(X)\mid\textstyle\sum_iX_i=t].
\]
\end{theorem}
\begin{proof}
First, convolution preserves the carrier inequality needed by the
two-variable result. The finite algebra behind this step is the identity
\[
 2\bigl(\langle A,C\rangle\langle B,D\rangle
       -\langle A,D\rangle\langle B,C\rangle\bigr)
 =\sum_{i,j}(A_iB_j-A_jB_i)(C_iD_j-C_jD_i),
\]
where $\langle A,C\rangle=\sum_iA_iC_i$.
For the appropriate shifted carrier arrays, each product on the right
is nonnegative: both factors are nonnegative for $i\leq j$ and both are
nonpositive for $j\leq i$. Zero extension outside each support makes all
sums finite and handles the endpoints. In Lean this proves preservation
of the inequalities
$a_i a_{j+d}\leq a_{i+d}a_j$ for $i\leq j$, $d\geq0$.
These inequalities are also proved directly for the capped factorial
carriers. Convolution has positive support exactly on the sum of the
two support intervals.

Now induct on the number of variables. Group the remaining variables
into their total $T$ and condition first on $(X_1,T)$. By induction the
conditional mean of $\phi(x,X_2,\ldots,X_n)$ increases with $T$;
it increases with $x$ by pointwise monotonicity. Extend this conditional
mean outside the feasible interval of $T$ by clamping $T$ to its nearest
endpoint. The extension is increasing, and the extra values have zero
weight in the original product law. The two-variable theorem, applied to
$X_1$ and $T$, completes the induction.

The formalization defines the finite product-weight sum recursively and
proves that its value at the constant function $1$ is exactly the
convolution weight. It also proves the normalized conditioning identity
used by the induction, including vanishing terms from infeasible fibers.
Thus the induction concerns the normalized product law itself.
The theorem \texttt{capped\_product\_efron} is in
\texttt{formal/ProductEfron.lean}; the general version permits arbitrary
finite carriers satisfying the explicit cross inequalities and positive
interval support. This is the Efron ingredient, not the conditional
association theorem or the final matching comparison.
\end{proof}

\begin{theorem}[A labelled-maximum comparison in Lean]\label{thm:leanpinned}
Let $n,q,T,a,b$ be nonnegative integers with
\[
 q\leq a\leq b\leq T,\qquad T-a\leq nq.
\]
For $0\leq s\leq nq$, let $\mu_s$ be the probability law on
$\{0,\ldots,q\}^n$ with total $s$ and weights proportional to
$\prod_i1/x_i!$. If $F$ is a rational-valued function increasing under
majorization on the vectors of length $n+1$ and total $T$, then
\[
 \mathbb E_{X\sim\mu_{T-a}}F(a,X)
 \leq \mathbb E_{X\sim\mu_{T-b}}F(b,X).
\]
Equivalently, take the single product law with first-coordinate cap $b$
and all remaining caps $q$, conditioned on total $T$. The displayed
inequality compares the conditional expectations given that the first
coordinate equals $a$ and $b$. Both conditioning events have positive mass.
\end{theorem}
\begin{proof}
For a residual vector $z$ in the box, put
$R(z)=(T-\sum_i z_i,z)$ and $\Phi(z)=F(R(z))$.
Consider $z\leq z'$ coordinatewise in the sum band
$T-b\leq\sum_i z_i\leq\sum_i z'_i\leq T-a$.
The first coordinate of $R(z')$ is at least $a\geq q$.
Every other coordinate of $R(z)$ is at most its counterpart in $R(z')$,
and their total sums agree. These facts imply
\[
 R(z')\prec R(z).
\]
To see this directly, consider a subset of coordinates of $R(z')$.
If it contains the first coordinate, compare complementary sums, all
of which lie in the residual vector. If it does not, and is nonempty,
replace one of its coordinates by the first coordinate; its sum can
only increase, and the first case applies. Thus every subset sum is
bounded by a subset sum of the same cardinality in $R(z)$.
It follows that $\Phi$ is coordinatewise nonincreasing on the sum band.

The Efron theorem applies to functions monotone only on this band.
For completeness, the formalization proves the extension step: bound
an increasing observable on the finite box by $[-B,B]$, use $-B$ below
the band and $B$ above it, and clamp coordinates to the box.
The extended function is globally increasing and agrees with the
original function on the two conditioning layers. Applying this to
$-\Phi$ gives the desired comparison. Exact support identities show
that changing the function off a conditioning layer does not change
its expectation.

Finally, fixing the first coordinate to $k$ contributes the common
factor $1/k!$ to every residual weight. The formal identity for the
restricted raw moment is
\[
 M_{h\otimes H,T}\!\left(\boldsymbol 1_{\{X_1=k\}}\phi(X_2,\ldots)\right)
 =h(k)M_{H,T-k}(\phi).
\]
The same factor occurs in the conditioning mass and cancels. The
formalization proves this identity, its equality to conditioning the
normalized law, and positivity of both masses in the theorem.
The main declarations are \texttt{pinned\_efron\_majorization} and
\texttt{labelled\_maximum\_regression} in
\texttt{formal/PinnedMajorization.lean} and \texttt{formal/PinnedLaw.lean}.
\end{proof}

This is the finite capped-factorial analogue of the comparison used in
Cohen--Sackrowitz, Theorem~2.8. It does not identify the conditional law
given an \emph{unlabelled} maximum: a labelled maximal coordinate
introduces a bias toward outcomes with more maximal coordinates.
The following theorem identifies the tie-biased law exactly. The cap
induction below uses it directly to prove association, bypassing the
earlier plan for simultaneous maximum/tie regressions. The labelled
comparison alone is not the full matching formalization.

\begin{theorem}[Exact tie law and deletion in Lean]\label{thm:leantielaw}
Let $\mu$ be the capped factorial law on $n+1$ coordinates with cap $q$
and total $q+r$, conditioned on its maximum being $q$, and suppose this
event has positive mass. Write $L(X)=\#\{i:X_i=q\}$. Let $\nu$ be the
capped factorial law on $n$ coordinates with cap $q$ and total $r$.
Then $\nu$ has positive normalizing mass and, for every symmetric
rational-valued $F$ on the original load vectors,
\[
 \mathbb E_{Y\sim\nu}F(q,Y)
 =\frac{\mathbb E_{X\sim\mu}[L(X)F(X)]}
        {\mathbb E_{X\sim\mu}L(X)}.
\]
In particular, every nondecreasing rational-valued $\phi$ satisfies
\[
 \mathbb E_{X\sim\mu}\phi(L(X))
 \leq \mathbb E_{Y\sim\nu}\phi(1+L(Y)).
\]
No regression hypothesis about general Schur-increasing functions is
needed for this tie-count comparison.
\end{theorem}
\begin{proof}
Coordinate permutations preserve the capped sample space and its
factorial weights. For a symmetric observable, the raw moment restricted
to $X_i=q$ therefore does not depend on $i$. Summing these moments over
$i$ counts a configuration exactly $L(X)$ times. Hence the moment under
the tie-biased maximum law equals $(n+1)$ times the moment with a fixed
coordinate pinned to $q$. The same identity holds for the masses, so
$n+1$ cancels from the expectation.

Deleting that fixed coordinate is an explicit bijection with the
$n$-coordinate capped space of total $r$. Its weight contributes the
constant factor $1/q!$, which again cancels from the expectation.
Positivity of the original maximum mass implies positivity of the
weighted mass, since $L\geq1$ there. The mass identity and $1/q!>0$
then prove positivity of the residual normalizing mass.
This proves the first identity for symmetric observables; it does not
assert equality of the full labelled-vector laws.

Weighting a finite distribution by its positive count $L$ increases
expectations of increasing functions of $L$, by the weighted Chebyshev
inequality. The deletion bijection changes the count to $1+L(Y)$, giving
the second assertion. All these identities and inequalities are proved
in \texttt{CappedSymmetry.lean}, \texttt{CappedPinnedLaw.lean} and
\texttt{CappedTieComparison.lean}.
\end{proof}

\begin{lemma}[The regression reduction in Lean]\label{lem:leantieregression}
Under the hypotheses of Theorem~\ref{thm:leantielaw}, suppose $F$ is
increasing under majorization and the conditional mean
\[
 \ell\longmapsto\mathbb E_\mu[F(X)\mid L(X)=\ell]
\]
is nondecreasing on the values $\ell$ with positive conditioning mass.
Then
\[
 \mathbb E_\mu F(X)\leq\mathbb E_\nu F(q,Y).
\]
\end{lemma}
\begin{proof}
Majorization monotonicity implies permutation invariance. Partition the
finite distribution by $L$ and apply the weighted Chebyshev inequality
to $\ell$ and the stated conditional mean. The exact finite conditional
expectation identities give $\operatorname{Cov}_\mu(L,F)\geq0$.
Combining this with Theorem~\ref{thm:leantielaw} proves the claim.
The formal partition identities handle zero-mass strata explicitly and
impose no monotonicity requirement there. The final declaration is
\texttt{maximum\_expectation\_le\_residual}.
\end{proof}

This is a proved conditional reduction corresponding to the capped
version of Cohen--Sackrowitz Lemma~5.2. Its monotonicity premise is now
proved by \texttt{capped\_tie\_regression\_all}, following the cap induction
below. The exact conditional distributions needed for that argument are
identified by the following theorem.

\begin{theorem}[Deleting every tied maximum in Lean]\label{thm:leanstrata}
Let $q\geq1$ and $n,r,\ell\geq0$ be integers. Give each load vector of
length $n+\ell$, cap $q$ and total $\ell q+r$ weight $\prod_i1/x_i!$.
Restrict to vectors with exactly $\ell$ coordinates equal to $q$.
Write $Z_{n,q-1,r}$ for the total factorial weight of $n$ coordinates
bounded by $q-1$ and summing to $r$. The restricted total weight is
\[
 Z_{\mathrm{ties}=\ell}
 =\binom{n+\ell}{\ell}\frac{Z_{n,q-1,r}}{(q!)^\ell}.
\]
In particular, it is positive if and only if $Z_{n,q-1,r}>0$.
When these masses are positive, let $\mu_\ell$ and $\nu$ be the respective
normalized laws. For every symmetric rational-valued statistic $F$ on
natural-valued vectors of length $n+\ell$,
\[
 \mathbb E_{\mu_\ell}F(X)
 =\mathbb E_{Y\sim\nu}
   F(\underbrace{q,\ldots,q}_{\ell\ \mathrm{coordinates}},Y).
\]
There is no bound on $n$, $q$, $r$ or $\ell$. The empty residual vector
is included. For cap zero, all coordinates are maximal; that endpoint
is handled separately, without using a negative cap.
\end{theorem}
\begin{proof}
First fix a positive number $k+1$ of maxima. Summing over possible
marked maximal coordinates multiplies every weight in that stratum by
exactly $k+1$. Exchangeability makes each of the $b$ labelled choices
contribute the same raw moment for a symmetric statistic. Deleting the
marked coordinate contributes a factor $1/q!$ and leaves precisely the
stratum with $k$ maxima in $b-1$ coordinates. Thus the one-step raw
moment relation is
\[
 (k+1)M_{b,q,T,k+1}(F)
 =\frac{b}{q!}\,M_{b-1,q,T-q,k}(F(q,\cdot)).
\]
The identical relation for masses cancels the prefactor in normalized
expectations. Symmetry is preserved when the first coordinate is fixed,
so one induction repeats this step $\ell$ times. With no maxima left,
the residual coordinates are all at most $q-1$. An explicit bijection
identifies this event with the smaller-cap sample space and preserves
every factorial weight. This proves the expectation identity.

The same induction on raw masses, using
$(k+1)\binom{n+k+1}{k+1}=(n+k+1)\binom{n+k}{k}$,
gives the displayed normalizing constant. Its prefactor is positive,
so positive conditioning mass is equivalent on both sides, including
infeasible parameter choices where both masses vanish.

The declarations are \texttt{tieStratum\_expectation\_allDelete},
\texttt{tieStratum\_total\_closed} and
\texttt{tieStratum\_all\_mass\_pos\_iff} in
\texttt{formal/CappedStrata.lean}. The one-step identity also handles
statistics defined only on the capped fixed-sum sample space. The
iterated statement uses an ambient symmetric statistic so its domain
remains explicit as dimensions and fixed sums change.
\end{proof}

This closes the capped-factorial tie-stratum identification used in
Cohen--Sackrowitz Lemma~5.1(i). It does not compare different values of
$\ell$: exact distribution identities do not imply monotonicity of
conditional means. The representation bridge and fixed-prefix comparison needed to
combine these results are established below. The cap induction in
Theorem~\ref{thm:leancapassociation} then proves the required regression
and association. The final matching comparison remains unproved in Lean.

\begin{theorem}[The two law representations coincide in Lean]\label{thm:leanrepresentation}
Let $n,q,s\geq0$ be integers and let $\mathcal C_{n,q,s}$ consist of the
natural-valued vectors of length $n$, bounded by $q$ and summing to $s$.
For every rational-valued list statistic $\phi$, the recursively defined
product-list moment satisfies
\[
 M_{[q,\ldots,q],s}(\phi)
 =\sum_{x\in\mathcal C_{n,q,s}}
   \left(\prod_{i=1}^{n}\frac1{x_i!}\right)
   \phi([x_1,\ldots,x_n]).
\]
Consequently the two normalizing constants and the two normalized
expectations agree. The constant is positive exactly when $s\leq nq$.
The raw identity and Lean's totalized rational identity also cover
infeasible sums, where both normalizing masses vanish. No symmetry or
monotonicity of $\phi$ is needed for this representation theorem.
\end{theorem}
\begin{proof}
Sum first over the whole finite box $\{0,\ldots,q\}^n$, using an indicator
for the required total. The head/tail bijection decomposes a vector
into its first coordinate and the remaining vector. Induction on $n$
identifies the box sum with the recursive product-list moment. The
list recurrence sums the first coordinate through $s$, whereas the box
sums it through $q$; in either expression the nonzero contributions are
exactly $0\leq i\leq\min(q,s)$. A finite sum identity equates these two
truncated ranges. The fixed-sum subtype is then precisely the support
of the indicator in the box sum, with the same factorial weights.

Applying this identity to the constant function $1$ identifies the
normalizing mass with the previously proved product carrier, whose
positive support is exactly $0,\ldots,nq$. Division gives equality of
the means, including the explicitly handled zero-mass cases.
The declarations \texttt{productMoment\_eq\_cappedMoment},
\texttt{cappedMean\_eq\_productMean} and
\texttt{capped\_total\_pos\_iff} are in
\texttt{formal/CappedProduct.lean}.
\end{proof}

\begin{corollary}[Efron comparison on the actual capped law]\label{cor:leanvectorpinned}
Let $q\leq a\leq b\leq T$ and $T-a\leq nq$. Let $Y_s$ have the
$n$-coordinate capped factorial law with cap $q$ and total $s$.
For every rational-valued $F$ that increases under majorization on
vectors of length $n+1$ and total $T$,
\[
 \mathbb E F(a,Y_{T-a})\leq \mathbb E F(b,Y_{T-b}).
\]
Both residual laws have positive mass. The comparison also holds after
prepending any fixed number of copies of a fixed value, with the total
in the hypothesis for $F$ adjusted accordingly.
\end{corollary}
\begin{proof}
The representation theorem transfers the proved labelled-maximum
comparison from product lists to the actual capped-vector laws.
Positivity follows from $T-b\leq T-a\leq nq$.
For the prefix extension, prepending the same coordinate to two vectors
preserves their majorization order: a subset of coordinates either
contains the new coordinate or it does not, and the corresponding
subset-sum witness is inherited from the remaining coordinates.
Induction gives the claim for every fixed prefix of equal values.
It also proves that a statistic increasing on the original fixed-total
slice remains increasing on the residual fixed-total slice.
The transferred comparison is
\texttt{capped\_vector\_pinned\_majorization}; its fixed-prefix version is
\texttt{capped\_prefixed\_pinned\_majorization} in
\texttt{formal/CappedPrefix.lean}.
\end{proof}

Together with the exact stratum mass, this also proves that the stratum
of $k$ maxima in $n+k$ coordinates with cap $q+1$ and total $(q+1)k+r$
has positive mass exactly when $r\leq nq$. Likewise, the maximum-$q$
event in $n+1$ coordinates of total $q+r$ has positive mass exactly when
$r\leq nq$. These are proved feasibility criteria, not additional
regression premises. The following induction uses these ingredients to
close the association obligation without a separate maximum regression.

\begin{theorem}[Capped factorial association by cap induction in Lean]\label{thm:leancapassociation}
For every $b,q,T\geq0$, let $w(x)=\prod_i1/x_i!$ on
$\mathcal C_{b,q,T}$. For any rational-valued $F,G$ increasing in
majorization on this finite space,
\[
 \left(\sum_xw(x)F(x)\right)\left(\sum_xw(x)G(x)\right)
 \leq\left(\sum_xw(x)\right)\sum_xw(x)F(x)G(x).
\]
This includes empty spaces. When the mass is positive, it is precisely
$\operatorname{Cov}(F,G)\geq0$ under the normalized capped factorial law.
\end{theorem}
\begin{proof}
Induct on $q$, simultaneously for every dimension and total. The zero-cap
space is empty or consists of the zero vector, so its covariance vanishes.
An increasing statistic on a finite capped space can be extended to the
ambient weak-majorization preorder: at $y$, take the maximum of its
values at capped points below $y$, together with a fixed common lower
bound. This extension is increasing and agrees with the statistic on
its original space. It is invariant under permutations. We may therefore
use ambient statistics when deleting or prepending coordinates.

Assume association for cap $q$ in every dimension. First let $X$ have
$n+1$ coordinates, cap $q$, and total $q+r$, with $r\leq nq$.
The count $L$ of coordinates at $q$ increases in majorization on this
space: the subset consisting of all maximal coordinates must be matched
by an equally large subset with at least the same sum, forcing all its
coordinates to equal $q$. Thus lower-cap association gives
$\operatorname{Cov}(L,F)\geq0$. Biasing by $L$ and deleting the distinguished
maximum identifies the actual residual capped law, giving
\begin{equation}\label{eq:capinduction-pin}
 \mathbb E F(X)\leq
 \frac{\mathbb E[L F(X)]}{\mathbb E L}
 =\mathbb E F(q,Y_r),
\end{equation}
where $Y_r$ has $n$ coordinates, cap $q$ and total $r$.
The feasibility condition proves all needed positive masses.

Now consider cap $q+1$ and condition on the number $K$ of coordinates
at $q+1$. Compare first $K=0$ and $K=1$ in dimension $n+1$, at total
$q+1+r$. The zero-tie law has cap $q$, and its positive mass implies
$r+1\leq nq$. Applying \eqref{eq:capinduction-pin} with residual total
$r+1$, then the proved pinned Efron comparison, gives
\[
 \mathbb E[F\mid K=0]
 \leq\mathbb E F(q,Y_{r+1})
 \leq\mathbb E F(q+1,Y_r)
 =\mathbb E[F\mid K=1].
\]
Both residual vectors here have cap $q$. For any adjacent supported
counts $\ell,\ell+1$, delete $\ell$ copies of $q+1$ using the exact
stratum identity. The remaining comparison is the zero/one comparison
just proved, applied to the statistic with this fixed prefix. Prepending
a fixed prefix preserves majorization.

The supported tie counts form an integer interval: a stratum $K=\ell$
at dimension $b$, cap $q+1$ and total $T$ has positive mass exactly when
\[
 \ell\leq b,\qquad (q+1)\ell\leq T\leq bq+\ell.
\]
Hence the adjacent comparisons prove that $\mathbb E[F\mid K=\ell]$
increases throughout the support. The same holds for $G$.
Within any supported stratum, deleting all maxima leaves a cap-$q$ law;
its conditional covariance is nonnegative by the induction hypothesis.
Finally,
\[
 \operatorname{Cov}(F,G)
 =\mathbb E\operatorname{Cov}(F,G\mid K)
  +\operatorname{Cov}(\mathbb E[F\mid K],\mathbb E[G\mid K])\geq0.
\]
The second term is nonnegative by the finite weighted Chebyshev
inequality. Zero-mass strata impose no conditional hypotheses. This
completes the cap induction for all dimensions and totals.

The final declaration is \texttt{cappedAssociation\_all} in
\texttt{formal/CappedAssociation.lean}. It has no association or
regression premise. Its corollaries prove tie regression for arbitrary
increasing statistics defined only on the capped space, and
\texttt{capped\_maxCount\_covariance}, the association specialization
needed by the matching hazard argument. The finite rational theorem
is now kernel-checked. The categorical next-choice step is closed by
Theorem~\ref{thm:leancategorical}, and the density and association
application by Theorem~\ref{thm:leandensity}. The refined matching law is
proved in Theorem~\ref{thm:leanhistory}, and its residual-row comparison
in Theorem~\ref{thm:leanrows}; accepted-count transitions and coupling remain formalization obligations. No mathematical
novelty or independent kernel reimplementation audit is claimed.
\end{proof}

\begin{theorem}[Pair-product coefficients in Lean]\label{thm:leanpair}
Let $s$ be any finite index set, $t=|s|$, and let $u_i\geq v_i\geq0$
be rational. Put $P(X)=\prod_{i\in s}(u_iX+v_i)$ and
$g_k=[X^k]P/\binom tk$ for $0\leq k\leq t$.
Then $g_k$ is nondecreasing and discretely convex. In particular,
if $a+c=t$ and $a\geq c+2$, then
\begin{align*}
 &(a-1)!(c+1)!\bigl([X^{a-1}]P+[X^{c+1}]P\bigr)\\
 &\hspace{2em}\leq a!c!\bigl([X^a]P+[X^c]P\bigr).
\end{align*}
If $q\leq t\leq2q$, then $[X^{t-q}]P\leq[X^q]P$; division by any
common positive mass preserves this inequality.
\end{theorem}
\begin{proof}
The formal proof uses an exact binomial-basis expansion. Set
\[
 w_A=\prod_{i\in A}(u_i-v_i)\prod_{i\in s\setminus A}v_i\geq0.
\]
Expanding $u_iX+v_i=(u_i-v_i)X+v_i(1+X)$ gives
\[
 P(X)=\sum_{A\subseteq s}w_A X^{|A|}(1+X)^{t-|A|},\qquad
 g_k=\sum_{A\subseteq s}w_A\frac{\binom{k}{|A|}}{\binom{t}{|A|}}.
\]
For fixed $j$, Pascal's identity makes $\binom{k}{j}$ increasing in
$k$ with nonnegative second differences. Hence the same holds for
$g_k$. Increasing first differences give
$g_{a-1}+g_{c+1}\leq g_a+g_c$. Multiplying by $t!$ and using
$a!c!\binom ta=t!$ gives the factorial inequality. Monotonicity and
$\binom t{t-q}=\binom tq$ give the boundary inequality. All denominators
used here are proved positive; zero values among the factors are allowed.

The exact polynomial expansion and its coefficient identity are
\texttt{pairPolynomial\_expansion} and
\texttt{normalizedPairCoeff\_eq\_basis} in
\texttt{formal/PairCoefficients.lean}. The two applications are
\texttt{pair\_factorial\_balancing} and
\texttt{pair\_boundary\_conditional} in
\texttt{formal/PairApplications.lean}.
These are identities and inequalities for actual polynomial coefficients.
The identification with the symmetrized categorical load law is closed
by Theorem~\ref{thm:leandensity}. The refined conditional matching law
is proved in Theorem~\ref{thm:leanhistory}; its residual-row comparison
is proved in Theorem~\ref{thm:leanrows}. The accepted-count transition
comparison remains a separate obligation.
\end{proof}

\begin{theorem}[Integer balancing criterion]\label{thm:leanbalancing}
For any $n\geq0$ and $F:\mathbb N^n\to\mathbb Q$, the following
conditions are equivalent:
\begin{enumerate}
 \item $F(x)\leq F(y)$ whenever $x,y$ have the same total and
 $y$ majorizes $x$.
 \item $F$ is invariant under coordinate permutations and
 \[
 F(y-e_i+e_j)\leq F(y)
 \quad\text{whenever }i\ne j\text{ and }y_i\geq y_j+2.
 \]
\end{enumerate}
The displayed gap condition ensures that the transfer is defined
in $\mathbb N^n$.
\end{theorem}
\begin{proof}
Necessity follows from invariance of majorization under permutations
and from the fact that a balancing move decreases majorization.
For the latter, a subset containing both changed coordinates or neither
has unchanged sum. A subset containing only the donor has smaller sum.
For a subset containing only the recipient, replace that recipient by
the donor to obtain a same-size subset in the original vector whose
sum is at least as large.

For sufficiency, initially sort both vectors decreasingly. The defining
subset-sum inequalities give
$\sum_{i\leq k}x_i\leq\sum_{i\leq k}y_i$ at every prefix, since
the first $k$ coordinates of a decreasing vector maximize the sum of
a $k$-element subset. Keep $x$ sorted throughout; no further sorting of
the changing vector $y$ is needed.

If $x\ne y$, equality of totals gives a first index $j$ with $y_j<x_j$.
For every $i<j$, we have $x_i\leq y_i$. The prefix inequality at $j$
forces some $i<j$ with $x_i<y_i$. As $x_i\geq x_j$, integrality gives
$y_i\geq y_j+2$. Transfer one unit from $i$ to $j$.
Prefixes before $i$ and at or after $j$ have unchanged sums. Each
intermediate prefix previously had a surplus of at least one, since
all its coordinate differences were nonnegative and the difference at
$i$ was positive. Thus every prefix inequality is preserved.

The nonnegative integer
\[
 E(x,y)=\sum_i\max(y_i-x_i,0)
\]
decreases by exactly one. Induction on $E$ terminates at $x=y$ and
composes the inequalities for the successive moves. This proves the
criterion in every finite dimension, with no bound on coordinate values.

The Lean argument is in \texttt{formal/LoadBalancing.lean} and
\texttt{formal/LoadMajorization.lean}, ending with
\texttt{loadLE\_mono\_iff\_balancing}. It uses the same
\texttt{LoadLE} relation as the capped association theorem.
This supplies the order-theoretic passage from the two-coordinate
coefficient inequality to majorization. The exact decomposition of
the symmetrized categorical density into those pair contributions
is proved in Theorem~\ref{thm:leandensity}.
\end{proof}

\begin{theorem}[Capped categorical next choice in Lean]\label{thm:leancategorical}
Let $b\geq1$, and let finitely many independent assignments have rational
probability rows $p_{r1}\geq\cdots\geq p_{rb}\geq0$ with
$\sum_i p_{ri}=1$. Write $N_i$ for the load of bin $i$ and suppose
$E=\{\max_i N_i\leq q\}$ has positive probability. Then
\[
 B_i=\Pr(N_i=q\mid E)\quad\text{satisfies}\quad
 B_1\geq\cdots\geq B_b.
\]
For any independent next choice with probabilities
$v_1\geq\cdots\geq v_b\geq0$, $\sum_i v_i=1$, we have
\[
 \Pr(N_{J_v}=q\mid E)=\sum_i v_iB_i
 \ \geq\ \frac1b\sum_i B_i
 =\Pr(N_{J_{\mathrm{unif}}}=q\mid E).
\]
Thus the uniform next choice minimizes the conditional probability of
selecting a bin at the cap among next rows with this common order.
\end{theorem}
\begin{proof}
For an actual assignment $d$, put $w(d)=\prod_r p_{r,d(r)}$.
The finite product-of-sums identity proves $\sum_d w(d)=1$.
Fix $i<j$. A pattern $o$ records, for each old choice, its destination
outside $\{i,j\}$ or the fact that it belongs to this pair. Let $S_o$
be the latter index set, $t_o=|S_o|$, and let $W_o\geq0$ be the product
of the recorded outside probabilities. Expansion gives the exact identity
\[
 \sum_{\substack{d:\,\pi_{ij}(d)=o\\N_i(d)=a}}w(d)
 =W_o[X^a]\prod_{r\in S_o}(p_{ri}X+p_{rj}).
\]
This is a sum over actual assignments, not an assumed conditional law.
The pattern also fixes all outside loads and $N_i+N_j=t_o$.
Within a fixed pattern, $E\cap\{N_i=q\}$ consists of all assignments
with $N_i=q$ if the outside loads are at most $q$ and
$q\leq t_o\leq2q$, and is empty otherwise. Its mass therefore uses the
coefficient at $q$. The event $E\cap\{N_j=q\}$ uses the same patterns and
the coefficient at $t_o-q$. The boundary inequality of
Theorem~\ref{thm:leanpair} compares these coefficients. Multiply by $W_o$,
sum over patterns, and divide by $\Pr(E)>0$ to obtain $B_j\leq B_i$.
No individual pattern needs positive mass or separate normalization.

Expanding the joint mass of $d$ and the independent next choice gives
$\Pr(N_{J_v}=q\mid E)=\sum_i v_iB_i$. The finite Chebyshev inequality
for the similarly ordered sequences $v$ and $B$ gives
$b\sum_i v_iB_i\geq(\sum_i v_i)(\sum_i B_i)=\sum_i B_i$.

The formal statements are in \texttt{CategoricalLaw.lean},
\texttt{CategoricalPair.lean}, \texttt{CategoricalBoundary.lean} and
\texttt{CategoricalNextChoice.lean} under \texttt{formal/}.
The final declaration is \texttt{uniform\_next\_minimizes\_saturation}.
Its algebraic statement is slightly stronger: it needs nonnegative old
row weights and positive capped mass, but not normalized old rows;
it also needs only order and total one for the next weights.
The separately proved row-normalization identity supplies the probability
interpretation above. This formalizes Lemma~\ref{lem:nextuniform} over
the rationals. The symmetrized density and its association application
are proved in Theorem~\ref{thm:leandensity}; the refined matching-law
identity is proved in Theorem~\ref{thm:leanhistory} and its row comparison
in Theorem~\ref{thm:leanrows}. Accepted-count transitions and final
coupling remain obligations.
\end{proof}

\begin{theorem}[Conditional categorical concentration in Lean]\label{thm:leandensity}
Let $p_{ri}\in\mathbb Q_{\geq0}$ be finitely many categorical probability
rows, all nonincreasing in the same order on $b$ bins. For an assignment
$d$, put $w(d)=\prod_r p_{r,d(r)}$ and let $N(d)$ be its load vector.
Define the actual load mass and its permutation average by
\[
 W_p(x)=\sum_{d:N(d)=x}w(d),\qquad
 \overline W_p(x)=\frac1{b!}\sum_{\pi\in S_b}W_p(x\circ\pi),
 \qquad L_p(x)=\left(\prod_i x_i!\right)\overline W_p(x).
\]
Then $L_p$ is symmetric and increases in majorization on every fixed-total
slice of $\mathbb N^b$. Consequently, for every rational-valued
majorization-increasing $F$ on $\Omega_{b,r,q}$ and every cap event
$E=\{\max_i N_i\leq q\}$ of positive probability,
\[
 \mathbb E_{\mu_{b,r,q}}F\ \leq\ \mathbb E(F(N)\mid E).
\]
Both the density statement and this comparison of the original conditional
assignment law are formalized without an assumed density-monotonicity,
association or concentration premise.
\end{theorem}
\begin{proof}
Fix the other coordinate loads and consider a pair of coordinates with
loads $a,c$, where $a\geq c+2$. The exact pattern identity of
Theorem~\ref{thm:leancategorical} identifies the mass of this load vector
with a sum over patterns fixing every outside destination. The allowed
patterns have exactly $a+c$ assignments in the pair and exactly the
specified outside loads. They are therefore the same for loads
$(a,c)$, $(c,a)$, $(a-1,c+1)$ and $(c+1,a-1)$.

For each such pattern, Theorem~\ref{thm:leanpair} compares the sum of the
two factorial-weighted coefficients before and after balancing. Common
row ordering allows one orientation of the pair; reversing that
orientation exchanges the two summands. The unchanged outside factorials
and pattern weights are nonnegative. Thus, writing
$A_p(x)=(\prod_i x_i!)W_p(x)$, we obtain
\[
 A_p(x^{a-1,c+1})+A_p(x^{c+1,a-1})
 \leq A_p(x^{a,c})+A_p(x^{c,a}).
\]
Here the superscripts replace just the distinguished pair of loads.

Sum this inequality over all permutations. Relabeling commutes with
replacing a pair of coordinates. Swapping the two distinguished
coordinates is a bijection of the permutation sum, so the two terms on
each side have equal sums. Cancel the factor two and divide by $b!>0$.
This proves that balancing decreases $L_p$. Its permutation invariance
is the same reindexing identity; Theorem~\ref{thm:leanbalancing} now gives
majorization monotonicity for every dimension and total.

On the capped fixed-total space, the exact weight identity is
\[
 \left(\prod_i\frac1{x_i!}\right)L_p(x)=\overline W_p(x).
\]
Permutation averaging preserves the total mass and all moments of $F$,
because a majorization-increasing function on a fixed-total space is
symmetric. Theorem~\ref{thm:leancapassociation} therefore gives
\[
 \frac{\sum_x (\prod_i1/x_i!)F(x)L_p(x)}
      {\sum_x (\prod_i1/x_i!)L_p(x)}
 \geq
 \frac{\sum_x (\prod_i1/x_i!)F(x)}{\sum_x\prod_i1/x_i!}.
\]
The left side is identified with the original capped assignment average:
each assignment satisfying the cap maps to exactly one capped load
vector, and summing its weight over that fiber gives $W_p(x)$. In
particular, its denominator is exactly $\Pr(E)>0$. The proof includes
zero-weight fibers; no separate conditional probabilities are assigned
to them.

The final declarations are \texttt{categoricalDensity\_majorization} in
\texttt{formal/CategoricalDensity.lean} and
\texttt{categorical\_conditional\_concentration} in
\texttt{formal/CategoricalConditioning.lean}. As with the next-choice
lemma, the algebraic statement allows nonnegative unnormalized old rows;
row normalization supplies the probability interpretation. The final combination is also formalized. For $b\geq1$ and an independent
next row $v_1\geq\cdots\geq v_b\geq0$ summing to one, put
$B(x)=|\{i:x_i=q\}|$. The declaration
\texttt{categorical\_saturation\_lower\_bound} in
\texttt{formal/CategoricalHazard.lean} proves
\[
 \Pr(N_{J_v}=q\mid E)\ \geq\ \frac1b\mathbb E_{\mu_{b,r,q}}B.
\]
It first identifies the uniform-next hazard with
$\mathbb E(B(N)\mid E)/b$, then applies conditional concentration and the
ordered next-choice comparison. Thus the categorical hazard bound has
no remaining association, density or hazard premise. The exact identity
between uniform categorical rows and the factorial reference law is now
formalized in Theorem~\ref{thm:leanuniform}. The refined conditional-neighborhood law is proved in
Theorem~\ref{thm:leanhistory} and its residual-row comparison in
Theorem~\ref{thm:leanrows}. Accepted-count transitions and final coupling
remain to be formalized.
\end{proof}

\begin{theorem}[Uniform categorical assignments give the reference law]\label{thm:leanuniform}
Let $r\geq0$ labelled assignments independently choose one of $b\geq1$
bins uniformly. Write $N$ for their load vector. For every $x\in\mathbb N^b$
with $\sum_i x_i=r$,
\[
 \Pr(N=x)=\frac{r!}{b^r}\prod_i\frac1{x_i!}.
\]
For a common cap $q\geq0$, let $E=\{\max_i N_i\leq q\}$ and
$Z_{b,r,q}=\sum_{x\in\Omega_{b,r,q}}\prod_i1/x_i!$. Then
\[
 \Pr(E)=\frac{r!}{b^r}Z_{b,r,q},\qquad
 \Pr(E)>0\ \Longleftrightarrow\ r\leq bq.
\]
Whenever $r\leq bq$, every rational-valued function $F$ on the capped
load space satisfies
\[
 \mathbb E(F(N)\mid E)=\mathbb E_{\mu_{b,r,q}}F.
\]
No symmetry or monotonicity of $F$ is required for this identity.
\end{theorem}
\begin{proof}
Expand the multivariate polynomial in two ways:
\[
 (z_1+\cdots+z_b)^r
   =\sum_{d:\{1,\ldots,r\}\to\{1,\ldots,b\}}
       z_1^{N_1(d)}\cdots z_b^{N_b(d)}.
\]
Taking the coefficient at $z^x$ counts exactly the assignments with load
$x$. Mathlib's multinomial coefficient theorem identifies this count as
$r!/\prod_i x_i!$. Every uniform assignment has probability $b^{-r}$,
which gives the claimed load mass. The formal proof first establishes
the more general identity for rows constant across destinations: their
common assignment weight is $\prod_j w_j$.

Restricting to the cap multiplies both the unnormalized observable moment
and the total reference mass by the same positive constant $r!/b^r$.
It cancels in their ratio. The already proved capped-space feasibility
criterion gives $Z_{b,r,q}>0$ exactly when $r\leq bq$. This includes
$r=0$ and $q=0$: a zero cap is feasible precisely for zero assignments.
The algebraic ratio identities also hold on empty capped spaces under
Lean's convention $0/0=0$; their interpretation as conditional
expectations is restricted to positive-mass events.

The coefficient identity is \texttt{constantRowLoadMass\_factorial} in
\texttt{formal/UniformLoadMass.lean}. The exact moments, feasibility and
conditional expectation are proved in \texttt{formal/UniformCappedLaw.lean}.
Combining the identity with Theorem~\ref{thm:leandensity} gives comparisons
directly between two actual assignment laws. If $p$ consists of commonly
ordered nonnegative probability rows, $v$ is an independent next
probability row in the same order, and its cap event $E_p$ has positive
probability, let $J_v$ have row $v$ and let $J_U$ be a uniform next
choice, each independent of its old assignments. Then the corresponding
uniform cap event $E_U$ is also positive and
\[
 \Pr(N^U_{J_U}=q\mid E_U)
 \leq \Pr(N^p_{J_v}=q\mid E_p).
\]
This is \texttt{ordered\_vs\_uniform\_saturation}; the analogous comparison
for every majorization-increasing $F$ is
\texttt{ordered\_vs\_uniform\_concentration}. These theorems assume neither
the reference-law identification nor a saturation inequality. The exact conditional-law identity for the actual matching kernels
is now proved in Theorem~\ref{thm:leanhistory}. Its residual-row comparison
is proved in Theorem~\ref{thm:leanrows}. Accepted-count transitions and
final coupling remain to be formalized.
\end{proof}

\begin{theorem}[The exact refined history law in Lean]\label{thm:leanhistory}
Fix a finite degree sequence $0\leq d_j\leq m$, a common capacity
$C=q+1$, and the priority order on $m$ bins. Both RANDOM-VERTEX and
fixed-priority RANKING have explicit, normalized finite history laws in
Lean. Their one-arrival kernels average the actual response to a
uniform $d_j$-subset of bins: RANDOM-VERTEX chooses uniformly among
available neighbors, and RANKING chooses the least available neighbor.
Rejection is a separate outcome. Every history of nonzero weight has
all loads at most $C$.

For a refinement $\rho$, let $D$ be its prescribed full bins and $J$ its
unspecified accepted jobs. All other arrivals have fixed outcomes $f_j$,
either rejection or a destination in $D$. Histories with these fixed
outcomes and residual destinations outside $D$, all of residual load at
most $q$, are in bijection with assignments
\[
 z:J\longrightarrow I,\qquad I=\{1,\ldots,m\}\setminus D,
 \qquad \max_{i\in I}N_i(z)\leq q.
\]
For either algorithm their weights satisfy exactly
\[
 \Pr(\rho,z)=\kappa_\rho\prod_{j\in J}u_j(z_j)
               \mathbf1_{\{\max_iN_i(z)\leq q\}},
\]
where $u_j(i)$ is the actual one-arrival kernel on the availability set
prescribed by the refinement. If the history fiber has positive mass,
then $\kappa_\rho>0$, its capped assignment mass is positive, and every
conditional observable moment equals that of this capped product law.
No path-factorization or conditional-law premise is assumed.
\end{theorem}
\begin{proof}
Write a history as a function with values in bins or the rejection
symbol. Its prefix loads determine the available set before each arrival.
The history weight is the product of the corresponding actual kernel
values. Each response kernel is nonnegative and sums to one; averaging
over the uniformly chosen neighborhoods preserves both properties. An
induction deleting the last arrival proves that the sum over all history
weights is one. The same decomposition and the fact that unavailable
bins receive zero choice weight prove the capacity bound for every
history of nonzero weight. Degree-zero arrivals and empty histories are
included.

Complete a refinement by an assignment $z$ to its free jobs. For a bin
in $D$, every prefix load is unchanged by $z$. For a bin outside $D$,
all its assignments come from $J$; its prefix load is at most its final
load, hence at most $q$. It is therefore available at every earlier
arrival. This proves that the entire availability path, not just its
endpoint, is fixed by the refinement.

The path product consequently splits into the factors for prescribed
outcomes and those for free destinations. The former give
$\kappa_\rho$; the latter are precisely $\prod_{j\in J}u_j(z_j)$.
The completion map is injective. Conversely, extracting the destinations
of the free jobs from any history in the fiber recovers a unique capped
assignment, and completion recovers the original history. This explicit
bijection proves the exact total-mass and observable-moment identities.
Their common nonnegative factor must be positive whenever the fiber
mass is positive. It can then be canceled, giving the conditional law.
The argument does not assign conditional probabilities to null fibers.

Every history supplies a consistent refinement by taking $D$ to be the
bins whose final loads exceed $q$ and $J$ to be its accepted jobs outside
$D$. For a history of nonzero weight, the capacity bound makes every
load in $D$ exactly $q+1=C$. These fixed contributions make the full-bin
set equal to $D$ throughout its final refined fiber.

The formal modules are \texttt{HistoryRefinement},
\texttt{NeighborhoodChoice}, \texttt{HistoryFactorization} and
\texttt{HistoryFiber}. The final identities are
\texttt{rv\_refined\_history\_law} and
\texttt{ranking\_refined\_history\_law}. The comparison of their actual
residual rows is now formalized in Theorem~\ref{thm:leanrows}, including
RV symmetry, RANKING ordering and the order-preserving relabeling of
residual bins.
The next-success waiting and refined mark laws are proved in
Theorem~\ref{thm:leanwaiting}. Canonical partition and averaging are
proved in Theorem~\ref{thm:leanpartition}. Stopping-state marginal
identities are proved in Theorem~\ref{thm:leanstopping}. Full priority
relabeling and the final coupling remain.
The complete matching comparison is now proved in Theorem~\ref{thm:leansource}.
\end{proof}

\begin{theorem}[The concrete residual saturation comparison in Lean]\label{thm:leanrows}
Fix $C=q+1$, residual bins $I=D^c$ with $b=|I|>0$, and free jobs $J$.
For either algorithm, consider a positive refined history fiber as in
Theorem~\ref{thm:leanhistory}. Normalize the fresh one-arrival kernel on
$I$, conditional on a successful choice; assume that success mass is
positive. Write $H_{\mathrm{RV}}$ and $H_{\mathrm{RK}}$ for the conditional
probabilities that the chosen residual bin currently has load $q$.
Let $H_U(b,q,|J|)$ denote the saturation probability for uniform
assignments conditioned on all loads being at most $q$, followed by an
independent uniform choice. Then Lean proves
\[
 H_{\mathrm{RV}}=H_U(b,q,|J|)\leq H_{\mathrm{RK}}.
\]
The two positive refinements may have different fixed outcomes and degree
sequences, but the direct comparison uses the same $D$ and $J$.
If the prescribed load in every bin of $D$ is $C$, then the actual final
availability set of every history in either fiber is exactly $I$.
No symmetry, row-order or conditional-saturation inequality is assumed.
\end{theorem}
\begin{proof}
For any permutation of the bins, mapping each neighborhood by that
permutation is a bijection of the fixed-size neighborhoods and preserves
its uniform weight. RANDOM-VERTEX is equivariant under this map. Swapping
two available bins therefore proves equality of their choice weights.

For RANKING, let $i\leq j$ be priority labels with $i$ available. If a
neighborhood selects $j$, swap $i$ and $j$ in that neighborhood. The new
neighborhood selects $i$: every unchanged available neighbor is at least
$j$, and the swapped labels cannot create an earlier selected neighbor.
Summing this pointwise implication over the neighborhood bijection gives
$K(j)\leq K(i)$. This argument also covers zero-weight destinations.

Every residual bin is available throughout the fixed partial history,
as proved by its zero fixed load. Thus each actual RV residual row is
constant, while all actual RANKING residual rows decrease in the same
inherited priority order. The same conclusions hold for the next row
after normalization by its positive success mass.

An increasing enumeration identifies $I$ with $\operatorname{Fin}(b)$.
An explicit assignment bijection preserves loads, the cap event and
saturation probability under this reindexing. The old RV rows need not
sum to one on $I$: their row constants multiply all assignment weights
by one common factor, which cancels because the conditioning mass is
positive. Hence the RV law gives precisely $H_U$. The actual RANKING rows
satisfy every hypothesis of Theorem~\ref{thm:leanuniform}, giving the lower
bound. Finally, the history-fiber moment identity, applied to
\[
 G(h)=\sum_{i\in I}v_i\,\mathbf1_{\{N_i(h)=q\}},
\]
identifies these categorical quantities with the asserted history
observables. Consistent full-bin loads make final availability $D^c$.

The final statements are
\texttt{rv\_refined\_saturation\_reference},
\texttt{ranking\_refined\_saturation\_bound}, and
\texttt{rv\_le\_ranking\_refined\_saturation} in
\texttt{HistorySaturation.lean}. The proof is uniform in all finite
parameters. Conditioning on the next-success index is now formalized in
Theorem~\ref{thm:leanwaiting}. Canonical partition and averaging are
proved in Theorem~\ref{thm:leanpartition}; stopping-state marginal
identities are proved in Theorem~\ref{thm:leanstopping}. The final
coupling remains an obligation.
The full matching theorem is still not a Lean certificate.
\end{proof}

\begin{theorem}[Next-success conditioning and acceptance accounting in Lean]\label{thm:leanwaiting}
Let a consistent positive refined history fiber end after arrival $s_0$,
with full-bin set $D$, $a=|D|$, capacity $C=q+1$, and free jobs $J$.
For either algorithm, the exact conditional weight of $t$ rejections
followed by an acceptance is
\[
 \left(\prod_{j=0}^{t-1}f_{s_0+j}(a)\right)
 (1-f_{s_0+t}(a)),\qquad
 f_s(a)=\frac{\binom a{d_s}}{\binom m{d_s}}.
\]
Here future indices are zero based, as in the Lean functions. After a
finite horizon $N$, the no-further-acceptance weight is
$\prod_{j=0}^{N-1}f_{s_0+j}(a)$. These weights sum to one.

On any positive next-success event, conditioning also on its index
preserves the entire old-history distribution in the fiber. The fresh
destination has its actual kernel restricted to $D^c$ and normalized by
its positive success mass, independently of that old history. Therefore
\[
 H_{\mathrm{RV}}(S=s)=H_U(b,q,|J|)
 \leq H_{\mathrm{RK}}(S=s)
\]
in the respective positive refined fibers, with $b=|D^c|>0$.
All null events are excluded explicitly when a conditional law is formed.
Degree-zero arrivals and the no-further-acceptance outcome are included.

Every history in the fiber also satisfies
\[
 K=aC+|J|,\qquad |J|=K-aC.
\]
Appending the first acceptance increases $K$ by exactly one and increases
the full-bin count by exactly
$\mathbf1_{\{\text{chosen bin previously had load }q\}}$.
\end{theorem}
\begin{proof}
A neighborhood rejects exactly when it is contained in $D$. The rejecting
neighborhoods are therefore the $d_s$-subsets of $D$, and their count
proves the displayed rejection probability for both algorithms.

The formal history extension recursively appends $t$ actual rejection
outcomes and then one accepted destination. Rejections preserve every
final load, hence availability. Repeated application of the exact
one-arrival path-product identity gives
\[
 w(h\mathbin{\|}\underbrace{\bot\cdots\bot}_{t}\mathbin{\|}i)
 =w(h)\left(\prod_{j=0}^{t-1}K_{s_0+j}(D^c,\bot)\right)
 K_{s_0+t}(D^c,i).
\]
The extension map is injective in the old history and final destination.
Thus these sums contain no duplicate paths. The formula is derived from
the literal history weights, with availability $D^c$ proved throughout
the fiber; it is not an assumed waiting kernel.

Summing over histories and available destinations factors the event
mass into the old fiber mass, the common rejection product, and the
fresh success mass. Nonnegativity shows that all three factors are
positive whenever the next-success event is positive. They can then be
canceled in any joint observable moment. This proves both preservation
of the old law and its product with the normalized fresh row.
Theorem~\ref{thm:leanrows} applies at the actual next-success index.
For the never-success outcome only the common rejection product remains.
The identity $W_t(1-f_t)=W_t-W_{t+1}$ telescopes, proving the complete
finite waiting normalization, including zero success probabilities.

Finally, the total accepted count equals the sum of all bin loads.
The loads in $D$ sum to $aC$; the residual loads sum to $|J|$. This gives
the exact accounting identity. One accepted destination adds one to its
load, and an available bin becomes full precisely when its previous
load was $q$. The resulting indicator is formally identified with the
actual full-bin count increment in the extended histories.

The main modules are \texttt{WaitingKernel}, \texttt{RejectionKernel},
\texttt{HistoryExtension}, \texttt{WaitingRefinement}, and
\texttt{AcceptedCount}. The preservation identity is
\texttt{fiber\_next\_old\_law}; the refined mark bounds are
\texttt{rv\_next\_success\_saturation\_reference} and
\texttt{ranking\_next\_success\_saturation\_bound}.
The no-further-success normalization is
\texttt{fiber\_next\_or\_never\_total}, and
\texttt{fiber\_next\_mark\_probability} identifies the actual count increment.

The canonical partition and averaging of these laws are proved in
Theorem~\ref{thm:leanpartition}. The finite-horizon stopping-state
identification is proved in Theorem~\ref{thm:leanstopping}. Full
priority relabeling and finite coupling remain.
No Markov property of the original projected process is assumed.
\end{proof}

\begin{theorem}[Canonical refinement and coarse prefix laws in Lean]\label{thm:leanpartition}
Fix a prefix length $s_0$, a capacity $C=q+1$, an accepted count $K$,
and a full-bin count $a<m$. Let $E_{s_0,K,a}$ be the event that the
prefix has exactly $K$ accepted jobs and $a$ full bins and that its last
arrival was accepted. At $s_0=0$ the last-arrival condition is vacuous.
The following statements hold for arbitrary finite parameters.

Every capacity-respecting prefix has a unique canonical refinement.
The event $E_{s_0,K,a}$ is a disjoint union of its consistent canonical
fibers. Restricting the original path sums to capacity-respecting
prefixes discards only zero weight.

After averaging the actual next-success weights over these fibers,
conditional on a positive event $E_{s_0,K,a}$ followed by $t$ rejections
and an acceptance, the full-bin increment has probability
\[
 H_{\mathrm{RV}}=h_{m-a,C}(K-aC)
 \leq H_{\mathrm{RK}}.
\]
The waiting probabilities on a positive $E_{s_0,K,a}$ are the same
explicit binomial products as in Theorem~\ref{thm:leanwaiting}.
The next-success masses plus the no-further-success mass sum exactly
to the old event mass over every finite future horizon.
\end{theorem}
\begin{proof}
For a history $h$, take $D$ to be its full bins and $J$ to be the
accepted jobs sent outside $D$. Record the outcomes outside $J$ in $f$,
and set $f(j)=\bot$ on $J$. The latter convention removes redundant
key data. The formal map from the disjoint union of fibers to legal
histories is bijective: every history supplies this key, and every
member of its fiber recovers exactly the same key. Transporting a
finite sum through the bijection proves the partition identity.

Both the acceptance pattern and the final full-bin set are constant
on a fiber. Thus $E_{s_0,K,a}$ selects whole fibers. Acceptance accounting
gives $|J|=K-aC$, and the complement of $D$ has $m-a$ bins. The exact
uniform reference law therefore has the same numerical value on every
selected fiber, even though their labels and prescribed outcomes differ.

On every positive next-success fiber, Theorem~\ref{thm:leanwaiting}
gives RV equality and the RANKING lower bound. If a fiber has zero mass,
nonnegativity forces every constituent path weight to vanish, so every
observable moment there is zero. Multiply each positive conditional
comparison by its mass, sum over the disjoint fibers, and divide by
the positive total event mass. This proves the coarse comparison
without dividing by the mass of a null fiber. The cap-boundary indicator
is the actual full-bin count increment, by the previously proved
extension identity.

The waiting factor depends only on $|D|=a$, so it factors out of the
same finite sum. Summing the refined next-success/never-success
normalization gives the asserted coarse normalization. Finally, an
illegal old history has zero path weight; the exact extension formula
makes every one of its extended weights zero as well. This proves the
identities with sums over the unrestricted path space.

The main statements are \texttt{history\_refinement\_partition},
\texttt{coarse\_rv\_saturation}, \texttt{coarse\_ranking\_saturation},
\texttt{coarse\_rv\_waiting\_law},
\texttt{coarse\_ranking\_waiting\_law}, and
\texttt{coarse\_next\_or\_never\_total}. The proof uses actual path
weights, with no assumed Markov property or transition comparison.

The explicit finite-horizon stopping-state random variable, its marginal
identities with these prefix events, and its never-reached outcome are
proved in Theorem~\ref{thm:leanstopping}. Full priority relabeling and
the finite ordered coupling induction remain to be formalized. Consequently the full
matching comparison is still not a Lean theorem.
\end{proof}

\begin{theorem}[Actual finite-horizon stopping-state laws in Lean]\label{thm:leanstopping}
For a full history $h$ of $N$ arrivals, define $X_K(h)$ to be the pair
consisting of the arrival index of its $K$th acceptance and the number
of full bins at that time. Put $X_0=(0,0)$ and use a separate value
$\dagger$ if the $K$th acceptance never occurs. Lean constructs this
function with values in the finite space
\[
 \{\dagger\}\ \cup\ \bigl(\{0,\ldots,N\}\times\{0,\ldots,m\}\bigr).
\]
It proves that $X_K=\dagger$ exactly when the total accepted count is
less than $K$, and this outcome is absorbing in the acceptance level.
For $s\leq N$, the event $X_K=(s,a)$ equals the event that the length-$s$
prefix satisfies $E_{s,K,a}$ from Theorem~\ref{thm:leanpartition}.

The actual history weights give a normalized state law $\mu_K$ and a
joint successive-state law $\pi_K$. Their marginals satisfy
\[
 \sum_y\pi_K(x,y)=\mu_K(x),\qquad
 \sum_x\pi_K(x,y)=\mu_{K+1}(y).
\]
For positive $\mu_K(x)$, let $Q_K(x,y)=\pi_K(x,y)/\mu_K(x)$.
Null rows contribute zero, and Lean proves
\[
 \mu_{K+1}(y)=\sum_x\mu_K(x)Q_K(x,y).
\]
These are identities for the original process and require no Markov
assumption. Also,
\[
 \Pr(M\geq K)=1-\mu_K(\dagger).
\]

For these actual stopping variables, the waiting law and the RV equality
and RANKING lower bound from Theorem~\ref{thm:leanpartition} hold
conditional on a positive current state and a positive next-hit event.
A finite next hit is strictly later than the old hit, and its full-bin
count differs by either zero or one. The joint mass of an old finite
state and no next hit equals the proved coarse never-success mass.
\end{theorem}
\begin{proof}
The function \texttt{stoppingState} reads a history recursively. A new
accepted outcome records the current index and full-bin count for its
acceptance level; all other levels retain their previous values.
Induction proves the endpoint characterization, the never-reached
criterion, and persistence of an earlier reached state. Restricting a
history to its first $s$ entries then proves the prefix characterization.
The recorded time is at most $N$ and the full-bin count at most $m$,
so \texttt{finiteStoppingState} has the displayed finite carrier.

For any prefix observable $F$, summing the full path weight times $F$
over all histories gives exactly the prefix-weighted sum. To prove
this, remove the last outcome: its fresh response kernel sums to one.
Induction removes the entire suffix. Applying this identity to the
stopping-state prefix event gives its exact coarse prefix mass.
Pushforward finite sums then prove normalization and both joint
marginal identities. Nonnegative null rows vanish, so division is used
only on positive rows in the total-probability formula.

Suppose the next acceptance occurs $t+1$ arrivals after the old one.
The accepted count has not increased in the intervening prefix. Since
each individual outcome increases it by zero or one, all those outcomes
must be rejections. Thus the history through the next acceptance is
uniquely an old prefix followed by $t$ rejections and an accepted bin.
The proved injectivity and exact converse identify the corresponding
joint-event sum with \texttt{coarseNextMoment}. Marginalizing the suffix
after that next acceptance gives the same identity at the full horizon.
If the next level is never reached, the entire suffix consists of
rejections and the analogous identity gives \texttt{coarseNeverMass}.

Summing over the bounded next full-bin count identifies all its
observable moments. The actual increment is supported on zero and one,
and the indicator of a new full bin equals the prior cap-boundary
indicator. The coarse waiting law and saturation comparison therefore
apply to the actual stopping variables. Finally, the never-reached
criterion identifies matching-size tails with $1-\mu_K(\dagger)$.

The principal statements are \texttt{stoppingMass\_finite},
\texttt{stoppingTransition\_total\_probability},
\texttt{stopping\_joint\_next\_identity},
\texttt{stopping\_joint\_never\_identity},
\texttt{stopping\_rv\_waiting}, \texttt{stopping\_ranking\_waiting},
\texttt{stopping\_rv\_saturation}, and
\texttt{stopping\_ranking\_saturation}.
Theorem~\ref{thm:leanquantile} below proves waiting-law monotonicity
and the ordered coupling of actual conditional next-hit times. The
joint time/mark coupling and fixed-priority matching-tail comparison are
proved in Theorem~\ref{thm:leanjoint}. Priority relabeling and
source-model averaging are proved in Theorem~\ref{thm:leansource}.
\end{proof}

\begin{theorem}[Formal ordered waiting coupling]\label{thm:leanquantile}
Fix a horizon $N$, a degree sequence $0\leq d_j\leq m$, a capacity
$C=q+1$, and an accepted count $K$. Let $(s,a)$ and $(t,b)$ be
positive-mass stopping states for RV and RANKING respectively, with
$s\leq t$ and $a\leq b$. There is a nonnegative rational coupling
of their actual conditional next-acceptance times supported on the
order ``RV no later than RANKING'', including the never-reached atom.
The construction has exactly the conditional marginals of the original
finite history laws. No Markov property is assumed.
\end{theorem}
\begin{proof}
Use indices $0,\ldots,N-1$ for arrivals and index $N$ for no further
acceptance. A finite index $i$ represents the next stopping time $i+1$.
For a current stopping time $s$, set
\[
 F_{s,a}(j)=
 \begin{cases}
 1,&j<s,\\
 \binom{a}{d_j}/\binom{m}{d_j},&j\geq s,
 \end{cases}
 \qquad R_{s,a}(k)=\prod_{j<k}F_{s,a}(j).
\]
The degree bound makes every denominator positive, even for $d_j=0$.
Binomial monotonicity gives
$0\leq F_{s,a}(j)\leq F_{t,b}(j)\leq1$ for all $j$.
This compares absolute arrival indices; it requires no monotonicity
of the degree sequence.

For any failure sequence $F$ in $[0,1]$, define the finite law
\[
 p_F(i)=
 \begin{cases}
 R_F(i)(1-F(i)),&i<N,\\
 R_F(N),&i=N,\\
 0,&i>N.
 \end{cases}
\]
Telescoping proves $\sum_{i<k}p_F(i)=1-R_F(k)$ for $k\leq N$
and $\sum_{i\leq N}p_F(i)=1$. Products preserve the pointwise
failure comparison, so the RV cumulative masses are at least the
RANKING cumulative masses, including the common terminal total.

Here is an explicit rational quantile construction. For any two
nonnegative laws $p,r$ on $0,\ldots,N$, write
$P_i=\sum_{u<i}p(u)$ and $Q_j=\sum_{v<j}r(v)$, and put
\[
 c_{ij}=\min(P_{i+1},Q_{j+1})-\min(P_i,Q_{j+1})
        -\min(P_{i+1},Q_j)+\min(P_i,Q_j).
\]
The rectangle inequality for $\min$ proves $c_{ij}\geq0$.
Telescoping proves $\sum_jc_{ij}=p(i)$ and
$\sum_ic_{ij}=r(j)$. If $Q_k\leq P_k$ for every $k$ and $j<i$,
then $Q_{j+1}\leq P_i$, so all four terms cancel and $c_{ij}=0$.
This proves ordered support without real-valued random variables,
division by point masses, or an external coupling theorem.

To identify the marginals with the actual algorithm laws, the earlier
next-hit identity gives the atom at each $i=s+u<N$ as
\texttt{countWaitingFactor}. Hits before the old stopping time have
zero joint mass. The exact rejection-extension identity and finite
refinement partition give the never-hit joint mass as the old-state
mass times $R_{s,a}(N)$. Dividing by the positive old-state mass
therefore identifies every atom, including the cemetery atom.
These are the statements \texttt{stopping\_rv\_time\_mass} and
\texttt{stopping\_ranking\_time\_mass}; applying the rational
construction proves \texttt{actual\_waiting\_ordered\_coupling}.

The Bernoulli part also has an explicit rational primitive. For
$0\leq u\leq v\leq1$, its nonzero masses on $(0,0),(0,1),(1,1)$
are $1-v,v-u,u$. Both Bernoulli marginals and ordered support are
proved. If old full-bin counts differ, increments in $\{0,1\}$
preserve their order under any coupling.
\end{proof}

\noindent\textbf{Historical stage assessment.}
This closes the waiting-law order and its coupling with the actual
conditional marginals for all finite parameters. The joint time/mark coupling and its induction are now formalized in
Theorem~\ref{thm:leanjoint}. The source-model transfer is proved separately
in Theorem~\ref{thm:leansource}; the generic Bernoulli table alone would
not certify it.

\begin{theorem}[Formal next-state coupling and fixed-priority comparison]\label{thm:leanjoint}
In the concrete finite uniform-neighborhood history model, fix any
number of bins $m$, horizon $N$, capacity $q+1$, and deterministic degree
sequence $d_j\leq m$. For every accepted count $K$, the actual RV and
fixed-priority RANKING stopping-state distributions have an ordered
rational coupling. In particular,
\[
 \sum_h W_{\RK}(h)\,\mathbf 1\{K\leq\operatorname{acc}(h)\}
 \ \leq\
 \sum_h W_{\RV}(h)\,\mathbf 1\{K\leq\operatorname{acc}(h)\}.
\]
This holds for all $K,N,m,q$, with no finite parameter cutoff and no
Markov assumption on the projected stopping-state process.
\end{theorem}
\begin{proof}
For each positive finite current state, the next full-bin count is
supported on $a,a+1$. If $D$ is the joint mass of a specified next hit
and $C$ is its mass with a new full bin, every count observable satisfies
\[
 \sum_b w_b F(b)=(D-C)F(a)+CF(a+1)
   =D\sum_{x\in\{0,1\}}\operatorname{Ber}_{C/D}(x)F(a+x).
\]
The identity includes $D=0$: nonnegative null rows vanish. The parameter
$C/D$ lies in $[0,1]$ even on a null row, where it equals zero.
This proves the entire conditional mark law, not only its mean.

Take the ordered waiting coupling $c_{ij}$ from
Theorem~\ref{thm:leanquantile}. At tied old counts and finite next hits,
attach the Bernoulli table with masses $1-v,v-u,u$ on
$(0,0),(0,1),(1,1)$. On a positive time pair, both time marginals are
positive, so the proved actual saturation comparison gives $u\leq v$.
Null time pairs contribute zero without imposing a comparison on them.
At unequal old counts, or when a next hit is never reached, use the
product Bernoulli table. Summing either bit recovers its exact marginal;
summing either time then recovers the full marked marginal.

Decode a marked outcome $(i,x)$ as $(i+1,a+x)$ for $i<N$ and as the
cemetery state for $i=N$. Splitting the original next-state sum into
its cemetery atom and finite times proves equality of \emph{every}
observable moment with this decoded marked law. Time zero has zero
next-hit mass. A bounded-state retraction sends out-of-range states to
the cemetery state; it is the identity on genuine states and preserves
the order. Pushing the rational table through it therefore gives a
coupling on the actual finite state space, with both exact conditional
transition marginals. If the current RANKING state is already cemetery,
its next state stays there and RV follows its own transition law.
This proves \texttt{actual\_ordered\_transition\_coupling}.

A generic finite composition lemma now finishes the induction. For each
positive entry of the old coupling, choose the proved next-state coupling
and multiply it by that old entry. Sum over all old state pairs. Zero
entries need no conditional coupling. Nonnegativity and ordered support
are preserved; summing a new coordinate and then an old coordinate gives
exactly the law-of-total-probability update for each marginal.
Level zero is the common deterministic state $(0,0)$, so induction proves
\texttt{stopping\_laws\_ordered\_coupling} for every natural $K$.

In an ordered state coupling, RV can be cemetery only when RANKING is
cemetery. Thus its cemetery mass is no larger. The previously proved
matching-tail identity converts this to
\texttt{fixed\_priority\_matching\_tail\_comparison}.
\end{proof}

\noindent\textbf{Historical stage assessment.}
The full coupling induction and every matching-size tail comparison are
now proved in the concrete fixed-priority history model. The following
theorem supplies the exact source-model transfer and closes the formal proof.

\begin{theorem}[Formal source-model comparison]\label{thm:leansource}
Let $n,m,q$ be arbitrary natural numbers and fix degrees
$d_j\leq m$ for all $j\in\{0,\ldots,n-1\}$. Independently sample
uniform neighborhoods of these degrees, a uniform arrival permutation,
and, for RANKING, one uniform bin-priority permutation. Give every bin
capacity $q+1$. For every natural threshold $K$,
\[
 \Pr(M_{\RK}\geq K)\leq\Pr(M_{\RV}\geq K).
\]
The complete statement is \texttt{source\_matching\_tail\_comparison}
in \texttt{formal/SourceModel.lean}. Its only data restriction is
$d_j\leq m$; it assumes no comparison, association or coupling lemma.
All positive integer capacities are covered by $q+1$.
\end{theorem}
\begin{proof}
A priority permutation $e$ maps original bin labels to numerical ranks.
The response \texttt{priorityResponse} selects $k\in A\cap N$ precisely
when $e(k)\leq e(i)$ for every $i\in A\cap N$. Relabel each accepted
history outcome by $e$, leaving rejections unchanged. The proved history
bijection preserves the number accepted, carries each bin load to its
renamed bin, and sends availability sets to their images under $e$.
Uniform neighborhoods are invariant under this bijection. Consequently
an arbitrary-priority history weight equals the numerical-priority weight
of the relabeled history. Summing over all histories proves equality of
matching-size tails for every fixed priority.

Next define the explicit mass of a neighborhood $S$ of job $j$ as
\[
 u_j(S)=\frac{\mathbf 1\{|S|=d_j\}}{\binom m{d_j}}.
\]
These masses are nonnegative and sum to one, including $d_j=0$.
The original job-labeled graph has mass $\prod_j u_j(G_j)$.
For a fixed arrival permutation $\sigma$, its algorithmic history weight
is the product of response probabilities using $G_{\sigma(j)}$ and the
actual availability set before arrival $j$. Reindexing graph profiles
by $\sigma$ is a bijection preserving the product mass, with degrees
reordered to $d_{\sigma(j)}$. Expanding a finite product of sums then
shows that averaging the graph gives exactly the already verified
history kernel. This proves the source identification algebraically,
without assuming independence of the evolving availability sets.

Finally, \texttt{sourceRVHistory} averages the explicit graph/history law
over the arrival permutation. \texttt{sourceRankingHistory} also averages
over a single priority permutation outside the entire history product.
Both laws are nonnegative and sum to one. RANDOM-VERTEX is written in
its sequential form, choosing uniformly from the available neighbors;
this is the same rule as taking the first such neighbor in a fresh
uniform job-specific preference permutation. Finite averaging commutes
with the matching-tail sum and preserves inequalities. Apply the
arbitrary-priority version of Theorem~\ref{thm:leanjoint} for each arrival
permutation and average. This is the claimed source-model inequality.
\end{proof}

\noindent\textbf{Stage assessment.}
The exact source-model transfer is now proved, including arbitrary-priority
relabeling, explicit independent graph sampling, the change of arrival
indices, normalization, and both finite permutation averages. This closes
the formal proof of the stated equal-capacity comparison. It does not
extend the statement to unequal capacities or domination on each fixed
graph. No independent human review or novelty claim is made.

\medskip\noindent\textbf{Finite audit of the matching bridge.}
An Azure run on 5 October 2026 checked the transition formulas against
direct enumeration of fresh neighborhoods in 1,536 rows. It checked
16,541 refined residual-load distributions and 26,943 conditional
mark hazards, including 1,365 cases with no full day and 7,818 with
one available day. All 1,439 matching-tail comparisons across 185
degree sequences passed. These checks used exact rational arithmetic.
They test the conditioning and transition formulas; they do not prove
their unbounded versions or formally verify the published association
theorem. The complete proof is the argument above, which no longer
depends on the first-saturation theorem or on any historical partial
matching regime below.

\medskip\noindent\textbf{Reading the research record below.}
The following sections retain the intermediate proofs and their stage
assessments. Statements about an unresolved regime or a next research
gate in those assessments describe the status \emph{before}
Section~\ref{sec:literatureclosure}; the current result is
Theorem~\ref{thm:fullconjecture}.

\section{Two days: arbitrary capacity and number of jobs}
For two days a positive degree is either one (a constrained job) or two
(a fully flexible job). Fix any deterministic sequence of these degrees,
independent of the random neighborhoods. The following theorem is therefore
stronger than the comparison averaged over random job order.

\begin{theorem}[Two days]
For $m=2$, any $C\geq1$, and any finite fixed sequence of degrees in
$\{1,2\}$, RANDOM-VERTEX stochastically dominates RANKING in the number
of jobs accepted. Consequently the original random-order conjecture holds
for two days and all capacities.
\end{theorem}
\begin{proof}
Let $T$ be the first step at which a day becomes full; put $T=\infty$
if no day becomes full during the given sequence. Until that step every
job succeeds. We first show
\begin{equation}\label{eq:survival}
 \Pr(T_\RV>t)\geq\Pr(T_\RK>t)\quad\text{for every prefix length }t.
\end{equation}
Fix a prefix containing $f$ flexible jobs and $s=t-f$ constrained jobs.
Extend its unconstrained preferred-day choices hypothetically, ignoring
capacity. No saturation has occurred by time $t$ precisely when both
day counts are at most $C-1$. In terms of the count for day A this is
\[
 L=t-C+1\ \leq\ N_A(t)\ \leq\ U=C-1.
\]
This endpoint condition suffices because each count only increases.
Under RANDOM-VERTEX, all preferred-day choices are independent fair
coins. Under RANKING, fix day A as the higher-ranked day; relabeling
symmetry justifies this. All $f$ flexible jobs choose A, while the
$s$ constrained choices are still independent fair coins.

Let $B\sim\operatorname{Bin}(s,1/2)$ and set
$g(k)=\Pr(L-k\leq B\leq U-k)$ for $0\leq k\leq f$.
If $L>U$ both survival probabilities vanish. Otherwise
$g(k)=g(f-k)$, since the binomial law is symmetric and $L+U=s+f$.
Moreover, writing $b_s(x)=\Pr(B=x)$, with zero outside $0,\ldots,s$,
\[
 g(k+1)-g(k)=b_s(L-k-1)-b_s(U-k)\geq0
 \qquad\left(0\leq k\leq\frac{f-1}{2}\right).
\]
Indeed, the midpoint of the two arguments is
$s/2+(f-2k-1)/2\geq s/2$, so the lower argument is at least as close
to the mode $s/2$ as the upper argument. Symmetry and unimodality of the
binomial mass function give the inequality. Thus $g(k)\geq g(f)$ for
every $k$. The RANKING survival probability is $g(f)$, whereas the
RANDOM-VERTEX survival probability is $\mathbb E[g(K)]$ with
$K\sim\operatorname{Bin}(f,1/2)$ independent of $B$. This proves
\eqref{eq:survival}.

It remains to justify that later first saturation improves the terminal
matching, not just a prefix statistic. Let $n$ be the total number of
jobs and $\tau=\min(T,n)$. Conditional on $\tau=t<n$, the first $t$
jobs have all succeeded and exactly one day has filled. In the remaining
day every future flexible job succeeds while there is space; each future
constrained job succeeds with probability $1/2$, independently. The
random feasible days of future jobs are independent of the saturation
history. Which day remains does not change their joint success law.

Take independent variables $X_j$, equal to one for a flexible job and
fair Bernoulli variables for constrained jobs, independently of $\tau$.
The terminal matching-size law is therefore the law of
\begin{equation}\label{eq:continuation}
 F(\tau,X)=\min\left\{2C,\ \tau+\sum_{j=\tau+1}^n X_j\right\}.
\end{equation}
For $\tau=n$ this still holds: all $n$ jobs succeeded. For fixed $X$,
the expression inside the minimum increases by $1-X_{t+1}\geq0$
when $t$ increases to $t+1$. Thus $F(t,X)$ is nondecreasing in $t$.
Equation~\eqref{eq:survival} implies stochastic dominance of
$\tau_\RV$ over $\tau_\RK$. Couple these finite random variables by
the same uniform quantile so that $\tau_\RV\geq\tau_\RK$, and use
the same independent $X$ in \eqref{eq:continuation}. The resulting
matching sizes satisfy $M_\RV\geq M_\RK$ pointwise in this coupling,
which proves every tail inequality. Finally average over the original
random job order. Zero-degree jobs, if any, can simply be deleted.
\end{proof}

\noindent\textbf{Why this does not yet prove the general conjecture.}
With two days, after the first fills, only one destination remains and
its future behavior is determined by the number of remaining slots.
With three or more days, the residual capacities can have different
shapes and remain correlated with the shared priority order. A scalar
first-saturation time no longer specifies the continuation law.

\section{First saturation for any number of days}
The binomial argument extends to any number of days by averaging adjacent
preference probabilities. The following elementary form makes the necessary
common ordering of the preferences explicit.

\begin{lemma}[A central interval for coins biased in the same direction]
Let $Y_1,\ldots,Y_N$ be independent Bernoulli variables with parameters
$q_j\in[1/2,1]$, and let $B\sim\operatorname{Bin}(N,1/2)$.
For any nonnegative integer $b$,
\[
 \Pr\left(N-b\leq\sum_jY_j\leq b\right)
 \leq \Pr(N-b\leq B\leq b).
\]
\end{lemma}
\begin{proof}
The assertion is immediate if $N>2b$ or $b\geq N$.
Otherwise fix all parameters except $q_i$ and let $R=\sum_{j\ne i}Y_j$.
The derivative of the central-interval probability with respect to $q_i$ is
\[
 \Pr(R=N-b-1)-\Pr(R=b).
\]
Put $r=N-1$, $a=N-b-1$. Then $0\leq a\leq b\leq r$ and $a+b=r$.
First suppose all other parameters are below one, and put
$w_j=q_j/(1-q_j)\geq1$. Up to the common factor
$\prod_{j\ne i}(1-q_j)$, the mass $\Pr(R=k)$ is the elementary symmetric
polynomial $e_k(w)$. Its normalized value $e_k(w)/\binom rk$ is the
mean product over a uniform $k$-subset. Choose a uniform permutation of
the $r$ indices: the product of its first $k$ entries is nondecreasing
in $k$ because each $w_j\geq1$. Taking expectations shows that these
normalized values are nondecreasing. Since $\binom ra=\binom rb$,
we obtain $\Pr(R=a)\leq\Pr(R=b)$. Parameters equal to one follow by
continuity. Thus lowering each $q_i$ to $1/2$ can only increase the
central-interval probability.
\end{proof}

\begin{lemma}[Aligned categorical choices]\label{lem:aligned}
Let independent choices $Z_1,\ldots,Z_N$ take values in $\{1,\ldots,m\}$,
with probabilities $p_{ji}=\Pr(Z_j=i)$ satisfying
$p_{j1}\geq\cdots\geq p_{jm}$ for every $j$.
For every common nonnegative occupancy bound $b$, the probability that
all $m$ occupancy counts are at most $b$ is no larger than for
$N$ independent uniform choices.
\end{lemma}
\begin{proof}
Choose adjacent columns $i,i+1$ and replace both entries of each row by
their average. This preserves the ordering of every row, all probabilities
outside the pair, and each row's total probability of choosing the pair.
Condition on which choices fall in the pair and on the exact destinations
of all other choices. If an outside count exceeds $b$, both conditional
probabilities are zero. Otherwise, let $K$ choices fall in the pair.
Their choices within it are independent Bernoulli variables with parameters
$p_{ji}/(p_{ji}+p_{j,i+1})\geq1/2$. Zero-denominator rows cannot occur
among these $K$ choices. Both pair counts are at most $b$ precisely when
their sum of Bernoulli successes lies in $[K-b,b]$. After averaging,
the same conditional choices are fair. The preceding lemma therefore
shows that averaging cannot decrease the required occupancy probability.

Repeat these averages cyclically over all adjacent pairs. The rows converge
to $(1/m,\ldots,1/m)$: each update decreases a row's sum of squared entries
by half the squared difference of the averaged entries. The total decrease
is finite, so all differences at their updates tend to zero. Each coordinate
change tends to zero as well; consequently all adjacent differences at
the start of a cycle tend to zero. The row sum remains one, which gives
the asserted limit. The occupancy probability is a polynomial in the
entries and hence continuous. Passing to the limit proves the claim.
\end{proof}

\begin{theorem}[First saturation, without a bound on $m,C,n$]\label{thm:first}
For every $m\geq2$, common capacity $C\geq1$, and fixed finite sequence
of positive degrees independent of the neighborhoods, the first saturation
times satisfy
\[
 \Pr(T_\RV>t)\geq\Pr(T_\RK>t)\qquad\text{for every prefix length }t.
\]
\end{theorem}
\begin{proof}
Ignore capacities when generating each job's preferred day. Until the first
saturation, these hypothetical choices agree with the actual assignments.
Under RANDOM-VERTEX the preferred days are independent uniform choices.
Under RANKING, conditional on the single priority order and after relabeling
days in that order, job $j$ chooses rank $r$ with probability
\[
 p_{d_j}(r)=\frac{\binom{m-r}{d_j-1}}{\binom m{d_j}},
\]
where the numerator is zero if its lower index exceeds its upper index.
These probabilities are nonincreasing in $r$, and the preferred choices
are independent because the neighborhoods are independent. The event
$T>t$ is exactly the event that every preferred-day occupancy in the
first $t$ jobs is at most $C-1$. Occupancies only increase, so checking
the final prefix counts checks all earlier steps as well.
Apply Lemma~\ref{lem:aligned} with $b=C-1$.
\end{proof}

\section{The matching-size comparison through threshold $2C$}
\begin{theorem}[Stop after $2C$ successful assignments]\label{thm:2C}
For every $m\geq2$, common capacity $C\geq1$, and fixed sequence of
$n$ positive degrees independent of the neighborhoods,
\[
 \Pr(M_\RV\geq k)\geq\Pr(M_\RK\geq k)
 \qquad\text{for every integer }k\leq2C.
\]
Equivalently, $\min(M_\RV,2C)$ stochastically dominates
$\min(M_\RK,2C)$. In particular, the full matching-size comparison holds
when $n\leq2C$, or when $m=2$ with arbitrary $n$.
All these statements also hold under uniformly random job order.
For $m=1$ both algorithms always agree.
\end{theorem}
\begin{proof}
Stop counting successful assignments when the count reaches $2C$.
Before that stopping point, two full days are impossible: they would
already contain $2C$ accepted jobs. Thus every arrival relevant to the
stopped count encounters either zero or one full day.

Let $\tau=\min(T,n)$. The first $\tau$ jobs all succeed. If $\tau=t<n$,
exactly one day is full at that point. When $t<2C$, this remains the
only full day until the stopped count reaches $2C$. Each future degree-one
job avoids that particular day with probability $(m-1)/m$, independently
of the history and the other future neighborhoods. Each job of degree
at least two has an available neighbor while the stopped count is below
$2C$, and therefore succeeds. The identity of the first full day does
not change this joint success law. After reaching $2C$, any subsequent
outcomes are irrelevant to the stopped count.

Consequently, for independent variables $X_j$, independent also of $\tau$,
with $X_j\sim\operatorname{Bernoulli}((m-1)/m)$ for degree-one jobs and
$X_j=1$ otherwise, the stopped matching-size law is
\begin{equation}\label{eq:2C-mixture}
 \min(M,2C)\ \overset{\mathrm{law}}=\ H(\tau,X)
 :=\min\left\{2C,\ \tau+\sum_{j=\tau+1}^nX_j\right\}.
\end{equation}
For $\tau\geq2C$ the stopped count is already $2C$; for $\tau=n$
all $n$ jobs succeed. These observations also cover the boundary cases
in the formula. For each fixed $X$, the expression inside the minimum
increases by $1-X_{t+1}\geq0$ when $t$ increases to $t+1$.
Theorem~\ref{thm:first} permits
a quantile coupling $\tau_\RV\geq\tau_\RK$. Using the same independent
$X$ proves the stopped-count dominance and hence every tail comparison
through $2C$. When $n\leq2C$ or $m=2$, no matching can exceed $2C$,
so this proves the entire matching-size comparison. Average over job
orders to recover the source model; degree-zero jobs can be deleted.
\end{proof}

\noindent\textbf{Why the cap cannot simply be omitted.}
With $m=3$, $C=3$, take six fully flexible jobs followed by one degree-two
job, in that fixed order. RANKING fills two days with the first six jobs,
so the last job is rejected with probability $1/3$.
Under RANDOM-VERTEX, fix an unordered pair of days. The probability that
exactly this pair is full after the first six jobs is
\[
 \sum_{t=3}^{5}
 \frac{2\binom{t-1}{2}}{3^t2^{6-t}}
 =\frac1{108}+\frac1{54}+\frac2{81}=\frac{17}{324}.
\]
Here $t$ is the first saturation time: the last job in that prefix fills
one of the two specified days, and all remaining flexible jobs must choose
the other specified day. There are three pairs; conditional on two full
days, the last degree-two job chooses exactly those days with probability
$1/3$. Therefore the probabilities of accepting all seven jobs are
\[
 \Pr(M_\RV=7)=\frac{307}{324},\qquad
 \Pr(M_\RK=7)=\frac23.
\]
Omitting the cap from \eqref{eq:2C-mixture} would wrongly predict
probability one for both, since there are no degree-one jobs. The stopped
formula is correct: both algorithms certainly accept at least six jobs.
This refutes the uncapped extension, not the matching conjecture;
the comparison still points in the conjectured direction. Also, this
example uses a fixed order, not the random-order numerical probabilities
for the same multiset.

\noindent\textbf{Later saturation times need not have the same ordering.}
Another tempting extension is that RANDOM-VERTEX delays every successive
day becoming full. This fails even from equal initial capacities.
Take $m=3$, $C=3$, with fixed degree sequence $(3,3,3,1,1,1)$.
RANKING fills its first day with the first three jobs. Two days are full
at the end exactly when the three constrained jobs all choose the same
one of the other two days, with probability $2/27$.

Under RANDOM-VERTEX, after the first three jobs the load shapes
$(3,0,0)$, $(2,1,0)$, $(1,1,1)$ have probabilities $1/9$, $2/3$, $2/9$.
The conditional probabilities of having two full days after the remaining
three constrained jobs are respectively $2/27$, $1/9$, and zero.
In the middle case the constrained jobs must send one job to the day
of load two and two to the day of load one, in any of three orders.
Consequently, writing $T_2$ for the second saturation time,
\[
 \Pr(T_{2,\RV}\leq6)=\frac19\frac2{27}+\frac23\frac19
 =\frac{20}{243}>\frac2{27}=\Pr(T_{2,\RK}\leq6).
\]
This concerns a stronger fixed-order auxiliary claim, not a reversal of
the original random-order matching comparison. More full days can reflect
more accepted jobs rather than worse performance. Thus a general proof
cannot merely repeat the first-saturation ordering for all later times.

\section{An explicit formula at the first rejection threshold}
Zero-degree jobs are always rejected and can be removed without changing
the relative order or outcomes of the other jobs. Let $n$ denote the number
of positive-degree jobs and $h_1$ the number of them having degree one.
Adopt $\binom ab=0$ when $b>a\geq0$.

\begin{lemma}[Products of decreasing nonnegative sequences]\label{lem:products}
If $f_1,\ldots,f_\ell$ are nonnegative, nonincreasing functions on
$\{1,\ldots,m\}$, then
\[
 \frac1m\sum_{r=1}^m\prod_{j=1}^\ell f_j(r)
 \;\geq\;\prod_{j=1}^\ell\left(\frac1m\sum_{r=1}^m f_j(r)\right).
\]
\end{lemma}
\begin{proof}
For two such functions $f,g$, the difference between the two sides is
\[
 \frac{1}{2m^2}\sum_{r,s=1}^m(f(r)-f(s))(g(r)-g(s))\geq0.
\]
A product of nonnegative nonincreasing functions is again nonnegative
and nonincreasing, so induction proves the assertion. The empty product
is one.
\end{proof}

\begin{theorem}[At most $C+1$ positive-degree jobs]
For every $m,C$ and every degree vector with $n\leq C+1$, the full
stochastic-dominance assertion holds. For $n\leq C$ both algorithms
accept all $n$ positive-degree jobs. For $n=C+1$, both matching sizes
belong to $\{C,C+1\}$ and
\[
 \Pr(M_\RV=C)=\frac{h_1}{C+1}\,m^{-C}.
\]
If $h_1=0$, neither algorithm rejects a job. If $h_1>0$, remove any one
degree-one job and define
\[
 p_d(r)=\frac{\binom{m-r}{d-1}}{\binom md},\quad
 Q=\sum_{r=1}^m\prod_{j\ne j_*}p_{d_j}(r),\qquad d_{j_*}=1.
\]
Then
\[
 \Pr(M_\RK=C)=\frac{h_1}{(C+1)m}Q
 \;\geq\;\frac{h_1}{C+1}m^{-C}.
\]
\end{theorem}
\begin{proof}
Before $C$ jobs have been assigned, no day can be full. Thus the first
$C$ positive-degree jobs all succeed. At that point a day is full if and
only if all $C$ assignments went to that one day. At most one day is full.
Consequently the $(C+1)$st job fails precisely when it has degree one
and its unique feasible day is that full day.

The last job has degree one with probability $h_1/(C+1)$. Conditional
on its identity, its feasible day is uniform and independent of all
earlier jobs. The earlier degrees form the remaining multiset; their
order does not matter for the product calculations below.

Under RANDOM-VERTEX, each of the first $C$ choices is uniform among all
$m$ days, regardless of its positive degree: none is yet full, and the
neighborhood and individual preference randomness are fresh. The choices
are independent. The probability of all $C$ choices coinciding is
$m^{1-C}$. Multiplication by the probability $1/m$ that the last
singleton chooses that day gives the displayed rejection probability.

For RANKING fix its priority order. By relabeling symmetry we may number
days in this order. A degree-$d$ job chooses day $r$ if its neighborhood
contains $r$ and avoids $1,\ldots,r-1$, giving exactly $p_d(r)$.
The neighborhoods of distinct jobs are independent, so the probability
that all first $C$ choices coincide is $Q$. Each $p_d$ is nonnegative
and nonincreasing, and $\sum_r p_d(r)=1$. Applying the lemma to the
$C$ factors shows $Q/m\geq m^{-C}$. This proves the required comparison
of rejection probabilities. Since the only possible outcomes are $C$
and $C+1$, that single comparison establishes every tail inequality.
\end{proof}

\begin{corollary}[One constrained job and $C$ fully flexible jobs]
For $m\geq2$ and degrees $(1,m,\ldots,m)$, with $C$ copies of $m$,
\[
 \Pr(M_\RV=C)=\frac{1}{(C+1)m^C},\qquad
 \Pr(M_\RK=C)=\frac{1}{(C+1)m}.
\]
\end{corollary}
\begin{proof}
Every flexible job chooses the first-ranked day under RANKING, so
$Q=1$. Substitute $h_1=1$ in the theorem.
\end{proof}

For $m=2,C=3$, the degree vector is $(1,2,2,2)$.
The probabilities of accepting all four jobs are respectively
$31/32$ and $7/8$. Under RANKING, rejection occurs when the constrained
job arrives last and asks for the top-ranked day. Under RANDOM-VERTEX,
the flexible jobs must additionally have all chosen that same day.

\section{Exact computation and the remaining obstruction}
For an arbitrary later state let $c_i$ be day $i$'s remaining capacity
and let $a=|\{i:c_i>0\}|$. An unprocessed degree-$d$ job fails with
probability $\binom{m-a}{d}/\binom md$. Under RANDOM-VERTEX each available
day receives probability $(1-\binom{m-a}{d}/\binom md)/a$ when $a>0$.
Under RANKING, the $r$th available day in priority order receives
$\binom{m-r}{d-1}/\binom md$. Full days remain in the sampling universe.
These transitions and the remaining degree multiplicities give a finite
exact recursion for each output distribution.

The first-saturation theorem holds for arbitrary numbers of days and jobs.
The continuation argument proves all tails through $2C$ for arbitrary $n$,
and therefore the full comparison for $n\leq2C$ or for $m=2$.
The marked-hitting arguments below extend this to threshold $2C+2$.
The three- and four-day results below also close every threshold for $m\leq4$.
The unresolved regime in this work therefore has $m\geq5$, $C\geq3$,
and thresholds $k\geq2C+3$ (necessarily also $n\geq2C+3$).
After further saturation, the remaining capacities depend on the common
ranking and the history. The arguments here do not yet control all of
this information.

The next research decision is whether a distributional comparison of
these residual states can survive further saturation. Passing larger
finite tests alone is not a substitute for such a comparison.

\subsection{An exact reduction for threshold $2C+1$}
There is a smaller description for the first threshold beyond
Theorem~\ref{thm:2C}.
Fix the positive degree sequence and let $\sigma$ be the arrival index
at which the $2C$th successful assignment is made, with $\sigma=\infty$
if this never happens. On $\{\sigma\leq n\}$, let $A$ be the number of
full days at that instant. Necessarily $A\in\{0,1,2\}$, since only $2C$
jobs have been accepted. For $a\in\{0,1,2\}$ define
\[
 Q_a(s)=\prod_{j=s+1}^{n}
          \frac{\binom a{d_j}}{\binom m{d_j}},
 \qquad Q_a(n)=1.
\]
With an expectation restricted to the event $\sigma\leq n$, the exact
identity for either algorithm is
\begin{equation}\label{eq:next-threshold}
 \Pr(M\geq2C+1)
 =\sum_{s=1}^{n}\sum_{a=0}^{2}
   \Pr(\sigma=s,A=a)\,[1-Q_a(s)].
\end{equation}
Indeed, after $2C$ successes, we need only one more success. If all
subsequent jobs fail, no capacities change. Each fails precisely when its
entire neighborhood lies in the same set of $a$ full days. Future
neighborhoods are independent of the past and of each other, giving
$Q_a(s)$. The identity does not depend on how the remaining free slots
are distributed among nonfull days, nor on their ranking.

For a nonempty suffix, $Q_0(s)=0$. The value $Q_1(s)$ is zero if any
future degree exceeds one, and otherwise is $m^{-(n-s)}$.
The value $Q_2(s)$ is zero if any future degree exceeds two, and otherwise,
with $h_1,h_2$ counting the future degrees one and two, equals
\[
 Q_2(s)=\left(\frac2m\right)^{h_1}
         \binom m2^{-h_2}.
\]
Thus the next comparison concerns a weighted joint distribution of just
$(\sigma,A)$, not an arbitrary continuation function of the full capacity
vector. The theorem already shows that $\sigma_\RV$ is stochastically
\emph{earlier} than $\sigma_\RK$, by applying the threshold-$2C$
comparison to every prefix. But an earlier $\sigma$ can come with a larger
$A$, as the preceding second-saturation example demonstrates.

The following argument supplies the required joint comparison.

\section{A marked hitting-time coupling and threshold $2C+1$}
Continue to use $T$ for first saturation and $(\sigma,A)$ for the first
arrival reaching $2C$ accepted jobs and the number of full days then.
All statements in this section concern a fixed positive-degree sequence.
The empty conditional events in the next lemma impose no requirement.

\begin{lemma}[Concentration outside the first full day]\label{lem:mark}
Let $m\geq2$ and $C\geq1$, and put $\lambda=(m-1)^{1-C}$.
For each $t<2C$ and finite $s$ such that the relevant conditioning event
has positive probability,
\[
 \Pr(A_\RV=2\mid T_\RV=t,\sigma_\RV=s)=\lambda,
 \qquad
 \Pr(A_\RK=2\mid T_\RK=t,\sigma_\RK=s)\geq\lambda.
\]
On these events $A\in\{1,2\}$.
\end{lemma}
\begin{proof}
Condition more finely on the first full day $i$, on which of the first
$t$ jobs were assigned to $i$, and on which jobs between arrivals $t+1$
and $s$ were rejected. The first set has exactly $C$ members, includes
arrival $t$, and each of its other members is earlier than $t$.
Before first saturation every job succeeds. Thus exactly $t-C<C$
assignments went to days other than $i$ before or at $t$; these cannot
have filled any other day.

Between $t$ and $s$, before the last success, fewer than $2C$ jobs have
been accepted. Therefore $i$ remains the only full day: filling a second
one would already require $2C$ successes. Rejections in this interval are
exactly degree-one jobs choosing $i$. There are $2C-t$ successful
assignments in this interval, including the arrival at $s$. In total,
exactly
\[
 (t-C)+(2C-t)=C
\]
successful assignments go to days other than $i$ by time $s$.
The event $A=2$ is exactly the event that all these $C$ assignments
choose the same remaining day.

For RANDOM-VERTEX, conditional on the finer information, their destinations
are independent and uniform among the $m-1$ days other than $i$.
Before $t$, this follows by conditioning independent uniform choices
to avoid $i$. After $t$, each successful job chooses uniformly among
those $m-1$ days; its success or failure depends only on whether a
degree-one job chose $i$. No other day fills before the last of these
assignments, so there is no additional state-dependent restriction.
Their coincidence probability is exactly $(m-1)^{1-C}$.

For RANKING, fix the original priority order and index the remaining days
in their inherited order. Before $t$, an assignment conditioned to avoid
$i$ has probabilities proportional to $p_d(r)$ with the entry for $i$
removed. This is still a nonincreasing probability vector. After $t$,
conditional on success, a degree-$d$ job's probability of choosing the
$r$th remaining day is proportional to
$\binom{m-r}{d-1}/\binom md$, again nonincreasing in the inherited order.
For a degree-one job this normalization gives the uniform vector.

The $C$ destinations are independent under the finer conditioning:
the prescribed assignments to $i$ and prescribed rejections constrain
individual jobs only. The bounds on the number of accepted jobs ensure
that no further saturation restrictions have been suppressed. Write
$q_\ell(r)$ for the $C$ nonincreasing destination probability vectors.
Lemma~\ref{lem:products}, on $m-1$ positions, gives
\[
 \Pr(A=2\mid\hbox{finer information})
 =\sum_{r=1}^{m-1}\prod_{\ell=1}^{C}q_\ell(r)
 \geq(m-1)^{1-C}.
\]
Averaging over the finer information proves both claims of the lemma.
\end{proof}

For example, with three days and capacity three, the constant in
Lemma~\ref{lem:mark} is $1/4$. Conditional on first saturation occurring
before the sixth accepted job and on when the sixth success occurs,
RANDOM-VERTEX has exactly a $1/4$ probability of having two full days
then. RANKING's conditional probability is at least $1/4$.
This conditional comparison is compatible with the earlier example in
which RANDOM-VERTEX has a larger \emph{unconditional} probability of
two full days after six arrivals: it reaches six successes more often.

\begin{theorem}[Joint comparison at $2C$ successes]\label{thm:joint}
For every fixed positive-degree sequence, there is a coupling of
$(\sigma_\RV,A_\RV)$ and $(\sigma_\RK,A_\RK)$ such that
\[
 \sigma_\RV\leq\sigma_\RK,
 \qquad A_\RV\leq A_\RK\quad\hbox{whenever }\sigma_\RK<\infty.
\]
An infinite hitting time is treated as failure to reach the level, and
its mark need not be defined.
\end{theorem}
\begin{proof}
If $n<2C$, both hitting times are infinite. Otherwise couple the first
saturation times so that $T_\RV\geq T_\RK$, using
Theorem~\ref{thm:first}. Independently generate the same variables $X_j$
as in \eqref{eq:2C-mixture}: a degree-one job has success parameter
$(m-1)/m$, and other degrees give $X_j=1$.

Given first saturation at $t<2C$, the hitting time of $2C$ successes is
the first $u\geq t$ for which $t+\sum_{j=t+1}^{u}X_j\geq2C$,
or infinity if there is no such arrival. For $t\geq2C$, the hitting
time is $2C$. These conditional laws are exact because no second day can
fill before $2C$ successes. For every prefix $u$, delaying $t$ replaces
some $X_j\leq1$ by certain successes. Hence the common-$X$ construction
gives $\sigma_\RV\leq\sigma_\RK$.

We now add the marks without changing these hitting-time marginals.
If both first saturation times are below $2C$ and both hitting times
are finite, Lemma~\ref{lem:mark} says that the conditional RANDOM-VERTEX
mark is $1+\operatorname{Bernoulli}(\lambda)$, whereas the RANKING mark
is $1+\operatorname{Bernoulli}(p)$ for some $p\geq\lambda$, conditional
on its own $(T,\sigma)$. A fresh shared uniform variable couples these
marks in the required order. We may use these conditional mark laws
even though the artificial $X$ variables have been shared: sampling the
mark according to its own $(T,\sigma)$ recovers each algorithm's correct
joint distribution.

If $T_\RV\geq2C>T_\RK$, then $A_\RV\leq1\leq A_\RK$ whenever
both hit the level. If both first saturation times are at least $2C$,
both hitting times equal $2C$ and
$A=\mathbf1_{\{T=2C\}}$. The order $T_\RV\geq T_\RK$ again gives
$A_\RV\leq A_\RK$. Finally, if RANKING never reaches $2C$, its
continuation payoff is zero and no mark comparison is needed; if
RANDOM-VERTEX never reaches it, neither does RANKING in this coupling.
\end{proof}

\begin{theorem}[All tails through $2C+1$]\label{thm:2Cplus1}
For every $m\geq1$, common capacity $C\geq1$, any number $n$ of jobs,
and every fixed positive-degree sequence independent of neighborhoods,
\[
 \Pr(M_\RV\geq k)\geq\Pr(M_\RK\geq k)
 \qquad\text{for every integer }k\leq2C+1.
\]
Consequently the full comparison holds for $n\leq2C+1$, and all these
statements hold after averaging over a uniformly random job order.
\end{theorem}
\begin{proof}
The case $m=1$ is equality. The tails through $2C$ are already proved,
so take $m\geq2$ and $k=2C+1$.
For finite $s$, the payoff $1-Q_a(s)$ in \eqref{eq:next-threshold}
is nonincreasing in both $s$ and $a$: an earlier $s$ adds factors in
$[0,1]$ to $Q$, while decreasing $a$ decreases each factor.
Define the payoff at $\sigma=\infty$ to be zero.
Theorem~\ref{thm:joint} therefore couples the two hitting states so that
the RANDOM-VERTEX payoff is at least the RANKING payoff pointwise.
Taking expectations in \eqref{eq:next-threshold} proves the remaining
tail. Thresholds above $n$ have probability zero, giving the stated full
comparison when $n\leq2C+1$. Averaging over degree orders and deleting
degree-zero jobs recovers the original source model.
\end{proof}

The next section controls the additional saturation that was missing
from this argument. The mark $A$ alone does not determine its probability;
a conditional occupancy inequality supplies the needed comparison.

\section{One more marked level and threshold $2C+2$}
\begin{lemma}[A near-full bin under aligned choices]\label{lem:nearfull}
Let $b,c\geq2$ be integers. Let $Z_1,\ldots,Z_c,Y$ be independent
categorical variables on $\{1,\ldots,b\}$, each having a nonincreasing
probability vector in this same order. Put
$N_r=\#\{j:Z_j=r\}$ and $E=\{\max_r N_r<c\}$.
If $\Pr(E)>0$, then
\[
 \Pr(N_Y=c-1\mid E)\ \geq
 \beta(b,c):=\frac{c(b-1)}{b(b^{c-1}-1)}.
\]
Equality holds when all $c+1$ vectors are uniform.
\end{lemma}
\begin{proof}
Write $q_j$ for the row of $Z_j$ and $q$ for the row of $Y$.
The denominator and numerator of the conditional probability are
\begin{align*}
 D&=1-\sum_{r=1}^b\prod_{j=1}^c q_j(r),\\
 H&=\sum_{r=1}^b q(r)\sum_{h=1}^c
             (1-q_h(r))\prod_{j\ne h}q_j(r).
\end{align*}
Indeed $E$ excludes exactly the outcomes in which all $c$ old draws
coincide, and $N_r=c-1$ requires exactly one draw to avoid $r$.
It suffices to prove $H-\beta(b,c)D\geq0$; if $D=0$, also $H=0$.

This expression is affine in each probability row separately. Every
nonincreasing probability row $p$ is a convex combination of the rows
$U_r$, uniform on the first $r$ positions and zero elsewhere:
\[
 p=\sum_{r=1}^b r(p(r)-p(r+1))U_r,\qquad p(b+1)=0.
\]
The coefficients are nonnegative and sum to one. Expanding in each
row therefore reduces the inequality to $q_j=U_{r_j}$ and $q=U_\ell$.
Permute the old rows so that $r_1\leq\cdots\leq r_c$, and put
$a=r_1$, $v=r_2$, $P=\prod_jr_j$, $S=\sum_jr_j$.
Then $D=(P-a)/P$ and
\[
 w_x:=\Pr(N_x=c-1)=
 \begin{cases}
 (S-c)/P,&x\leq a,\\
 a/P,&a<x\leq v,\\
 0,&x>v.
 \end{cases}
\]
For the middle range, only the unique shortest row can be the draw
outside $x$; that range is empty if the shortest length is repeated.
The sequence $w_x$ is nonincreasing. In fact, if $a<v$, then
$S-c\geq a$: this follows from $v\geq2$ and $c\geq2$.
Consequently averaging $w_x$ on any initial segment gives at least
its average on all $b$ positions. Taking $Y$ uniform on all $b$
positions can only decrease $H$.

Suppose $D>0$. With this last uniform choice we obtain
\[
 \frac HD\ \geq
 \frac{S-c+v-a}{b(P/a-1)}
 =\frac{2(v-1)+\sum_{j=3}^c(r_j-1)}
        {b\bigl(v\prod_{j=3}^c r_j-1\bigr)}.
\]
For $c=2$ the last expression is $2/b=\beta(b,2)$.
For $c\geq3$, increase the lengths $r_3,\ldots,r_c$ one at a time
to $b$. Holding the others fixed, the ratio without the factor $1/b$
has the form $(A+x-1)/(Bx-1)$, where $A$ is a nonnegative integer,
$B\geq1$, and $A=0$ implies $B=1$.
Its derivative has the sign of $B(1-A)-1\leq0$.
The denominator stays positive, since the starting $D$ is positive.
After these increases the ratio without $1/b$ is
\[
 R(v)=\frac{2(v-1)+(c-2)(b-1)}{v b^{c-2}-1}.
\]
Its derivative has the sign of
$b^{c-2}(2-(c-2)(b-1))-2$, which is nonpositive:
for $c=3$ it is $-(b-1)(b-2)$, and for $c\geq4$ the expression
in parentheses is nonpositive. Thus increasing $v$ to $b$ also
cannot increase the ratio. The resulting value, after division by $b$,
is $\beta(b,c)$. This proves the inequality at all extreme rows and
hence at all rows. Substitution of uniform rows gives equality.
\end{proof}

Continue to write $(\sigma,A)$ for the marked first hitting time of
$2C$ successes. Let $\rho$ be the arrival of the $(2C+1)$st success,
and $B$ the number of full days then, when $\rho$ is finite.

\begin{lemma}[Conditional risk of filling the second day]\label{lem:fillsecond}
Let $m\geq3$, $C\geq2$, and $\beta=\beta(m-1,C)$.
For every positive-probability conditioning event
$\{T=t,\sigma=s,A=1,\rho=u\}$, with $u$ finite,
\[
 \Pr(B_\RV=2\mid T=t,\sigma=s,A=1,\rho=u)=\beta,
 \qquad
 \Pr(B_\RK=2\mid T=t,\sigma=s,A=1,\rho=u)\geq\beta,
\]
where each conditional law refers to its own algorithm.
\end{lemma}
\begin{proof}
Necessarily $t\leq2C$. Condition further on the identity $i$ of the
first full day, the jobs assigned to $i$ up to $t$, and the rejections
between $t$ and $s$. As in Lemma~\ref{lem:mark}, exactly $C$ accepted
assignments outside $i$ have occurred by $s$. Their conditional law can
be represented by independent categorical draws restricted to the event
$E$ that those $C$ destinations do not all coincide. Before that
restriction, the rows are uniform among the other $b=m-1$ days for
RANDOM-VERTEX, and nonincreasing in the inherited priority order for
RANKING. For $t<2C$, the condition $A=1$ imposes exactly this restriction.

For clarity this also covers the boundary $t=2C$, not covered by
Lemma~\ref{lem:mark}. In that case $s=t$, all jobs up to $t$ succeed,
and $C$ of their destinations avoid $i$. The requirement that no other
day filled earlier is precisely $E$: these $C$ assignments fill another
day if and only if they all coincide. There are no further restrictions
on their destinations. If $t<2C$, earlier saturation elsewhere is
already impossible by the accepted-job count, as in that lemma.

Until $u$, every intervening rejected job must be a singleton choosing
$i$. These rejections do not change loads; their probabilities depend
on $i$ but not on the destinations outside $i$. Conditional on success
at $u$, the new destination $Y$ is independent of the old destinations,
uniform for RANDOM-VERTEX and nonincreasing for RANKING. In the latter
case its rank probabilities are proportional to
$\binom{m-r}{d_u-1}/\binom m{d_u}$ for $1\leq r\leq b$.
This is true also for a singleton after normalization.
The extra success fills the second day exactly when $N_Y=C-1$.
Lemma~\ref{lem:nearfull} proves the RANKING bound and its equality
case gives the RANDOM-VERTEX value. Averaging over the finer
conditioning preserves both conclusions.
\end{proof}

\begin{theorem}[Joint comparison at $2C+1$ successes]\label{thm:jointplus1}
For $m\geq3$, $C\geq2$ and any fixed positive-degree sequence, there
is a coupling of $(\rho_\RV,B_\RV)$ and $(\rho_\RK,B_\RK)$ such that
\[
 \rho_\RV\leq\rho_\RK,\qquad
 B_\RV\leq B_\RK\quad\hbox{whenever }\rho_\RK<\infty.
\]
\end{theorem}
\begin{proof}
Use the construction of Theorem~\ref{thm:joint}, retaining the coupled
first saturation times $T_\RV\geq T_\RK$ as well as $(\sigma,A)$.
When $A\geq1$, $T$ is already in the past. Until the next success,
failure at arrival $j>\sigma$ has probability
\[
 f_{d_j}(A)=\frac{\binom{A}{d_j}}{\binom{m}{d_j}},
\]
independently across these arrivals. This follows because only the
set of full days can cause a failure, and its size remains $A$.
When $A=0$, the next arrival succeeds deterministically, even if the
retained value of $T$ gives additional future information.

Use fresh shared independent uniforms indexed by arrival to generate
these first subsequent successes. Since $\sigma_\RV\leq\sigma_\RK$
and $A_\RV\leq A_\RK$, this gives $\rho_\RV\leq\rho_\RK$:
if RANDOM-VERTEX has not succeeded by $\sigma_\RK$, its success
probability at each common future arrival is at least RANKING's.
If RANKING's hitting time is infinite, no mark order is required.

Suppose both $\rho$ are finite. Add the marks according to their own
conditional laws given $(T,\sigma,A,\rho)$. If $A_\RV<A_\RK$, one
success can increase the number of full days by at most one, so
$B_\RV\leq A_\RV+1\leq A_\RK\leq B_\RK$ for any mark coupling.
If both $A$ equal two, the two full days contain all $2C$ previous
successes. Every other day is empty. Since $C\geq2$, one success
cannot fill a third day, so both $B$ equal two.

If both $A$ equal zero, both $\sigma$ equal $2C$, and both $\rho$
equal $2C+1$. Here $B=\mathbf1_{\{T=2C+1\}}$; the retained order
of $T$ gives the required order of $B$. Finally if both $A$ equal one,
Lemma~\ref{lem:fillsecond} gives conditional marks
$1+\operatorname{Bernoulli}(\beta)$ for RANDOM-VERTEX and
$1+\operatorname{Bernoulli}(p)$ for RANKING, with $p\geq\beta$.
A fresh shared uniform gives the ordered marks. Sampling these marks
conditional on each algorithm's own $(T,\sigma,A,\rho)$ recovers
its correct marginal joint law; the artificial shared uniforms do not
impose any additional conditioning on the original process.
\end{proof}

\begin{theorem}[All tails through $2C+2$]\label{thm:2Cplus2}
For every $m\geq1$, common capacity $C\geq1$, any number $n$ of jobs,
and every fixed positive-degree sequence independent of neighborhoods,
\[
 \Pr(M_\RV\geq k)\geq\Pr(M_\RK\geq k)
 \qquad\text{for every integer }k\leq2C+2.
\]
Consequently the full comparison holds for $n\leq2C+2$, as well as
for $m=2$ with arbitrary $n$. All conclusions hold after averaging
over uniformly random job order.
\end{theorem}
\begin{proof}
For $C=1$, after $z$ successes there are exactly $z$ full days.
The next success probability is $1-\binom z{d_j}/\binom m{d_j}$
under either algorithm, so their matching counts have identical laws.
The cases $m\leq2$ have already been proved. Otherwise use
Theorem~\ref{thm:jointplus1} and the one-additional-success payoff
\[
 1-\prod_{j>\rho}\frac{\binom B{d_j}}{\binom m{d_j}}.
\]
As before, this payoff is nonincreasing in the finite hitting index
$\rho$ and the mark $B$, and is zero if $\rho=\infty$.
Its expectation is $\Pr(M\geq2C+2)$, so the coupled order proves
that tail. The earlier theorem covers all smaller thresholds.
Deleting degree-zero jobs and averaging fixed orders gives the
source-model conclusions.
\end{proof}

For $m=C=3$, the new conditional risk is $\beta(2,3)=1/2$:
given exactly one full day at the sixth success, the seventh success
fills a second day with probability $1/2$ for RANDOM-VERTEX and
at least $1/2$ for RANKING, under the conditioning stated above.
The new theorem controls the probability of accepting at least eight
jobs, regardless of how many jobs arrive.

This proves an additional tail of the target comparison for unbounded
parameters. The next section goes further for three days by handling
any feasible number of assignments outside the first full day.

\section{A conditional hazard comparison for two remaining days}
The preceding near-full inequality used exactly $C$ assignments outside
the first full day. With two remaining days, the comparison extends to
every feasible number of outside assignments. This will close the whole
three-day problem rather than a single additional threshold.

\begin{lemma}[Two-bin conditional saturation hazard]\label{lem:twohazard}
Let $C\geq2$, $0\leq r\leq2C-2$, and let $Z_1,\ldots,Z_r,Y$ be
independent choices of one of two bins. Suppose every choice gives the
first bin probability at least $1/2$. Write $X$ for the number of old
choices of the first bin, and condition on both old bin loads being
strictly below $C$. If this event has positive probability, the chance
that the new choice fills a bin is at least
\[
 h_C(r)=
 \begin{cases}
 0,&r<C-1,\\[2pt]
 \displaystyle\frac{\binom r{C-1}}
 {\sum_{k=r-C+1}^{C-1}\binom rk},&C-1\leq r\leq2C-2.
 \end{cases}
\]
Equality holds when all choices are fair.
\end{lemma}
\begin{proof}
The claim is immediate for $r<C-1$. Otherwise put $L=r-C+1$ and
$U=C-1$, so that the conditioning event is $E=\{L\leq X\leq U\}$
and $L+U=r$. If $L=U$, both bins already have load $C-1$, and the
next choice fills one with probability one, equal to $h_C(r)$.
We may therefore suppose $L<U$.

Any choice with first-bin probability $p\geq1/2$ is a mixture of a
deterministic first-bin choice (weight $2p-1$) and a fair choice
(weight $2(1-p)$). Independently expose these mixture components for
the $r$ old choices. It suffices to prove the same conditional bound in
each component with $\Pr(E)>0$: conditioning on $E$ simply reweights
the mixture components. In a component with $h$ deterministic choices,
$X$ has law $h+\operatorname{Binomial}(r-h,1/2)$.

Let $\mu$ be the fair $\operatorname{Binomial}(r,1/2)$ law conditioned
on $[L,U]$. It is symmetric under $k\mapsto r-k$. Using
\[
 \binom{r-h}{k-h}=\frac{\binom rk\binom kh}{\binom rh},
\]
the component's conditional law is $\mu$ tilted by the nonnegative
weight $g_h(k)=\binom kh$, with the convention that it is zero for
$k<h$. Its normalizing constant $\mathbb E_\mu g_h$ is positive.

The integer sequence $g_h(k)$ is convex on $0,\ldots,r$:
for $h\geq2$, its second forward difference is $\binom k{h-2}\geq0$;
for $h=0,1$ it is constant or linear. Thus
$f_h(k)=g_h(k)+g_h(r-k)$ is convex and symmetric. On $[L,U]$ its
maximum is attained at the endpoints, whose values agree. Symmetry of
$\mu$ gives
\[
 \mathbb E_\mu g_h
 =\tfrac12\mathbb E_\mu f_h
 \leq\tfrac12\bigl(g_h(L)+g_h(U)\bigr).
\]
Since $\mu(L)=\mu(U)=h_C(r)$, it follows that in this component
\[
 \Pr(X=L\mid E)+\Pr(X=U\mid E)\geq2h_C(r).
\]
Also $g_h(U)\geq g_h(L)$, so the upper endpoint has probability at
least that of the lower endpoint. The next independent choice has
first-bin probability $p\geq1/2$. Its conditional saturation probability
is therefore
\[
 p\Pr(X=U\mid E)+(1-p)\Pr(X=L\mid E)
 \geq\tfrac12\bigl(\Pr(X=U\mid E)+\Pr(X=L\mid E)\bigr)
 \geq h_C(r).
\]
The bound survives mixing over $h$ and the exposed components.
For fair old choices, the endpoints have equal mass $h_C(r)$;
a fair next choice gives the claimed equality.
\end{proof}

\begin{lemma}[One full day, with the scalar history retained]\label{lem:historyhazard}
Let $m=3$, $C\geq2$, and $C\leq K\leq3C-2$. At the first time
$\sigma_K$ of $K$ accepted jobs, suppose exactly one day is full.
Condition on the first saturation time $T$ and on any retained history
of accepted-job hitting times and their numbers of full days up to
$\sigma_K$. Further condition on the arrival index of the next success,
when it is finite. The probability that this success fills a second day
is exactly $h_C(K-C)$ for RANDOM-VERTEX and at least $h_C(K-C)$ for
RANKING, on every conditioning event of positive probability.
\end{lemma}
\begin{proof}
All the conditioned information is in the past of $\sigma_K$, except
the waiting time to its next success; the first saturation has already
occurred. Refine it by specifying the unique full day $i$, which jobs
were assigned to $i$, and the acceptance or rejection of every arrival
through $\sigma_K$. Exactly $C$ accepted jobs went to $i$; their last
arrival is $T$. The remaining $r=K-C$ accepted jobs chose one of the
other two days.

Fix the priority order for RANKING. Before $T$, each of these remaining
destinations has the original rank-choice law conditioned to avoid $i$.
After $T$, its law is the choice law with $i$ unavailable, conditional
on success. Both laws give the earlier of the two surviving ranks
probability at least $1/2$. For RANDOM-VERTEX both laws are fair.
The underlying choices for different jobs are independent.

The complete refinement prescribes all choices of $i$ and all rejected
jobs separately. It therefore fixes the availability path of $i$.
The only remaining restriction on the other destinations is that neither
of their final loads reaches $C$. Since their loads are nondecreasing,
this is equivalent to neither having filled at any earlier point.
It also imposes all the retained full-day marks: they are zero before
$T$ and one from $T$ through $\sigma_K$. The refinement fixes the
success indices, so it imposes all the retained hitting times as well.
Thus the two residual loads have precisely the conditional independent
choice representation in Lemma~\ref{lem:twohazard}.

Between $\sigma_K$ and the next success, a rejection is a singleton
whose feasible day is $i$. These events leave the residual loads
unchanged and do not bias their distribution. Conditional on success,
the next destination is independent of the old ones and is fair for
RANDOM-VERTEX or favors the earlier surviving rank for RANKING.
Applying Lemma~\ref{lem:twohazard} and averaging over the refinement
proves the claim. The equality value depends only on $K,C$, so the
averaging preserves it even when the retained scalar histories differ.
\end{proof}

\begin{theorem}[Three days, arbitrary capacity and job count]\label{thm:three}
For $m=3$, every common integer capacity $C\geq1$, and every finite
fixed degree sequence independent of neighborhoods,
\[
 \Pr(M_\RV\geq k)\geq\Pr(M_\RK\geq k)
 \qquad\text{for every integer }k.
\]
In particular, the full random-order comparison holds for three days,
without a bound on the capacity or on the number of jobs.
\end{theorem}
\begin{proof}
Delete degree-zero jobs. The case $C=1$ was handled in
Theorem~\ref{thm:2Cplus2}, so take $C\geq2$.
Write $(\sigma_K,A_K)$ for the first arrival reaching $K$ successes
and its number of full days, with $\sigma_K=\infty$ if the level is
not reached. Put $(\sigma_0,A_0)=(0,0)$.
Couple the actual first saturation times using Theorem~\ref{thm:first}
so that $T_\RV\geq T_\RK$. We inductively extend this coupling to the
scalar hitting histories through every level $K\leq3C-1$, preserving
\[
 \sigma_{K,\RV}\leq\sigma_{K,\RK},\qquad
 A_{K,\RV}\leq A_{K,\RK}
 \quad\hbox{whenever }\sigma_{K,\RK}<\infty.
\]
The base level zero is deterministic. Histories retain $T$ throughout.

Suppose the order has been established at $K\leq3C-2$. If neither
algorithm has hit, there is nothing to extend. If only RANDOM-VERTEX
has hit, sample its continuation with its correct conditional law;
RANKING will not hit any higher level. Otherwise both have hit.
Until the next success, a state with $a\geq1$ full days has failure
probability $\binom a{d_j}/\binom3{d_j}$ at each later arrival $j$.
This depends only on $a$ and leaves the state unchanged on failure.
If $a=0$, the next arrival succeeds with certainty, including when
the retained $T$ gives future information. Fresh shared uniforms indexed
by arrival thus couple the next-success indices in the same order:
RANDOM-VERTEX starts no later and has no larger failure probability
at any common future arrival before its next success.

If RANKING has no next success, no new mark comparison is needed.
Otherwise both do, and marks can be coupled according to their own
conditional scalar histories and next-success indices as follows.
If the old marks differ, one success increases a mark by at most one,
so the new marks are automatically ordered. If both old marks are two,
both new marks remain two, since $K+1<3C$ rules out three full days.
If both old marks are zero, both hitting times equal $K$ and the next
ones equal $K+1$. Their new marks are $\mathbf1_{\{T=K+1\}}$;
the retained order of $T$ orders these marks as well.
If both old marks are one, Lemma~\ref{lem:historyhazard} gives new
marks $1+\operatorname{Bernoulli}(h_C(K-C))$ for RANDOM-VERTEX and
$1+\operatorname{Bernoulli}(p)$ for RANKING, where
$p\geq h_C(K-C)$. A fresh shared uniform orders the new marks.
At $K<C$, the case of one full day cannot occur.

At each step we sample from the conditional law given each algorithm's
own retained history. Hence both marginal histories are correct;
the shared auxiliary uniforms introduce no extra restrictions on them.
This proves the induction through $3C-1$. The hitting-time order gives
every matching-size tail up to that level. At level $3C-1$, use the
decreasing one-further-success payoff
\[
 1-\prod_{j>\sigma_{3C-1}}
          \frac{\binom{A_{3C-1}}{d_j}}{\binom3{d_j}},
\]
with value zero at an infinite hitting time, to obtain the tail at
$3C$. Larger tails are zero because the total capacity is $3C$.
Finally average over random job orders to recover the source model.
\end{proof}

This theorem uses the first-saturation order and the two-bin conditional
hazard lemma. It does not require extrapolating the finite matching
tests to larger capacities or longer sequences. The case $m=C=3$ now
includes the previously missing ninth-success threshold, even when
arbitrarily many jobs arrive.

\section{A precise sufficient inequality for the full conjecture}
For $b\geq1$, $C\geq2$, and $0\leq r\leq b(C-1)$, let
\[
 F_{b,C}(r)=r!\,[z^r]\left(\sum_{j=0}^{C-1}\frac{z^j}{j!}\right)^b,
 \qquad
 h_{b,C}(r)=1-\frac{F_{b,C}(r+1)}{b F_{b,C}(r)}.
\]
Here $F_{b,C}(r)$ counts assignments of $r$ labeled choices to $b$
bins with every load below $C$, and is positive in the stated range;
$F_{b,C}(b(C-1)+1)=0$. Thus $h_{b,C}(r)$ is the conditional
saturation probability for fair independent old choices and a fair
new choice. For $b=2$ it is $h_C(r)$ above.

The following assertion is denoted $\mathcal H(b,C)$; it is proved
below for $b\leq3$ at every capacity, and for $C\leq5$ at every
number of bins. The remaining target has $b\geq4$, $C\geq6$.
For every $r$ in the stated range,
let $Z_1,\ldots,Z_r,Y$ be independent categorical choices on $b$ bins
whose probability vectors are nonincreasing in the same order. With
$N_i=\#\{j:Z_j=i\}$ and $E=\{\max_iN_i<C\}$, require
\begin{equation}\label{eq:generalhazard}
 \Pr(N_Y=C-1\mid E)\ \geq\ h_{b,C}(r)
 \qquad\text{whenever }\Pr(E)>0.
\end{equation}
Lemma~\ref{lem:twohazard} proves $\mathcal H(2,C)$ for every $C$.
The case $b=1$ is deterministic. Lemma~\ref{lem:nearfull} proves
the slice $r=C$ for all $b\geq2$.

\begin{theorem}[Conditional reduction of the remaining conjecture]\label{thm:hazardreduction}
Fix $m\geq2$ and $C\geq2$. If $\mathcal H(b,C)$ holds for every
$2\leq b\leq m-1$, then the full matching-size dominance holds for
these $m,C$, for every finite fixed degree sequence and hence for
uniform random job order. In particular, proving
\eqref{eq:generalhazard} for all $b,C$ would settle the original
equal-capacity conjecture.
\end{theorem}
\begin{proof}
The essential observation is the conditional residual-load law at
a success level $K$ with $a\geq1$ full days and $b=m-a$ nonfull
days. Retain any scalar hitting history through this level and the
first saturation time. Refine this information by specifying the set
of full days, their individual filling times, every accepted job
assigned to one of those full days, and all acceptance/rejection
indicators through the current hitting time.

This refinement fixes the full-day set at each earlier arrival.
For an accepted job whose final destination is one of the $b$ still
nonfull days, its individual destination law is the choice law at
that arrival conditioned to land in this final set. Under
RANDOM-VERTEX this is uniform on these $b$ days. Under RANKING,
the available-day rank probabilities are nonincreasing; removing
the entries of other days and renormalizing preserves their order.
These rows can depend on the job and the prescribed filling history.

There are exactly $r=K-aC$ such jobs. The prescribed assignments to
eventually full days and the prescribed rejected jobs each constrain
individual jobs only. All remaining path restrictions are equivalent
to every final nonfull load being below $C$: their loads never
decrease. The prescribed filling times of the full days are already
fixed by their assigned job sets. Consequently the residual loads
are independent categorical choices with these rows conditioned on
$E=\{\max_iN_i<C\}$. This description incorporates the retained
scalar hitting history as well as the finer data.

Before the next success, rejections leave the state unchanged. Their
probabilities depend only on the current full-day set and the incoming
degrees, so conditioning on the next-success arrival index does not
further restrict the residual loads. Conditional on success, its new
destination is uniform for RANDOM-VERTEX and nonincreasing for RANKING.
Therefore, under the hypothesis of the theorem, the chance of filling
another day is exactly $h_{b,C}(r)$ for RANDOM-VERTEX and at least
$h_{b,C}(r)$ for RANKING. The constant depends only on $K,a,m,C$;
averaging over the finer data preserves this bound. For $b=1$ the same
statement is deterministic.

Now repeat the hitting-history induction from
Theorem~\ref{thm:three}, through level $mC-1$, keeping the initial
coupling $T_\RV\geq T_\RK$. The failure probability with $a$ full
days is $\binom a{d_j}/\binom m{d_j}$, so the next-success times
couple in order. Unequal old marks remain ordered after increments
of at most one. For equal positive old marks $a$, use the just-proved
conditional bound with the same $b=m-a$ and $r=K-aC$ to couple
the Bernoulli mark increments. Equal zero marks are treated using
the retained first-saturation times, exactly as before. A final
one-success payoff gives the tail at $mC$; all higher tails vanish.
\end{proof}

The unnormalized form of \eqref{eq:generalhazard},
$\Pr(E,N_Y=C-1)-h_{b,C}(r)\Pr(E)\geq0$, is affine in each of
the $r+1$ probability rows separately. The prefix-uniform decomposition
from Lemma~\ref{lem:nearfull} therefore reduces its full proof to
prefix-uniform rows, without a bound on the number of rows, their
prefix lengths, or $C$. Finite checks of such rows cannot replace
this unbounded inequality.

The reduction above explains exactly how the general conditional
inequality would reach the full target. The next section proves the
three-bin obligation and thereby closes the four-day case.

Related work on conditional dependence in urn models includes
Kahn and Neiman's
\href{https://arxiv.org/abs/1001.0610}{\emph{Conditional negative association for competing urns}}
and Chan's
\href{https://arxiv.org/abs/2509.06273}{\emph{Normalized matching property for competing urn models}}.
These are possible sources of tools for the next stage. No implication
from their results to \eqref{eq:generalhazard} has been established here,
and none of the proofs above relies on such an implication.

\section{Three residual bins: a proof using FKG}
We now prove $\mathcal H(3,C)$ without a bound on $C$ or on the
feasible number of old draws. The proof uses the classical
Fortuin--Kasteleyn--Ginibre (FKG) inequality: a log-supermodular
probability measure on a finite distributive lattice positively
correlates increasing functions. The precise result used is
Proposition~1 of
\href{https://math.bme.hu/~balint/oktatas/perkolacio/percolation_papers/fortuin_kasteleyn_ginibre.pdf}
{Fortuin, Kasteleyn and Ginibre, \emph{Correlation inequalities on some partially ordered sets}},
Communications in Mathematical Physics 22 (1971), 89--103.
We verify its hypotheses explicitly below.



\begin{lemma}[Association of capped uniform three-bin profiles]\label{lem:profilefkg}
Take $r$ independent uniform choices on three bins and condition on
all loads being at most $q\geq1$, with $0\leq r\leq3q$.
Sort the loads as $x\geq y\geq z$. On these profiles use the order
in which the maximum $x$ increases and the minimum $z$ decreases.
Any two increasing functions of the profile are positively correlated.
\end{lemma}
\begin{proof}
Use coordinates $(x,t)$, where $t=-z$ and $y=r-x+t$.
The profile set is
\[
 \mathcal L_{r,q}=\{(x,t)\in\mathbb Z^2:
 0\leq x\leq q,\ -q\leq t\leq0,\ x-2t\leq r\leq2x-t\}.
\]
It is closed under coordinatewise minimum and maximum. Indeed its
additional bounds on $t$ are $(x-r)/2\leq t\leq2x-r$, both with
nondecreasing right-hand sides. These inequalities persist under
both operations. Thus it is a finite distributive sublattice of
a product of two chains.

Up to a constant independent of the profile, its probability weight is
\[
 w(x,t)=\frac{o(x,y,z)}{x!\,y!\,z!},
 \qquad o(x,y,z)=\frac{6}{\prod_v m_v!},
\]
where $m_v$ are the multiplicities of equal entries. For comparable
pairs the log-supermodular inequality is equality. For an incomparable
pair write the two points as $(x_1,t_2)$ and $(x_2,t_1)$ with
$x_1<x_2$, $t_1<t_2$. Their meet $(x_1,t_1)$ and join $(x_2,t_2)$
both have three distinct loads. For example, at the meet,
\[
 t_1<t_2\leq2x_1-r,\qquad
 t_1\geq(x_2-r)/2>(x_1-r)/2;
\]
the join satisfies the corresponding strict inequalities as well.
Hence the orbit factors at the meet and join are both six, while
those at the original pair are at most six.

The factors $x!$ and $z!$ cancel in the product comparison. Among the
four middle loads, $r-x_1+t_2$ is the largest and $r-x_2+t_1$ the
smallest. The meet and join middle loads lie between them and have
the same combined sum. Log-convexity of factorials gives
\[
 (r-x_1+t_1)!\,(r-x_2+t_2)!
 \leq(r-x_1+t_2)!\,(r-x_2+t_1)!.
\]
Combining this with the orbit factors proves
$w(u\wedge v)w(u\vee v)\geq w(u)w(v)$.
The cited FKG inequality now gives the assertion.
\end{proof}

\begin{lemma}[Preference restrictions increase with profile imbalance]\label{lem:prefweight}
Fix $r=h+k+\ell$ old choices: $h$ are uniform on the first bin,
$k$ on the first two bins, and $\ell$ on all three bins. Relative to
the capped uniform three-bin profile law, their profile law has a
nonnegative density proportional to an increasing function $W$ in
the order of Lemma~\ref{lem:profilefkg}.
\end{lemma}
\begin{proof}
Start instead with $r$ uniform labeled choices on all three bins and
condition on the event $A$ that a prescribed group of $h$ choices
landed in bin one and a disjoint group of $k$ choices avoided bin
three. This gives exactly the stated independent prefix-uniform
rows, before imposing the cap. Given labeled loads $(n_1,n_2,n_3)$,
the conditional probability of $A$ is
\[
 \frac{\binom{n_1}h\binom{r-h-n_3}k}
      {\binom rh\binom{r-h}k}.
\]
Zero numerator terms are taken as zero. Averaging over the six
permutations of a profile (including repetitions when entries agree)
shows that its density, up to a positive constant, is
\[
 W(n_1,n_2,n_3)=
 \sum_{i\ne j}\binom{n_i}h\binom{r-h-n_j}k.
\]

Here is a coefficientwise proof of its monotonicity. Put
$A_1=1+s+t$, $A_2=1+t$, $A_3=1$ and define the symmetric polynomial
\[
 P_n(s,t)=\sum_{\pi\in S_3}\prod_{i=1}^3 A_{\pi(i)}^{n_i}.
\]
Then $W(n)=[s^h t^k]P_n(s,t)$. Balancing two exponents
$(a,b)$ to $(a-1,b+1)$ with $a\geq b+2$ cannot increase any
coefficient of $P_n$. To see this, pair permutation terms which
exchange the variables in these two positions. With these variables
denoted $X,Y$, choose their order so that $X-Y$ has nonnegative
coefficients. Their paired difference is
\[
 X^aY^b+X^bY^a-X^{a-1}Y^{b+1}-X^{b+1}Y^{a-1}
 =(XY)^b(X-Y)^2\sum_{u=0}^{a-b-2}X^{a-b-2-u}Y^u.
\]
Its coefficients are nonnegative: the possible differences $X-Y$
are $s$, $t$, or $s+t$. Multiplication by the remaining variable
power preserves this property, as does summation over pairs.

Integer balancing transfers and permutations generate the dominance
order on equal-sum triples. For three sorted entries, dominance is
equivalent to a larger maximum and a smaller minimum: the two proper
partial-sum conditions are $n_1\geq n'_1$ and
$n_1+n_2=r-n_3\geq r-n'_3$. Thus $W$ is increasing in precisely
the lattice order used above. Conditioning on the cap restricts the
profiles but does not change their density ratio, proving the lemma.
\end{proof}

\begin{theorem}[Three-bin conditional hazard]\label{thm:threehazard}
The assertion $\mathcal H(3,C)$ holds for every $C\geq2$ and every
$0\leq r\leq3(C-1)$.
\end{theorem}
\begin{proof}
Put $q=C-1$. By Lemma~\ref{lem:nextuniform}, replacing the next
choice by a uniform one cannot increase its saturation probability.
Decompose each old nonincreasing row into prefix-uniform rows, using
the decomposition in Lemma~\ref{lem:nearfull}. Independently expose
these components. Conditional on the cap event, the mixture weights
change, but a uniform lower bound for every positive-probability
component remains valid for the mixture.

In a fixed component there are $h,k,\ell$ old rows of prefix lengths
one, two, and three. Let $\mu$ be the capped uniform profile law and
$W$ the increasing density from Lemma~\ref{lem:prefweight}.
Write $B(n)=\#\{i:n_i=q\}$. This is increasing on the capped
profile lattice: a balancing transfer cannot create a new entry $q$
without having started with a larger entry, forbidden by the cap,
and it may remove an entry $q$.
Since the component has positive cap probability,
$\mathbb E_\mu W>0$. Lemma~\ref{lem:profilefkg} yields
\[
 \mathbb E(B\mid\hbox{prefix component and cap})
 =\frac{\mathbb E_\mu(BW)}{\mathbb E_\mu W}
 \geq\mathbb E_\mu B.
\]
Divide by three to obtain the uniform-next-choice saturation
probability. The right side divided by three is exactly
$h_{3,C}(r)$. This proves the bound in every mixture component,
and hence in the original rows. Uniform rows give equality.
\end{proof}

\begin{corollary}[Four days, arbitrary capacity and job count]\label{cor:four}
For $m=4$, every common integer capacity $C\geq1$, and every finite
fixed degree sequence independent of neighborhoods,
\[
 \Pr(M_\RV\geq k)\geq\Pr(M_\RK\geq k)
 \qquad\text{for every integer }k.
\]
The same full comparison holds for uniformly random job order.
\end{corollary}
\begin{proof}
For $C\geq2$, apply Theorem~\ref{thm:hazardreduction} with
Lemma~\ref{lem:twohazard} and Theorem~\ref{thm:threehazard}.
The capacity-one case has identical laws as already proved.
\end{proof}

\noindent\textbf{Stage assessment and the next obstruction.}
This closes every threshold for four days, without bounds on capacity
or job count. The remaining sufficient obligation for five days is
$\mathcal H(4,C)$. The coefficientwise balancing argument extends
to nested prefix polynomials in any number of bins, but the lattice
step above does not extend automatically.

For example, among four-entry profiles of total seven and cap four,
let $a=(4,1,1,1)$, $b=(3,3,1,0)$, and $c=(3,2,2,0)$.
In dominance order,
\[
 a\vee c=(4,2,1,0),\quad b\wedge a=(3,2,1,1),\quad b\wedge c=c.
\]
Thus $b\wedge(a\vee c)=b$ whereas
$(b\wedge a)\vee(b\wedge c)=c\ne b$: this lattice is not
distributive. This prevents a direct application of the cited FKG
theorem; it does not refute the needed correlation inequality or the
matching conjecture.
Kerner and N\'emethi's
\href{https://arxiv.org/abs/1412.8200}{\emph{A generalized FKG-inequality for compositions}}
proves an association inequality for uniform averages over compositions
(Theorem~3.1), and displays this nondistributivity obstruction.
Our next step would require the appropriate result for capped
multinomial weights, not the uniform composition measure.
No such extension is claimed here.

More precisely, full association of this measure would be stronger
than necessary. For $b$ bins let $\mu_{b,r,q}$ be the sorted uniform
multinomial profile law conditioned on the cap $q=C-1$, and let
$\alpha_j$ count the old rows of prefix length $j$, so
$\sum_j\alpha_j=r$. Define
\[
 V_i(z)=\sum_{j=i}^b z_j,\qquad
 W_\alpha(n)=[z_1^{\alpha_1}\cdots z_b^{\alpha_b}]
             \sum_{\pi\in S_b}\prod_{i=1}^b V_{\pi(i)}(z)^{n_i}.
\]
The same labeled-choice counting shows that $W_\alpha$ is proportional
to the density of the prefix-restricted profile law relative to
$\mu_{b,r,q}$. The paired-variable factorization above proves its
monotonicity for every $b$, since each difference $V_i-V_j$ for $i<j$
has nonnegative coefficients. With $B(n)=\#\{i:n_i=q\}$, the exact
remaining condition is
\begin{equation}\label{eq:prefixcovariance}
 \operatorname{Cov}_{\mu_{b,r,q}}(B,W_\alpha)\geq0
 \quad\text{for every feasible }r,\alpha.
\end{equation}
Indeed division by $b\mathbb E W_\alpha$ gives the uniform-next-choice
hazard comparison for this prefix component. Conversely that comparison
gives \eqref{eq:prefixcovariance}; zero-density components are harmless.
The old-row mixture reduction and Lemma~\ref{lem:nextuniform} show
that these covariance inequalities for all $r,\alpha$ are equivalent
to $\mathcal H(b,C)$. Thus a sufficient gate for five days is specifically
\eqref{eq:prefixcovariance} with $b=4$, rather than association for
all increasing functions on an arbitrary nondistributive lattice.
The next sections establish it for every $b$ when $C\leq5$.

\section{Capacities three and four: arbitrarily many days}
The number of bins need not be the parameter we bound. If the cap
$q=C-1$ is small, a profile can be represented by the numbers of bins
at each load. This gives a chain for $q=2$, and another distributive
lattice with log-supermodular weights for $q=3$.

\begin{lemma}[Profiles with cap two form a chain]\label{lem:captwochain}
For arbitrary $b\geq1$ and $0\leq r\leq2b$, the sorted profiles of
$r$ draws on $b$ bins with every load at most two form a chain in
dominance order. Under any probability measure on these profiles,
any two increasing functions are positively correlated.
\end{lemma}
\begin{proof}
Let $t$ be the number of entries equal to two. The profile has
$t$ twos, $r-2t$ ones, and $b-r+t$ zeros, where
\[
 \max(0,r-b)\leq t\leq\lfloor r/2\rfloor.
\]
Increasing $t$ by one replaces two ones by a two and a zero,
an unbalancing transfer. Hence the order is exactly the order of $t$.
For increasing functions $f,g$ and independent copies $T,T'$ of any
law on this chain,
\[
 2\operatorname{Cov}(f(T),g(T))
 =\mathbb E[(f(T)-f(T'))(g(T)-g(T'))]\geq0.
\]
\end{proof}

\begin{lemma}[Association of uniform profiles with cap three]\label{lem:capthreefkg}
For arbitrary $b\geq1$ and $0\leq r\leq3b$, take $r$ independent
uniform choices on $b$ bins, conditioned on all loads being at most
three. Any two functions of the sorted profile which increase in
dominance order are positively correlated.
\end{lemma}
\begin{proof}
Let $u$ count entries equal to three and $v$ entries equal to zero.
The numbers of entries equal to one and two are forced to be
\[
 m_1=2b-r+u-2v,\qquad m_2=r-b-2u+v.
\]
Consequently the profile domain in these coordinates is
\[
 \mathcal D_{b,r}=\{(u,v)\in\{0,\ldots,b\}^2:
             2u+b-r\leq v\leq(u+2b-r)/2\}.
\]
Both bounds on $v$ increase with $u$. Therefore coordinatewise
minimum and maximum preserve the domain, making it a finite
distributive sublattice of a product of two chains.

Increasing both coordinates increases the profile in dominance
order. To see this directly, for a sorted profile $n$ put
$S_n(a)=\sum_i(n_i-a)_+$. At $a=0,1,2,3$ its values are
$r$, $r-b+v$, $u$, and zero. The sum of its largest $k$ entries is
\[
 \sum_{i=1}^k n_i
 =\min_{a\in\{0,1,2,3\}}\{ka+S_n(a)\}.
\]
Thus increasing $u,v$ cannot decrease any of these partial sums.
In particular every dominance-increasing function is increasing in
our coordinate order.

Up to the fixed positive factor $b!r!$, the uniform multinomial
profile weight is
\[
 w(u,v)=\frac{1}{v!\,m_1!\,m_2!\,u!\,2^{m_2}6^u}.
\]
The multiplicity factorials count the orbit of the sorted profile;
the powers of two and six are the within-bin factorials. We now
verify log-supermodularity for the actual weight, including both.

Comparable pairs give equality. Write an incomparable pair as
$(u_1,v_2),(u_2,v_1)$ with $u_1<u_2$, $v_1<v_2$. Its meet and
join are $(u_1,v_1),(u_2,v_2)$. The factors $u!$, $v!$ cancel
in the product comparison, and so do $2^{m_2}6^u$, since their
exponents are affine in $u,v$. For each of $m_1$ and $m_2$,
the two meet/join values lie between the two original values and
have the same sum. Indeed $m_1$ has coefficients $(1,-2)$ and
$m_2$ has coefficients $(-2,1)$, so its extrema among the four
corners occur at the original incomparable pair. Log-convexity of
factorials therefore gives, for $j=1,2$,
\[
 m_j(u_1,v_1)!\,m_j(u_2,v_2)!
 \leq m_j(u_1,v_2)!\,m_j(u_2,v_1)!.
\]
It follows that
$w(x\wedge y)w(x\vee y)\geq w(x)w(y)$.
The classical FKG theorem already cited now proves the assertion.
\end{proof}

\begin{theorem}[Conditional hazard for capacities at most four]\label{thm:smallcapacityhazard}
For every $b\geq1$ and $C\in\{2,3,4\}$, the assertion
$\mathcal H(b,C)$ holds for every feasible old-draw total and
every collection of independent probability rows decreasing in
the same order.
\end{theorem}
\begin{proof}
Use the equivalent prefix-density covariance formulation
\eqref{eq:prefixcovariance}. The coefficientwise balancing argument
above proves that $W_\alpha$ increases in dominance order for
every number of bins. The function $B(n)=\#\{i:n_i=C-1\}$
also increases on the capped domain.

For $C=2$ the capped profile is unique, so $B$ is constant.
For $C=3$, apply Lemma~\ref{lem:captwochain}; here $B=t$.
For $C=4$, apply Lemma~\ref{lem:capthreefkg}; here $B=u$.
In every case $\operatorname{Cov}_\mu(B,W_\alpha)\geq0$.
Positive-density components give the conditional inequality by
division by $b\mathbb E_\mu W_\alpha$; zero-density components
do not occur after conditioning. The prefix mixture and the
uniform-next-choice lemma complete the proof for all ordered rows.
\end{proof}

\begin{corollary}[Capacities three and four, arbitrary day and job counts]\label{cor:capacityfour}
For every $m\geq1$, every common capacity $C\in\{3,4\}$,
and every finite fixed degree sequence independent of neighborhoods,
\[
 \Pr(M_\RV\geq k)\geq\Pr(M_\RK\geq k)
 \qquad\text{for every integer }k.
\]
The same holds after averaging over uniformly random job order.
Together with the cases $C=1,2$, this proves the comparison for
every common capacity at most four and arbitrarily many days.
\end{corollary}
\begin{proof}
Theorem~\ref{thm:smallcapacityhazard} supplies every residual-bin
inequality required by Theorem~\ref{thm:hazardreduction}, for any
chosen $m$. Apply that reduction and average over job order if desired.
\end{proof}

\noindent\textbf{Stage assessment.}
Capacities three and four are settled for every number of days and
jobs. The next section closes capacity five by proving an additional
coefficient inequality and using association of endpoint counts.
For higher capacities, neither that coefficient argument nor its
conditional-order step is established here. The non-distributivity
of the full profile order cannot be ignored when applying FKG.

\section{Capacity five: endpoint layers and an unbounded coefficient proof}
This section proves the capacity-five comparison for arbitrarily many
days and jobs. We first reduce it to a single coefficient inequality,
then prove that inequality by binomial identities and nonnegative
interval decompositions. Earlier finite certificates are retained
below as computational checks.

For $N\geq0$ and integer $S$, put
\[
 f(z)=z+\frac{z^2}{2}+\frac{z^3}{6},\qquad
 A(N,S)=\frac{[z^S]f(z)^N}{N!}.
\]
Coefficients outside the support are zero; in particular
$A(0,S)=\mathbf1_{\{S=0\}}$. Consider the following explicit assertion:
\begin{equation}\label{eq:fivecoefficient}
 A(N,S)A(N-2,S-4)
 \geq A(N-1,S)A(N-1,S-4)
 \qquad(N\geq2,\ S\in\mathbb Z).
\end{equation}

\begin{theorem}[A sufficient coefficient inequality for capacity five]\label{thm:fiveconditional}
If \eqref{eq:fivecoefficient} holds for every $N,S$, then any two
dominance-increasing functions of a capped uniform multinomial profile
with cap four are positively correlated, for every number of bins and
every feasible total. Consequently $\mathcal H(b,5)$ holds for all $b$,
and Arnosti's full comparison holds for common capacity five, with any
number of days and jobs and any fixed degree sequence.
More generally, if it holds for $2\leq N\leq L$, the same argument
gives association for $b\leq L$ and matching dominance for $m\leq L+1$.
\end{theorem}
\begin{proof}
Fix $b$ and $0\leq r\leq4b$. For a capped profile write $m_j$ for
the number of entries equal to $j$, and use the three coordinates
\[
 T=m_4,\qquad V=m_0,\qquad X=m_3+2m_4.
\]
The remaining multiplicities are
\begin{equation}\label{eq:fivecounts}
 m_1=2b-r+X-2V,\quad m_2=r-b-2X+V+T,\quad m_3=X-2T.
\end{equation}
Increasing all three coordinates increases the profile in dominance
order: its stop-loss sums $\sum_i(n_i-a)_+$ at $a=0,1,2,3,4$
are respectively $r,r-b+V,X,T,0$. Apply the largest-$k$-entries
identity used in Lemma~\ref{lem:capthreefkg}.

First marginalize over $X$. Up to a positive constant independent
of $t,v$, the weight of $(T,V)=(t,v)$ is
\begin{equation}\label{eq:fivemarginal}
 G(t,v)=\frac{A(b-t-v,r-4t)}{t!v!24^t}.
\end{equation}
Indeed, after prescribing the zero and four multiplicities, the
remaining $N=b-t-v$ loads belong to $\{1,2,3\}$ and sum to
$S=r-4t$. Expanding $f(z)^N/N!$ sums exactly
$1/(m_1!m_2!m_3!2^{m_2}6^{m_3})$ over those possibilities.

The positive support of $G$ is
\[
 \{(t,v)\in\{0,\ldots,b\}^2:
       b+3t-r\leq v\leq b+(t-r)/3\}.
\]
Both bounds increase with $t$, so this is a distributive sublattice.
For a unit square in its support, the factorial and exponential
factors in \eqref{eq:fivemarginal} cancel from the log-supermodular
inequality. What remains is precisely \eqref{eq:fivecoefficient}
with $N=b-t-v$, $S=r-4t$. All rectangles between an incomparable
pair lie in the support: their lowest $v$ is at least the lower
bound at their highest $t$, and their highest $v$ is at most the
upper bound at their lowest $t$. Multiplying the unit-square
inequalities therefore proves log-supermodularity for every pair.
Classical FKG gives association of $(T,V)$.

Next fix $(t,v)$. The conditional support of $X$ is the integer interval
\[
 \max(2t,r-2b+2v)\leq X\leq
                  \left\lfloor\frac{r-b+v+t}{2}\right\rfloor.
\]
Both endpoints are nondecreasing in $t,v$. These conditional laws
are stochastically increasing in both coordinates, jointly. Here
are details, including when a monotone unit-step path between two
feasible endpoint pairs does not exist.

Compare $(t,v)$ and $(t+a,v+d)$ for integers $a,d\geq0$.
On the common $X$ support, the ratio of their unnormalized
conditional weights, up to a factor independent of $X$, is
\[
 \frac{m_1!}{(m_1-2d)!}\,
 \frac{m_2!}{(m_2+a+d)!}\,
 \frac{m_3!}{(m_3-2a)!},
\]
where the $m_j$ are the old counts from \eqref{eq:fivecounts}.
The first and third factors increase with $X$, since $m_1,m_3$
increase. The middle factor also increases, since $m_2$ decreases.
Together with the ordered interval endpoints, this gives an increasing
likelihood ratio (with zero/infinite extensions at the endpoints),
and hence stochastic order. The latter implication follows directly
from the single-crossing of the two normalized probability masses.

Let $F,H$ be increasing functions of the profile. At each fixed
$(t,v)$, they are increasing functions of the single coordinate $X$,
so their conditional covariance is nonnegative. Moreover
$\overline F(t,v)=\mathbb E[F\mid T=t,V=v]$ and $\overline H(t,v)$
are increasing: couple the ordered one-dimensional conditional
$X$ laws by their quantiles, and use coordinatewise monotonicity
of $F,H$. The covariance decomposition and association of $(T,V)$
give
\[
 \operatorname{Cov}(F,H)
 =\mathbb E\operatorname{Cov}(F,H\mid T,V)
    +\operatorname{Cov}(\overline F(T,V),\overline H(T,V))\geq0.
\]
Apply this to the increasing functions $B=m_4=T$ and $W_\alpha$.
The prefix-mixture and uniform-next-choice reductions give
$\mathcal H(b,5)$, and Theorem~\ref{thm:hazardreduction} gives the
matching assertion.
For the final finite-range statement, every unit-square parameter
$N=b-t-v$ is at most $b$. The proof thus only uses coefficient
inequalities with $N\leq L$ when $b\leq L$; the matching reduction
uses at most $m-1$ residual bins.
\end{proof}

The coefficient question is readily stated using integers. If
\[
 c_{N,S}=[z^S](6z+3z^2+z^3)^N,
\]
then \eqref{eq:fivecoefficient} is equivalent to
\begin{equation}\label{eq:fiveinteger}
 (N-1)c_{N,S}c_{N-2,S-4}
 \geq N c_{N-1,S}c_{N-1,S-4}.
\end{equation}
The factor $N/(N-1)$ matters; ordinary log-concavity of one row
of coefficients does not by itself give this strengthened inequality.
The coefficients have the elementary recurrence
\[
 c_{0,S}=\mathbf1_{\{S=0\}},\qquad
 c_{N,S}=6c_{N-1,S-1}+3c_{N-1,S-2}+c_{N-1,S-3}.
\]

Related central-trinomial results do not automatically settle this
obligation. For example, Theorem~3.8 of
\href{https://arxiv.org/abs/2303.06922}
{Liang, Wang and Wang, \emph{Analytic aspects of generalized central
trinomial coefficients}} requires $\beta^2\geq2\gamma$ for
log-convexity of the central coefficients of
$(z^2+\beta z+\gamma)^N$. Here $(\beta,\gamma)=(3,6)$ fails
that condition. Indeed the first central coefficients are
$1,3,21,135$, and $3\cdot135=405<441=21^2$.
This rules out that direct application; it does not refute
\eqref{eq:fiveinteger}, which uses displaced coefficient indices.

A possible algebraic certificate is particularly explicit. Set
\[
 T_{N,K}(h)=[x^K](x^{-1}+h+x)^N,
\quad
 D_{N,K}(h)=(N-1)T_{N,K}(h)T_{N-2,K}(h)
       -N T_{N-1,K-2}(h)T_{N-1,K+2}(h).
\]
Substituting $z=\sqrt6\,x$ above shows that
\eqref{eq:fiveinteger} is exactly $D_{N,K}(\sqrt{3/2})\geq0$,
up to a positive factor, with $K=S-2N$. The polynomial $D_{N,K}$
is even in $h$. If the polynomial $Q_{N,K}(y)$ obtained by replacing
$h^2$ with $1+y$ had nonnegative coefficients for every $N,K$,
evaluation at $y=1/2$ would prove the required inequality.
The unbounded proof below establishes this stronger assertion.
For example, $Q_{3,0}(y)=2y^2+16y+11$.

\subsection{An unbounded proof of the coefficient inequality}
We now prove the missing coefficient assertion for every $N$.
The argument also proves the stronger shifted-polynomial assertion
above. No finite cutoff enters this proof.

Write
\[
 D_{n,k}(h)=\sum_{\ell\geq0}\delta_{n,k,\ell}h^{2\ell}.
\]
The following elementary binomial lemma supplies the sign of every
coefficient except possibly the constant one. Throughout, binomial
coefficients with an out-of-range lower index are zero.

\begin{lemma}[Two parity windows]\label{lem:paritywindows}
For integers $\ell\geq1$ and $R\geq0$, set
$w_\ell(u)=\binom{2\ell}{\ell+u}$. Then
\begin{align}
 &\sum_{t=-R,-R+2,\ldots,R}
       \bigl(w_\ell(t-1)-w_\ell(t+2)\bigr) \notag\\
 &\hspace{25mm}=
 \binom{2\ell-1}{\ell-R}
 -\binom{2\ell-1}{\ell-R-3}\geq0.
 \label{eq:paritywindow}
\end{align}
\end{lemma}
\begin{proof}
By symmetry, the sum of $w_\ell(t+2)$ is also the sum of
$w_\ell(t-2)$. Interleave the latter terms with $w_\ell(t-1)$.
The difference is the alternating sum of consecutive binomial
coefficients from lower index $a=\ell-R-2$ to
$b=\ell+R-1$, starting with a minus sign and ending with a plus.
Pascal's identity cancels the interior terms, leaving
\[
 \binom{2\ell-1}{b}-\binom{2\ell-1}{a-1}.
\]
Symmetry changes the first lower index to $\ell-R$ and gives
\eqref{eq:paritywindow}. The sequence
$\binom{2\ell-1}{j}$ is nondecreasing for $j\leq\ell$,
including its zero extension to negative $j$; its two middle
values at $\ell-1,\ell$ are equal. This proves the sign.
\end{proof}

\begin{lemma}[A binomial product inequality]\label{lem:binomialproduct}
For nonnegative integers $A,B$ of the same parity, put
\[
 b_A(t)=
 \begin{cases}
 \displaystyle\binom A{(A+t)/2},&t\equiv A\pmod2,\\
 0,&\text{otherwise}.
 \end{cases}
\]
For every integer $\ell\geq1$,
\begin{equation}\label{eq:Hpositive}
 H_\ell(A,B):=\sum_{u\in\mathbb Z}w_\ell(u)
 \bigl(b_A(u+1)b_B(u+1)-b_A(u-2)b_B(u+2)\bigr)\geq0.
\end{equation}
\end{lemma}
\begin{proof}
On its parity class, $b_A$ is symmetric and nonincreasing with
$|t|$. Hence it has the nonnegative interval decomposition
\[
 b_A(t)=
 \sum_{\substack{0\leq R\leq A\\R\equiv A\ (2)}}
 \bigl(b_A(R)-b_A(R+2)\bigr)
 \mathbf1_{\{|t|\leq R,\ t\equiv A\ (2)\}}.
\]
Use the analogous decomposition for $b_B$. By bilinearity it is
enough to prove \eqref{eq:Hpositive} for a pair of such interval
indicators, of radii $R,Q$ and the same parity.

Suppose first that $R\leq Q$. The positive sum then equals
$\sum_{t=-R,-R+2,\ldots,R}w_\ell(t-1)$.
In the negative sum, discard the restriction $|u+2|\leq Q$.
Since the weights are nonnegative, this bounds the negative sum
above by $\sum_{t=-R,-R+2,\ldots,R}w_\ell(t+2)$.
Lemma~\ref{lem:paritywindows} proves the required comparison.
If $Q<R$, instead discard $|u-2|\leq R$ and use symmetry of
$w_\ell$ to obtain the same bound with radius $Q$.
Summing the nonnegative interval coefficients finishes the proof.
\end{proof}

\begin{theorem}[Signs of the displaced trinomial coefficients]\label{thm:trinomialsigns}
For every $n\geq2$, $k\in\mathbb Z$ and $\ell\geq1$,
\[
 \delta_{n,k,\ell}\geq0.
\]
If $n-k$ is even, $\delta_{n,k,0}\geq0$ as well. If $n-k$ is
odd, then
\begin{equation}\label{eq:constantcompensation}
 \delta_{n,k,0}\leq0,\qquad
 \delta_{n,k,1}=-(n+1)\delta_{n,k,0}.
\end{equation}
Consequently $D_{n,k}(h)\geq0$ whenever $h^2\geq1/4$,
and every coefficient of $Q_{n,k}(y)$ is nonnegative.
\end{theorem}
\begin{proof}
The parity of every power of $h$ in $T_{n,k}$ is $n-k$;
both products defining $D$ therefore have only even powers.
If $|k|\geq n-1$, both products are zero, so it suffices to
consider $|k|\leq n-2$.

Fix $\ell$ and set $A=n-1-\ell+k$, $B=n-1-\ell-k$.
If either is negative, the coefficient of $h^{2\ell}$ is zero:
$A$ and $B$ are the total numbers of $x$ and $x^{-1}$ factors
in either product. Otherwise, factorial expansion gives the exact identity
\begin{equation}\label{eq:trinomialHidentity}
 \delta_{n,k,\ell}
 =\frac{n((n-1)!)^2}{(2\ell)!A!B!}\,H_\ell(A,B).
\end{equation}
Here are the index details. In the first product, if the first
factor supplies $c$ copies of $h$, it supplies
$(n-c+k)/2$ copies of $x$ and $(n-c-k)/2$ copies of $x^{-1}$.
After extracting the common factorial factor, its contribution is
\[
 \binom{2\ell}{c}
 \binom A{(n-c+k)/2}\binom B{(n-c-k)/2}.
\]
Set $u=\ell-c$. The last two factors are $b_A(u+1)b_B(u+1)$.
In the second product the first factor has degree $n-1$ and
index $k-2$; the corresponding factors are
$b_A(u-2)b_B(u+2)$. Both common factorial factors agree because
$(n-1)n!(n-2)!=n((n-1)!)^2$.
Nonintegral lower indices mean zero. This proves
\eqref{eq:trinomialHidentity}, and
Lemma~\ref{lem:binomialproduct} proves the claimed sign for $\ell\geq1$.

If $n-k$ is even, the second product has zero constant term and
the first has nonnegative constant term. If $n-k$ is odd, the
first has zero constant term. In this latter case with nonzero
$D$, necessarily $|k|\leq n-3$. Put
\[
 a=\frac{n-1+k}{2},\quad b=\frac{n-1-k}{2},\quad
 U=\binom{n-1}{a};\qquad a,b\geq1.
\]
Direct extraction of the constant and quadratic terms gives
\begin{align*}
 \delta_{n,k,0}&=-\frac{nab}{(a+1)(b+1)}U^2,\\
 \delta_{n,k,1}&=nabU^2\left(
     1-\frac12\frac{a-1}{a+1}-\frac12\frac{b-1}{b+1}\right)\\
 &=\frac{nab(n+1)}{(a+1)(b+1)}U^2.
\end{align*}
For example, the first product contributes $nabU^2$ to the
quadratic term; the two quadratic contributions of the second
product are the two subtracted fractions above. The formulas
also cover $a=1$ or $b=1$, with the corresponding term zero.
This proves \eqref{eq:constantcompensation}.

If the constant is negative, $n\geq3$ and
\[
 \delta_{n,k,0}+\delta_{n,k,1}h^2
 =(-\delta_{n,k,0})\bigl((n+1)h^2-1\bigr)\geq0
 \qquad(h^2\geq1/4).
\]
All remaining terms are nonnegative. Finally, on substituting
$h^2=1+y$, each nonconstant coefficient of $Q$ is a
nonnegative sum of the $\delta_{n,k,\ell}$ for $\ell\geq1$.
Its constant coefficient is $D_{n,k}(1)\geq0$.
\end{proof}

\noindent\textbf{Sharpness for the auxiliary family.}
The uniform threshold $1/4$ can even be lowered to
$\alpha=(\sqrt{42}-6)/2$. For $n\geq4$ the same argument works
at $h^2\geq1/5$, and $\alpha>1/5$. For $n=2$ every coefficient
is nonnegative. At $n=3$ the only possibly negative case is
$k=0$, where $D_{3,0}(h)=2h^4+12h^2-3$. Its nonnegative root
as a polynomial in $h^2$ is exactly $\alpha$. Thus, for a fixed
$h\geq0$, the inequalities $D_{n,k}(h)\geq0$ hold for all $n,k$
if and only if $h^2\geq\alpha$. Only the much larger value
$h^2=3/2$ is needed for matching.

\begin{corollary}[Capacity five, arbitrary day and job counts]\label{cor:capacityfive}
Inequalities \eqref{eq:fivecoefficient} and \eqref{eq:fiveinteger}
hold for every $N\geq2$ and every integer $S$.
The capped uniform multinomial profile with cap four is associated
for every number of bins and every feasible total.
For every $m\geq1$, common capacity $C=5$, every finite fixed
degree sequence independent of neighborhoods and every integer $k$,
\[
 \Pr(M_\RV\geq k)\geq\Pr(M_\RK\geq k).
\]
Uniform random job order follows by averaging. Together with the
earlier arguments this proves the comparison whenever $m\leq4$
or $C\leq5$, with arbitrarily many jobs.
\end{corollary}
\begin{proof}
The trinomial substitution has the explicit positive factor
\[
 (N-1)c_{N,S}c_{N-2,S-4}-N c_{N-1,S}c_{N-1,S-4}
 =6^{N-K-1}D_{N,K}(\sqrt{3/2}),\qquad K=S-2N.
\]
Indeed $c_{N,S}=6^{(3N-S)/2}T_{N,S-2N}(\sqrt{3/2})$.
Apply Theorem~\ref{thm:trinomialsigns}, then
Theorem~\ref{thm:fiveconditional} with no finite cutoff, followed
by the general hazard reduction. No computer enumeration is a
premise of this corollary.
\end{proof}

\subsection{Retained finite certificate}
The following certificate preceded the unbounded proof. It remains
an independent computational check of the coefficient inequality
and of the finite-range implication, but its cutoff is no longer
a restriction on the proved capacity-five result.

\begin{corollary}[Computer-assisted capacity-five range]\label{cor:fivefinite}
For $C=5$, $m\leq401$, every finite fixed degree sequence independent
of neighborhoods, and every integer $k$,
\[
 \Pr(M_\RV\geq k)\geq\Pr(M_\RK\geq k).
\]
The number of jobs is not bounded. Uniform random job order follows
by averaging. This is a written reduction with an exhaustive finite
integer verification; it is not a Lean certificate.
\end{corollary}
\begin{proof}
The accompanying Azure programs exhaustively verify
\eqref{eq:fiveinteger} for $2\leq N\leq400$ and
$0\leq S\leq3N+4$. For all other integer $S$, both sides are zero.
There are 242,592 comparisons, of which 159,201 are strict and
83,391 are equalities.

The first program constructs $c_{N,S}$ by the displayed convolution
recurrence across $N$. The independent program fixes each $N$ and
constructs the coefficients $d_j$ of $(6+3z+z^2)^N$ using
\[
 d_0=6^N,\qquad
 6j d_j=3(N-j+1)d_{j-1}+(2N-j+2)d_{j-2},
 \quad 1\leq j\leq2N,
\]
where $d_{-1}=0$, and then sets $c_{N,S}=d_{S-N}$, with
out-of-range coefficients zero. This recurrence follows by
differentiating $D(z)=(6+3z+z^2)^N$ and comparing coefficients
in $(6+3z+z^2)D'=N(3+2z)D$. The implementation checks every
division for zero remainder; it also checks $\sum_jd_j=10^N$
and $d_{2N}=1$. Direct multinomial sums independently reproduce
722 coefficients for $2\leq N\leq20$.

Both exhaustive implementations pass. Their source hashes, returned
results and case hashes are retained with the Azure reports under
\texttt{capacity-five/}; the independent verifier is
\texttt{independent-audit/check\_certificate.py}. Apply the
finite-range part of Theorem~\ref{thm:fiveconditional} with $L=400$.
This uses at most $m-1\leq400$ residual bins, regardless of the
number of arrivals.
\end{proof}

\noindent\textbf{Stage assessment.}
This step removes the day-count restriction from capacity five by a
written proof valid for all $N$, rather than by enlarging a finite
screen. The proved regimes are now $m\leq4$ or $C\leq5$.
The remaining region has $m\geq5$, $C\geq6$, $k\geq2C+3$.
It therefore advances an unbounded slice of the original conjecture,
but it does not prove the conjecture for every capacity.

The next objective is now a uniform-capacity mechanism, developed in
Section~\ref{sec:uniformgate}. Proving another individual capacity
without such a mechanism would not settle the universal quantifier.
The earlier suggested capacity-six endpoint calculation is therefore
not the current research gate.

\section{A uniform-capacity proof gate and failed shortcuts}\label{sec:uniformgate}
The next objective is a statement valid for every capacity, rather
than another individual capacity. The proved matching regimes remain
$m\leq4$ or $C\leq5$. We record a sufficient mechanism that has no
capacity cutoff, together with exact obstructions to two stronger
versions of that mechanism. These obstructions concern auxiliary
methods, not the matching conjecture.

\subsection{A boundary potential suffices}
Fix $b\geq1$, $q\geq1$ and $0\leq r\leq bq$. On the labeled state space
\[
 \Omega_{b,r,q}=\{x\in\{0,\ldots,q\}^b:\textstyle\sum_i x_i=r\}
\]
let $\mu(x)$ be proportional to $\prod_i1/x_i!$, and put
$B(x)=\#\{i:x_i=q\}$. Consider the continuous-time generator
\begin{equation}\label{eq:urngenerator}
 (Lf)(x)=\sum_{i\ne j}x_i\mathbf1_{\{x_j<q\}}
       \bigl(f(x-e_i+e_j)-f(x)\bigr),
\end{equation}
where zero-rate terms with $x_i=0$ are omitted. This chain selects
an individual ball and attempts to move it to a different bin;
a full target blocks the move. A common rate normalization would
not change the statements below.

The chain is irreducible: if $x\ne y$, choose an index with
$x_i>y_i$ and another with $x_j<y_j$, and move a ball from $i$ to $j$.
The target is below $q$, and this decreases the distance to $y$.
It is reversible for $\mu$, since for $y=x-e_i+e_j$,
\[
 \mu(x)x_i=\mu(y)(x_j+1).
\]
There is therefore a unique mean-zero solution of the finite Poisson
equation
\begin{equation}\label{eq:boundarypoisson}
 -L\psi=B-\mathbb E_\mu B,\qquad \mathbb E_\mu\psi=0.
\end{equation}
Existence follows because the right side has mean zero and the kernel
of the irreducible reversible generator consists of constants.
Uniqueness and permutation symmetry imply that $\psi$ is symmetric.
For a singleton state space the solution is simply zero.

Equivalently,
\[
 \psi(x)=\int_0^\infty
       \bigl(\mathbb E_x B(X_t)-\mathbb E_\mu B\bigr)\,dt.
\]
The integral converges for a finite irreducible continuous-time chain;
differentiating its semigroup verifies the Poisson equation.
Thus $\psi$ measures cumulative excess boundary occupation during
random redistribution, starting from $x$. The open question is whether
starting from a more unequal profile can ever decrease this quantity.

\begin{theorem}[A uniform boundary-potential criterion]\label{thm:boundarypotential}
If the solution $\psi$ of \eqref{eq:boundarypoisson} is increasing
under majorization, then
\[
 \operatorname{Cov}_{\mu}(B,F)\geq0
\]
for every symmetric majorization-increasing function $F$ on
$\Omega_{b,r,q}$. If this potential monotonicity holds for every
$b,r,q$, then Arnosti's full conjecture follows for every capacity.
\end{theorem}
\begin{proof}
By reversibility and \eqref{eq:boundarypoisson},
\begin{align*}
 \operatorname{Cov}_\mu(B,F)
 &=\langle-L\psi,F\rangle_\mu\\
 &=\frac12\sum_{x\in\Omega_{b,r,q}}\sum_{i\ne j}
 \mu(x)x_i\mathbf1_{\{x_j<q\}}
 \bigl(\psi(x-e_i+e_j)-\psi(x)\bigr)
 \bigl(F(x-e_i+e_j)-F(x)\bigr).
\end{align*}
Every move changes only two entries with fixed sum. The old and new
profiles are comparable by majorization: the move either balances
the pair, unbalances it, or just permutes the pair. Thus the two
differences in each summand have the same sign, or are both zero.
This proves the covariance inequality. Apply it to each prefix
weight $W_\alpha$, whose majorization monotonicity is already proved,
and use the uniform-next-choice lemma and
Theorem~\ref{thm:hazardreduction}.
\end{proof}

This criterion asks about one potential, not preservation of every
increasing function by the chain. Its sufficiency is proved here;
its monotonicity premise for arbitrary parameters is not assumed
or asserted. A finite computation can falsify the premise, or check
an algebraic proposal for it, but cannot establish it for all $q$.

\subsection{The simplest globally monotone chain does not work}
One might try to prove that the semigroup $P_t=\exp(tL)$ preserves
every majorization-increasing function. That stronger assertion is
false. If it were true, then for comparable $x\preceq y$ and an
increasing $f$ with $f(x)=f(y)$, differentiating
$P_tf(x)\leq P_tf(y)$ at $t=0$ would give $Lf(x)\leq Lf(y)$.

For every $q\geq4$, take four bins and
\[
 x=(q-1,3,1,0),\qquad y=(q,2,1,0),\qquad r=q+3.
\]
The profile $y$ majorizes $x$. The symmetric function
$f(x)=\#\{i:x_i=0\}$ increases in majorization and equals one at
both profiles. Write $m_1(x)$ for the number of singleton bins.
Directly from \eqref{eq:urngenerator},
\[
 Lf(x)=m_1(x)\bigl(b-B(x)-1\bigr)-r f(x).
\]
A singleton source creates a zero when it moves to any allowed
different target; filling a zero removes one, with total incoming
rate $r$ per zero. These contributions also correctly cancel for
a singleton moving into a zero.
In the displayed pair $m_1(x)=m_1(y)=1$, $B(x)=0$, $B(y)=1$,
so $Lf(x)=-q$ whereas $Lf(y)=-q-1$. This reverses the necessary
inequality. The obstruction therefore persists for arbitrarily large
caps, not merely for one exceptional parameter value.

\subsection{Resampling a pair does not repair full monotonicity}
A second natural generator resamples each unordered pair of bins,
keeping its sum $s$ fixed, according to the capped binomial law
\[
 \Pr\bigl((x_i,x_j)=(a,s-a)\bigr)
 =\frac{\binom sa}{\sum_{u=\max(0,s-q)}^{\min(q,s)}\binom su},
 \qquad \max(0,s-q)\leq a\leq\min(q,s).
\]
Each unordered pair has rate one. This heat-bath chain is also
reversible for the same capped multinomial law. Nevertheless it
fails to preserve majorization monotonicity.

Take $b=3$, $q=4$, $r=6$, and comparable profiles
$x=(4,1,1)\preceq y=(4,2,0)$. The feasible sorted profiles are
\[
 (4,2,0),\ (4,1,1),\ (3,3,0),\ (3,2,1),\ (2,2,2).
\]
Let $f$ equal one except at the balanced profile $(2,2,2)$.
This is increasing, with $f(x)=f(y)=1$.
No resampling of one pair can reach $(2,2,2)$ from $x$, because
none of its untouched entries equals two. Thus $L_{\rm hb}f(x)=0$.
From $y$, resampling the pair $(4,0)$ gives $(2,2)$ with probability
$\binom42/2^4=3/8$, while retaining the untouched two.
Consequently $L_{\rm hb}f(y)=-3/8$, another violation.
This example lies within an already proved matching regime, so it
cannot be a counterexample to the matching claim or to the stationary
association result. It rejects this particular dynamic proof route.

\subsection{Heterogeneous rows cannot simply be mixed into iid rows}
A different shortcut would replace an ordered collection of different
categorical rows by a positive mixture of identical-row models,
after symmetrizing the bin labels. This representation fails even
for two bins and four choices.

Take two choices deterministically in the first bin and two fair
choices, then randomly exchange the bin labels. The first-bin count
$K$ has probabilities
\[
 \Pr(K=0,1,2,3,4)=\tfrac18(1,2,2,2,1).
\]
Suppose this were a mixture $K\mid P\sim\operatorname{Bin}(4,P)$,
with $0\leq P\leq1$. The factorial-moment identities force
\[
 \mathbb EP=\tfrac12,\quad
 \mathbb EP^2=\tfrac7{24},\quad
 \mathbb EP^3=\tfrac3{16},\quad
 \mathbb EP^4=\tfrac18.
\]
They imply simultaneously
$\mathbb E(P-1/2)^2=1/24$ and $\mathbb E(P-1/2)^4=0$.
The latter forces $P=1/2$ almost surely, contradicting the former.
For any $q\geq4$, imposing the cap does not change this example.
Thus a theorem for iid tilted multinomial models cannot be extended
to our heterogeneous prefix rows by this positive-mixture shortcut.

\noindent\textbf{Literature check.}
Vitale's
\href{https://doi.org/10.1016/S0167-7152(99)00064-4}
{\emph{Majorization and Gaussian correlation}} (1999) contains a
multinomial correlation argument. Cohen and Sackrowitz's
\href{https://doi.org/10.1214/aos/1176350376}
{\emph{Unbiasedness of tests for homogeneity}} (1987) treats
log-concave exponential families and Schur-convex tests.
These are potentially relevant leads, but the full hypotheses needed
for our capped heterogeneous model have not been verified in this
stage: the accessible preview and scanned text were insufficient for
that inference. Neither result is used as a premise of our proofs.
Chan's normalized matching theorem concerns a different order on sets
of balls; no transfer to capped-profile majorization is asserted.

\noindent\textbf{Subsequent source verification.}
The full Cohen--Sackrowitz paper was subsequently obtained from the
publisher and all twelve scanned pages were inspected. Its Lemma 4.1
does apply to the capped iid reference measure, with heterogeneous
rows handled separately by their increasing density. This closes
the missing premise in Section~\ref{sec:literatureclosure}; the earlier
source-access limitation above is retained only as a stage record.

\subsection{Exact diagnostics and the stage decision}
The Azure transport audit rejected full generator monotonicity for
both chains, with the witnesses above (and a smaller three-bin witness
for the one-ball chain). It checked the four-bin obstruction formula
at 97 caps, $4\leq q\leq100$; the displayed calculation proves it
for every $q\geq4$. Separate finite diagnostics found no failure in
8,693 principal-event association tests or 756 boundary-layer order
tests. Those successes do not prove stationary association in general.

A second Azure job solved \eqref{eq:boundarypoisson} exactly on
850 profile spaces. Its ranges were $b=3,4$ with $1\leq q\leq12$,
$b=5$ with $1\leq q\leq8$, and $b=6$ with $1\leq q\leq6$,
at all feasible totals except the 40 spaces having more than
55 sorted profiles. It fixed one potential value to zero instead of
its mean; this adds a constant and has no effect on monotonicity.
All 9,997 equation rows were checked by substitution, and all
258,093 detailed-balance comparisons used the actual multinomial
orbit weights. None of 102,759 comparable-profile comparisons
violated potential monotonicity.

\noindent\textbf{Stage assessment.}
No new matching regime and no proof of the full conjecture is claimed
by these diagnostics. This stage discards two false global dynamics
arguments and the false iid-mixture shortcut, and proves a sufficient
criterion involving one boundary potential for arbitrary parameters.
The monotonicity of that potential remains open here. The next gate
is a uniform proof of its majorization monotonicity (or an exact
counterexample to it), not a larger collection of solved finite
linear systems and not the next individual capacity.

\subsection{Why even the boundary observable cannot be compared at every time}
A later diagnostic also rejects the narrower attempt to show
$P_tB(x)\leq P_tB(y)$ at every time whenever $x\preceq y$.
For $b=3$, $q=4$, $r=7$, list the sorted states as
\[
 (4,3,0),\quad (4,2,1),\quad (3,3,1),\quad (3,2,2).
\]
The generator and boundary vector are
\[
 L=\begin{pmatrix}
 -7&3&4&0\\1&-9&4&4\\2&6&-14&6\\0&4&4&-8
 \end{pmatrix},\qquad B=(1,1,0,0)^{\mathsf T}.
\]
Put $x=(3,3,1)\preceq y=(4,2,1)$. Direct solution of this
four-dimensional system gives
\[
 (e^{tL}B)(y)-(e^{tL}B)(x)
 =\frac17\left[4e^{-18t}+e^{-10t}
 \left(3\cosh(2\sqrt2\,t)-\frac5{\sqrt2}\sinh(2\sqrt2\,t)\right)\right].
\]
For example, this formula is checked by the eigenvalues
$0,-18,-10\pm2\sqrt2$ and the initial difference and its first
two derivatives $1,-16,260$. The coefficient of the slowest mode
$e^{-(10-2\sqrt2)t}$ is $(3-5/\sqrt2)/14<0$, so the difference
eventually becomes negative. Nevertheless its integral is
\[
 \psi(y)-\psi(x)=\frac{13}{207}>0.
\]
In fact an anchored Poisson solution on these four states is
$\psi=(3/23,2/23,5/207,0)$, with
$\mathbb E_\mu B=4/9$, and direct substitution checks every row.
Thus an ordered integral need not arise from an ordered integrand.

The Azure job recorded a discrete precursor of this obstruction:
for $P=I+L/42$, $(P^{15}B)(y)-(P^{15}B)(x)
=-14003064062062976/68122318582951682301$.
Its other diagnostics found no failure of pairwise convexity of the
Poisson solution or of the tested discounted resolvents; those are
finite observations only. None is used in the general proof, which
now follows from the published stationary association theorem.

\section{Why balancing residual capacities is not enough}
A tempting route would assert that moving one free slot from a larger
residual capacity to a smaller one improves every RANDOM-VERTEX
completion probability. That assertion is false, even with three days
and capacities at most three.

Consider the same future degree sequence $(3,3,2)$ from the two residual
capacity vectors
\[
 u=(0,2,3),\qquad v=(1,1,3).
\]
Both have five free slots, and $v$ is obtained from $u$ by balancing
the first two entries. The first two jobs are fully flexible and always
succeed. The third job fails only if just one day remains available;
it then misses that day with probability $1/3$.

Starting from $u$, a single day remains after the first two jobs exactly
when both went to the day with two free slots. Each choice has probability
$1/2$, so this event has probability $1/4$. Starting from $v$, a single
day remains exactly when the two one-slot days were used: the probability
is $(2/3)(1/2)=1/3$. Thus
\[
 \Pr_u(\hbox{all three accepted})=1-\frac14\frac13=\frac{11}{12},
 \qquad
 \Pr_v(\hbox{all three accepted})=1-\frac13\frac13=\frac89.
\]
The more balanced vector is worse by $1/36$.
These are residual states, not the equal initial capacities required by
Arnosti's conjecture; both computations concern RANDOM-VERTEX alone.
Therefore this is a counterexample to a proposed auxiliary lemma, not
to the original conjecture. A general argument must account for the
distribution of residual states actually produced by each algorithm.

\section{Computational checks and their scope}
The accompanying standard-library Python programs ran in bounded Azure
workers. The main integer-weight recursion checked every degree multiset
with $C=3$, $2\leq m\leq4$ and $1\leq n\leq12$: 2,363 instances and
22,295 nonzero-threshold comparisons. No violation was found. Another
1,404 instances with $C=1,2$ reproduced the known equality and dominance
controls. Three small original-model enumerations, including all labeled
job orders and all shared day rankings, agreed with the recursion.
The first-rejection formulas were checked on 108 degree multisets.

A separate program tested both two-day first-saturation formulas and
the continuation mixture against direct state propagation for every
fixed sequence of degrees one or two of lengths zero through ten, and
for $1\leq C\leq6$: 12,282 cases. All comparisons agreed. It separately
verified the random-order example $31/32$ versus $7/8$, the corresponding
fixed-order example $(2,2,2,1)$ with probabilities $7/8$ versus $1/2$,
and the residual-capacity example $11/12$ versus $8/9$ by enumerating
the original neighborhoods.

The next Azure audit checked every fixed positive-degree sequence of
length zero through six for $m\in\{2,3,4\}$ and $1\leq C\leq4$:
26,724 cases. It verified the independent preferred-choice survival formula
against actual neighborhood transitions, and the uncapped continuation
formula whenever $n\leq2C$. It also checked 2,431 central-interval
inequalities and 720 adjacent-averaging steps. All required checks passed.
An exploratory check of later saturation ordering found violations;
one is the explicit $20/243$ versus $2/27$ example above.

A separate audit tested the stopped identity \eqref{eq:2C-mixture} beyond
$n=2C$: 1,013 structured and fixed-seed sampled degree sequences, using
$m\in\{2,3,4\}$, $1\leq C\leq5$, and lengths $2C+1$, $2C+2$, and
$3C+3$. Both algorithms' stopped laws agreed exactly with direct
neighborhood-based state propagation in all 2,026 comparisons, and every
tail through $2C$ had the required order. The sampled objects were degree
sequences; their outcome probabilities were computed exactly, not estimated
by Monte Carlo.

To test the stronger marked hitting-state order, another Azure job
enumerated 269,319 fixed degree sequences in the following batches:
$(m,C,n_{\max})=(3,2,8),(3,3,10),(4,3,8),(3,4,10)$,
retaining lengths $2C\leq n\leq n_{\max}$ in each batch. It independently
propagated the original neighborhood transitions and verified the
one-additional-success identity in 538,638 algorithm instances.
For each sequence it checked every lower set of the marked state space,
of the form $\sigma\leq t_A$ with $t_0\geq t_1\geq t_2$; all had the
required probability order. All tested tails at threshold $2C+1$ did too.
These finite results motivated the unbounded marked coupling; they are
not its proof.

A separate conditional audit then tracked first saturation and first
arrival at $2C$ successes directly, stopping each trajectory on arrival.
Across 59,681 retained degree sequences, it verified 118,331 exact
RANDOM-VERTEX equalities and 114,951 RANKING lower bounds from
Lemma~\ref{lem:mark}; 106,240 of the latter were strict.
The parameter batches were $(m,C,n_{\max})=(2,2,7),(3,2,7),(3,3,8),
(4,3,7),(3,4,9)$. A negative control rejected the incorrect assertion
that RANKING always has equality: for six fully flexible jobs with
$m=C=3$, conditional on $T=3$ and $\sigma=6$, the two-full-day
probabilities are $1/4$ for RANDOM-VERTEX and one for RANKING.

The audit for the extension to $2C+2$ checked
Lemma~\ref{lem:nearfull} on 174,785 prefix-uniform configurations
with $2\leq b,c\leq8$ and on 480 fixed-seed rational ordered-row
configurations. In 597 of these configurations, direct occupancy
propagation independently checked the numerator and denominator formulas.
Actual neighborhood propagation then checked 425,014 RANDOM-VERTEX
conditional equalities and 401,063 RANKING conditional lower bounds
from Lemma~\ref{lem:fillsecond}, including the boundary $T=2C$.
Of the RANKING inequalities, 379,571 were strict. The batches were
$(m,C,n_{\max})=(3,2,6),(3,3,8),(4,3,7),(3,4,9)$; every possible
next positive degree was checked at each retained conditional state.

Finally the same audit checked the marked hitting order at $2C+1$
successes and the tail at $2C+2$ on 257,852 fixed degree sequences.
It used the earlier marked-diagnostic parameter batches, now retaining
lengths $2C+1\leq n\leq n_{\max}$. Every lower-set comparison passed,
and 515,704 one-further-success identities agreed exactly with direct
matching-state propagation. No tested marked or tail inequality failed.
The proof above, rather than these finite checks, establishes the
unrestricted-parameter extension.

The three-day audit checked the two-bin hazard formula in 85,848
extreme-row cases for $2\leq C\leq50$ and every
$C-1\leq r\leq2C-2$. In 54 small cases, independent enumeration
of labeled fair choices verified the binomial counts. The matching
checks covered 617,385 degree-sequence instances: exhaustive positive
degree sequences with $2C\leq n\leq n_{\max}$ for
$(C,n_{\max})=(2,9),(3,11),(4,10),(5,11)$, plus 1,056 structured
or fixed-seed sampled sequences at $2\leq C\leq12$ with lengths
$3C$, $3C+1$, $4C$, and $6C$.
All 7,153,740 tail comparisons passed. At the tested success levels
from $2C$ through $\min(3C-1,n)$, 1,487,727 marked lower-set
comparisons passed and 2,975,454 one-further-success identities agreed
with direct neighborhood-based state propagation.

The general conditional-hazard diagnostic checked 144,186
prefix-uniform cases, with $2\leq b\leq6$, $2\leq C\leq5$, and
$C-1\leq r\leq\min(b(C-1),2C+1)$. No inequality violation was
found. Increasing one old prefix length by one also never increased
the hazard for a uniform next choice in the tested pairs. These are
finite diagnostics only: neither the general inequality nor this
additional monotonicity assertion is proved here.

The FKG audit checked 1,822,977 pairs of capped three-bin profiles
for $2\leq C\leq32$ and all feasible totals. Their coordinatewise
meet and join remained in the profile set, and the actual multinomial
orbit weights satisfied the log-supermodular inequality in every case.
It checked the monotonicity of the prefix-restriction weight on
742,048 balancing transfers with totals through 30. Another 35,959
conditional hazard comparisons covered $2\leq C\leq15$ and all
feasible prefix-count triples. Direct occupancy propagation independently
verified the tilted-profile formula in 494 small cases and all three
next-prefix comparisons in 1,482 cases.

The four-day matching audit checked 763,472 degree-sequence instances
and 10,125,888 tail inequalities, all successfully. The exhaustive
batches retained every positive degree sequence with
$2C\leq n\leq n_{\max}$ for $(C,n_{\max})=(2,8),(3,9),(4,9)$.
Another 336 structured or fixed-seed sampled sequences used
$2\leq C\leq8$ and lengths $4C$, $4C+1$, and $6C$; their output
probabilities were computed exactly. These finite audits support the
implementation and the algebra; Corollary~\ref{cor:four} rests on the
written argument and the cited general FKG theorem.

The next diagnostic tested boundary association against \emph{every}
increasing function on each of 499 finite profile spaces: $b=4$ with
$1\leq q\leq12$, $b=5$ with $1\leq q\leq6$, and $b=6$ with
$1\leq q\leq4$, in each case at every feasible total $r$.
An integer minimum-cut computation finds the minimum of
$\sum_{n\in U}w(n)(B(n)Z-Z_B)$ over all dominance upper sets $U$,
where $Z=\sum_n w(n)$ and $Z_B=\sum_n w(n)B(n)$.
Nonnegative covariance for all such indicators implies it for all
increasing functions. No negative value was found. The same job
checked all 76,488 feasible prefix-count configurations for
$b=4$, $1\leq q\leq8$, and $b=5$, $1\leq q\leq3$.
All structured covariance inequalities passed; direct occupancy
propagation agreed in 1,284 cases and checked 5,136 next-prefix
comparisons. In 197,209 tested single-prefix extensions, the
uniform-next hazard did not increase. This last monotonicity remains
unproved without parameter bounds.

The capacity-three/four proof audit checked profiles for every
$1\leq b\leq40$ and every feasible total. It verified 129,800
cap-two chain comparisons, 7,563,485 cap-three lattice and actual
weight comparisons, and 15,126,970 comparisons between coordinate
order and dominance. Independent occupancy propagation checked
767 structured or fixed-seed sampled configurations on two through
six bins, including 3,690 next-prefix hazard comparisons.
Direct original-neighborhood matching propagation checked 96
structured or fixed-seed sampled degree sequences for
$m\in\{5,6\}$, $C\in\{3,4\}$, with lengths $mC$, $mC+2$,
and $2mC$. Across their 1,952 retained prefixes it checked
38,680 exact tail inequalities. All passed. These checks are
supplementary: Corollary~\ref{cor:capacityfour} follows from the
unbounded lattice proof and the general hazard reduction.

The capacity-five layer diagnostic checked the coefficient inequality
in 242,592 cases with $2\leq N\leq400$, all successfully.
For $2\leq N\leq30$ and $|K|\leq N-2$, all 841 shifted
polynomials $Q_{N,K}$ had nonnegative coefficients. This is a
symbolic finite check, not a proof for arbitrary $N$.
Across 560 profile spaces with cap four, $1\leq b\leq16$, and
every feasible total, it checked 124,901 endpoint-marginal weight
comparisons and 347,616 conditional tail comparisons, all successfully.
An exact minimum-cut check also found no failure of dominance order
between consecutive nonempty $B$-layers for $b\leq10$.
The independent coefficient audit for
Corollary~\ref{cor:fivefinite} recomputed all 242,592 inequalities
by the differential recurrence and checked the stated direct sums.

The unbounded capacity-five proof received a separate finite algebraic
audit. It checked 15,100 instances of the alternating-binomial window
identity, 14,430 interval-indicator comparisons, 25,513 instances of the
coefficient-extraction identity, and 3,953 trinomial determinants for
$2\leq n\leq60$, including zero-support boundary cases. The exact
substitution into the original integer inequality agreed in all 3,953
cases. All nonconstant raw coefficients and all shifted coefficients
had the stated signs; 1,947 constant/quadratic relations and 1,711
explicit negative-constant formulas agreed. A negative control rejected
the false assertion that every raw coefficient is nonnegative:
$\delta_{3,0,0}=-3$. The Azure worker exited successfully and returned
source-verified results. These tests check the algebra; the universal
quantifier in Corollary~\ref{cor:capacityfive} follows from the written
binomial argument, not the finite audit.

These are exact finite checks of the implementations and formulas.
The computer-assisted result is explicitly identified as such;
the other unrestricted-parameter statements rest on the written proofs above.
These historical numerical audits do not constitute Lean proofs.
The current partial formalization is described after
Theorem~\ref{thm:fullconjecture}; no independent human review is claimed.
Machine-readable stage assessments and cloud audit records accompany
the source under \texttt{checks/greedy-matching-capacity/}.
\end{document}
